\chapter{Inter-Constituent Dynamics and the b-Field}
\label{ch:ch09_bfield}

The question of this chapter is what can act between constituents, and with what rates. It treats the connection part of the
shift and the gauge transport of the two energies, the $b$-field between manifolds, the $b$-field up the ladder, the closure at two
constituents, the $b$-field on the temporal triple with $c$ on the spatial triple, and, in \cref{sec:ladder}, the ladder table and what changes
up the ladder. The connection part of the shift splits the binding (its part commuting with $G$) or destabilises it (its anticommuting part),
and never creates it. A compact field between constituents is a one-body operator: it moves quanta and keeps their number, so it does not
take one particle to two. It is removable between equal rates and, with the prescribed binding tensor kept fixed, non-removable exactly between unequal ones, and it changes no rate. From three
constituents it carries time-odd plaquette invariants. Copy isospin is integral at two and three constituents under a common Lorentz action;
doublets appear from four. The algebra is proved and checked. Reading a rate as a mass, $mc^2=\hbar\kappa$, is an identification, and reading the
plaquette as CP violation is an analogy; neither is a result.

Several letters carry more than one meaning in this chapter, each fixed by its context: $b$ is an index and the connection; $S_n$ is the spinor
module and, in \cref{e6:thm:ladder}, the spatial volume; $N$ is the number operator, the inertias $N_k$ and a nilpotent part; $H$ is the Noether
energy, the copy Hamiltonian and the operator $\mathsf H$; $C$ is the conversion, the copy rotations $C_\pm$, $C^{kl}$, $\hat C$, $C_v$, $C_b$ and a complex
symmetric matrix; and $M$ is the phase-space generator and a mass.

% --- source: tex/s17c_bfield#0 ((head))


% --- source: tex/s17c_bfield#1 (The connection part of the shift)
\section{The connection part of the shift}
\label{sec:bfield}

A constant shift of the interior Lagrangian has two summands (\cref{thm:EAdichotomy}): an $\eta$-antisymmetric one, in $\so{3,3}$, and an
$\eta$-self-adjoint one. \Cref{sec:neutral,sec:closure,sec:restclosure} were built on the second, the binding term $\kappa G$, and met the first only as transport along a path (\cref{prop:tear}) and in \cref{thm:EAdichotomy}. This section adds the
first, a constant $b\in\so{3,3}$, and asks what it does to the binding. Only $b$ is a connection, and where it is gauged it is the value of a
Yang--Mills potential along the worldline, the ``$b$-field'' in the notation of Yang and Mills \cite{YangMills54}. In brief, the answers are as follows. For
every $b$ the field equation factorises covariantly and the binding term transforms as a tensor, so no gauge transformation creates it
(\cref{bf:covariant}). An admissible gauge potential is compact. Its part commuting with $G$ splits the binding frequency into Zeeman lines that a
moving frame removes (\cref{bf:larmor,bf:lG}); a purely anticommuting part tears the binding (\cref{bf:pG}); a mixed one keeps every orbit bounded
exactly when the two triples rotate at equal rates about axes that are not opposite, outside one exceptional ratio (\cref{bf:screen}). On its own
a connection does not bind, and nothing in the interior dynamics ties the mass $m$ to the rate $\kappa$ (\cref{bf:alone}). Wherever the motion is
bounded the zero-point energy cancels frequency by frequency (\cref{bf:krein}). At three constituents the part of the Pati--Salam algebra
compatible with the binding splits the frequency by $\su2$ weights (\cref{bf:so4spec}). None of this supplies a scale. That the scale of a quantum
Yang--Mills theory is the binding length is stated as a conjecture (\cref{bf:transmute}).

Units are $\hbar=c=1$ unless a symbol is shown. The parameter $\tau$ has the dimension of a time, $\kappa=c/\ell$ of \eqref{eq:kappa} and every
eigenvalue of $b$ are rates, and $\hbar\kappa$ is a quantum of interior energy.

\paragraph{Notation.} $V=\R^{3,3}$, $F^\dagger:=\eta^{-1}F^{\!\top}\eta$ is the $\eta$-adjoint, $G$ is the binding root \eqref{eq:Gpair} and $\cK$
the grading root of \cref{prop:twoK}; on vectors $\cK=\tfrac12\ad\bigl(\sum_ie_ie_{i+3}\bigr)$ and $G=-\cK\eta$ (\cref{lem:Gclifford}). For indices $a<b$ (here $b$ is an index, not the connection) the
element $R_{ab}\in\so{3,3}$ is defined by $R_{ab}e_b=\eta_{bb}e_a$, $R_{ab}e_a=-\eta_{aa}e_b$, and zero on the other basis vectors; it is
$\tfrac12\ad(e_ae_b)$ on vectors. The generators $L_i,K_i$ of $\lG$ (\cref{thm:emfield}) act on $z=\bB+i\bE$ as $z\mapsto\mathbf e_i\times z$ and
$z\mapsto i\,\mathbf e_i\times z$, so $K_i=GL_i$, the $L_i$ are rotations and the $K_i$ boosts, and
\begin{equation}
\label{bf:eq:w}
b=\textstyle\sum_i\bigl(w^{\rm r}_iL_i+w^{\rm b}_iK_i\bigr)\in\lG\quad\text{acts as}\quad z\mapsto w\times z,\qquad w=w^{\rm r}+i\,w^{\rm b}\in\C^3 .
\end{equation}
In particular $R_{12}-R_{45}=-L_3$. We write $\sigma_j$ for the boost exchanging $e_j$ and $e_{j+3}$, $[\mathbf v]_\times$ for the matrix of
$\mathbf x\mapsto\mathbf v\times\mathbf x$ on $\R^3$, $\mathcal E_\mu(T)=\ker(T-\mu)$ for the eigenspaces of an operator on $V\otimes\C$, and
$h_\eta(y,y')=\bar y^{\!\top}\eta\,y'$ for the Hermitian extension of $\eta$. The interior Lagrangian with a constant shift is
\begin{equation}
\label{bf:eq:L}
\mathcal L=\frac m2\,\eta(u,u),\qquad u=\dot X+SX,\qquad S=b+\kappa G,\qquad b\in\so{3,3},
\end{equation}
with $p=m\,\eta u$ and $\pi=\eta^{-1}p=mu$. For $b=0$ it is the split-form Lagrangian \eqref{eq:splitL}.

% --- source: tex/s17c_bfield#2 (Two summands, and which one is a connection)
\subsection{Two summands, and which one is a connection}
\label{sec:bf-split}

By \cref{thm:EAdichotomy} at $n=1$, $\End(V)=\so{3,3}\oplus\{F:F=F^\dagger\}$, of dimensions $15+21$, and by \cref{prop:twoK} the binding root
lies in the second summand and the grading root in the first. The $\eta$-antisymmetric part of $S$ is therefore $b$, and its self-adjoint part
is $\kappa G$. The decomposition is linear algebra, valid for any nondegenerate form; the next theorem gives it a dynamical meaning.

\begin{theorem}[Covariant factorisation]
\label{bf:covariant}
Let $b\in\so{3,3}$ be constant and $\nabla:=\frac{d}{d\tau}+b$.
\begin{enumerate}[label=(\roman*),leftmargin=2.2em]
\item The Euler--Lagrange equations of \eqref{bf:eq:L} are $(\nabla-\kappa G)\pi=0$ with $\pi=m(\nabla+\kappa G)X$, that is
\begin{equation}
\label{bf:eq:covEL}
(\nabla-\kappa G)(\nabla+\kappa G)X=0\qquad\Longleftrightarrow\qquad \nabla^2X+\kappa^2X+\kappa\,(\nabla_{\!\tau}G)\,X=0,\qquad \nabla_{\!\tau}G:=[b,G].
\end{equation}
\item Let $X=g(\tau)Y$ with $g(\tau)\in O(3,3)$. Then $\bigl(\tfrac d{d\tau}+S\bigr)(gY)=g\bigl(\tfrac d{d\tau}+S^g\bigr)Y$ with
$S^g=g^{-1}Sg+g^{-1}\dot g$. The $\so{3,3}$ part $b^g=g^{-1}bg+g^{-1}\dot g$ of $S^g$ transforms as a connection, its self-adjoint part
$\kappa\,g^{-1}Gg$ as a tensor, and $\tfrac d{d\tau}(g^{-1}Gg)+[b^g,g^{-1}Gg]=g^{-1}(\nabla_{\!\tau}G)g$.
\end{enumerate}
\end{theorem}

\begin{proof}
(i) $\partial\mathcal L/\partial\dot X=p$ and $\partial\mathcal L/\partial X=S^{\!\top}p$, so $\dot p=S^{\!\top}p$, that is
$\dot\pi=S^\dagger\pi=(-b+\kappa G)\pi$, which is $(\nabla-\kappa G)\pi=0$. Moreover $\pi=m(\nabla+\kappa G)X$, $\nabla(GX)-G\nabla X=[b,G]X$ and
$G^2=-\id$. (ii) Differentiating $g^{\!\top}\eta g=\eta$ gives $(g^{-1}\dot g)^\dagger=-g^{-1}\dot g$, and $F\mapsto g^{-1}Fg$ commutes with the
$\eta$-adjoint, so it preserves both summands; the last identity is a direct differentiation.
\end{proof}

Part~(ii) is inherited: it is the transformation law of a connection and of an endomorphism field in any gauge theory. Part~(i) is fixed. It is
the internal-connection, split-form version of the curved second-order equation at the end of \cref{sec:shifted}, which, as there, adds to the case $b=0$ only the term
$\kappa(\nabla_{\!\tau}G)X$. That the binding part transforms as a tensor restates $G=G^\dagger$: no gauge transformation creates, removes or
rescales it. The connection is the transport part of the shift, and no other part of the shift can be one.

\begin{proposition}[Commuting and anticommuting parts]
\label{bf:lp}
$\Ad_G:b\mapsto GbG^{-1}$ is an involution of $\so{3,3}$, so $\so{3,3}=\lG\oplus\mathfrak p_G$, where $\lG=\{b:[b,G]=0\}$ has dimension $6$ and
$\mathfrak p_G=\{b:bG=-Gb\}$ has dimension $9$. In the variable $z$,
\[
\lG=\{z\mapsto(P+iQ)z\}\cong\so{3,\C},\qquad \mathfrak p_G=\{z\mapsto(C+i\Sigma)\bar z\},
\]
with $P,Q,C$ real antisymmetric and $\Sigma$ real symmetric. The Killing form $4\Tr$ has signature $(3,3)$ on $\lG$. On $\mathfrak p_G$ it is
negative on the three directions $\Sigma=0$, compact counter-rotations of the two triples such as $R_{12}+R_{45}$, and positive on the six boosts
$C=0$. Both $\cK$, which acts as $z\mapsto-i\bar z$, and the pairing boosts $\sigma_j$ lie in $\mathfrak p_G$.
\end{proposition}

\begin{proof}
$G^{-1}=-G$ and $G^\dagger=G$ give $(GbG^{-1})^\dagger=-GbG^{-1}$, and $\Ad_G^2=\id$. In $(\bB,\bE)$ blocks, commuting with $G$ means the form
$\bigl(\begin{smallmatrix}P&-Q\\Q&P\end{smallmatrix}\bigr)$ and anticommuting means $\bigl(\begin{smallmatrix}C&\Sigma\\\Sigma&-C\end{smallmatrix}\bigr)$;
$b^\dagger=-b$ gives the symmetry types. The Killing form of $\so{3,3}$ is $4\Tr$, and on the two block forms $\Tr b^2=2\Tr(P^2-Q^2)$, respectively $2\Tr(C^2+\Sigma^2)$, which gives the signatures.
\end{proof}

This is fixed; the description of $\lG$ restates \cref{thm:GLorentz}. The pairing boosts that tear the binding at a threshold in \cref{prop:tear}
are anticommuting, and so is the grading root.

\paragraph{Gauge admissibility.} A gauge-admissible potential is compact in the sense fixed after \cref{cor:noboost}: its generators are
unitary for a positive form on the module (\cref{thm:sig}). It therefore takes values in the compact part $\so3\oplus\so3$ of $\so{3,3}$, spanned
by the $R_{ab}$ with $a,b$ in the same triple. On the two triples,
\begin{equation}
\label{bf:eq:compact}
b=\operatorname{diag}(A+D,\,A-D),\qquad A=[\mathbf a]_\times,\quad D=[\mathbf d]_\times,
\end{equation}
so the temporal triple rotates about $\mathbf a+\mathbf d$ and the spatial triple about $\mathbf a-\mathbf d$. The diagonal part $(A,A)$ lies in
$\lG$, with $w=\mathbf a$ real in \eqref{bf:eq:w}; the anti-diagonal part $(D,-D)$ lies in $\mathfrak p_G$, with $\Sigma=0$ (it acts as
$z\mapsto\mathbf d\times\bar z$). The boost-valued connections used below ($K_i$, $\sigma_j$, $\cK$, the $\Sigma$-part of $\mathfrak p_G$) are metric
connections in the sense of \cref{prop:tear}, legitimate in \eqref{bf:eq:L}, but they are not gauge fields. What a compact $b$ does pointwise to
the undivided energy and to the rest energy is \cref{en:transport}; here it acts through the interior dynamics.

% --- source: tex/s17c_bfield#3 (The phase-space generator and a criterion for boundedness)
\subsection{The phase-space generator and a criterion for boundedness}
\label{sec:bf-criterion}

\begin{lemma}[The phase-space generator]
\label{bf:M}
On $(X,\pi)$ the flow of \eqref{bf:eq:L} is generated by
\[
M=\begin{pmatrix}-S&\id/m\\0&S^\dagger\end{pmatrix},\qquad S^\dagger=-b+\kappa G,
\]
since $\dot X=\pi/m-SX$ and $\dot\pi=S^\dagger\pi$. It is the Hamiltonian flow of the Noether energy $H(X,p)=\tfrac1{2m}p^{\!\top}\eta^{-1}p-p^{\!\top}SX$:
in the canonical variables $(X,p)$ the generator is $\mathbb J\mathcal H$, with $\mathcal H$ the Hessian of $H$ and
$\mathbb J=\bigl(\begin{smallmatrix}0&\id\\-\id&0\end{smallmatrix}\bigr)$, and it is similar to $M$. $H$ vanishes on the invariant Lagrangian subspace
$p=0$, and $\det(\lambda-M)=\chi_S(\lambda)\chi_S(-\lambda)$ with $\chi_S(\lambda)=\det(\lambda-S)$, because $S^\dagger$ is similar to $S^{\!\top}$.
\end{lemma}

This is Hamiltonian bookkeeping, inherited. The classification below rests on the following criterion.

\begin{theorem}[When the interior motion stays bounded]
\label{bf:criterion}
Let $S$ be any constant real $6\times6$ matrix.
\begin{enumerate}[label=(\roman*),leftmargin=2.2em]
\item $M$ is semisimple if and only if $S$ is semisimple and, for every $\mu\in\C$, $\mathcal E_\mu(S^\dagger)$ and $\mathcal E_{-\mu}(S^\dagger)$ are
orthogonal for the bilinear extension of $\eta$.
\item Every solution is bounded on $\R$ if and only if, in addition, $\spec S\subset i\R$. Then (i) says that every eigenspace of $S^\dagger$ is
totally isotropic for $h_\eta$. Equivalently, every momentum normal mode $\pi=\Re(y\,e^{i\omega\tau})$ with $\omega\neq0$, superpositions within a
degenerate frequency included, has period mean $\langle\eta(\pi,\pi)\rangle=\half h_\eta(y,y)=0$, and the static momenta ($\omega=0$) span a totally
$\eta$-null space.
\end{enumerate}
\end{theorem}

\begin{proof}
$M$ is semisimple if and only if $-S$ (its restriction to $\pi=0$) and $S^\dagger$ (the induced map on the quotient) are semisimple and $M$ is
similar to $\operatorname{diag}(-S,S^\dagger)$, which by Roth's theorem \cite{Roth52} (for an invariant proof see \cite{FlandersWimmer77}) means that $-S\Xi-\Xi S^\dagger=-\id/m$ is solvable. With
$-S=F_1\Lambda_1F_1^{-1}$ and $S^\dagger=F_2\Lambda_2F_2^{-1}$ diagonal, this holds if and only if $(F_1^{-1}F_2)_{ij}=0$ whenever
$(\Lambda_1)_i=(\Lambda_2)_j=:\mu$. A row $r^{\!\top}$ of $F_1^{-1}$ with eigenvalue $\mu$ has $S^{\!\top}r=-\mu r$, and $S^{\!\top}=\eta S^\dagger\eta^{-1}$
puts $\eta^{-1}r$ in $\mathcal E_{-\mu}(S^\dagger)$; the condition is $r^{\!\top}y=\eta(\eta^{-1}r,y)=0$ for $y\in\mathcal E_\mu(S^\dagger)$. (ii) A linear
flow is bounded on $\R$ exactly when its generator is semisimple with imaginary spectrum, and then $\mathcal E_{-\mu}(S^\dagger)=\overline{\mathcal E_\mu(S^\dagger)}$.
\end{proof}

The theorem is fixed, although it holds for any nondegenerate $\eta$ and any $S$. Its necessity half is the split virial identity
(\cref{thm:splitvirial}: a bound orbit has null momentum on average); what it adds is that the averaged virial condition, imposed mode by mode, is
also sufficient. For $S=\kappa G$ it is the isotropy of \cref{prop:howmuch}.

% --- source: tex/s17c_bfield#4 (Connections commuting with G: a Zeeman pattern without a sca)
\subsection{\texorpdfstring{Connections commuting with $G$: a Zeeman pattern without a scale}{Connections commuting with G: a Zeeman pattern without a scale}}
\label{sec:bf-commuting}

\begin{proposition}[The commuting part is a moving frame]
\label{bf:larmor}
Let $b\in\lG$. Then $\bigl(\tfrac d{d\tau}+b+\kappa G\bigr)(e^{-b\tau}Y)=e^{-b\tau}(\dot Y+\kappa GY)$. The exponents are therefore those of the
binding, $\pm i\kappa$, shifted by the eigenvalues of $-b$, and with $p_X=e^{b^{\!\top}\tau}p_Y$ the energies in the two frames are related by
\[
H_{\rm lab}=H_{\rm co}-p_Y^{\!\top}bY.
\]
\end{proposition}

\begin{proof}
$e^{-b\tau}$ is an isometry that commutes with $b$ and $G$; moreover $e^{b\tau}\eta^{-1}e^{b^{\!\top}\tau}=\eta^{-1}$ and
$e^{b\tau}(b+\kappa G)e^{-b\tau}=b+\kappa G$.
\end{proof}

For compact $b$ this is Larmor's theorem \cite{Larmor1897}, inherited: a rotating frame trades the field for an angular velocity. It is also the
first statement of \cref{prop:tear}. For a boost in $\lG$, $e^{-b\tau}$ is a uniformly boosted frame instead. $Q_b=p_Y^{\!\top}bY$ is the Noether
charge of the $\lG$ generator, so $H_{\rm lab}$ carries a chemical potential for it.

\begin{theorem}[Connections commuting with $G$]
\label{bf:lG}
Let $b\in\lG$ act as in \eqref{bf:eq:w}.
\begin{enumerate}[label=(\roman*),leftmargin=2.2em]
\item $\Xi=-G/(2\kappa m)$ solves $-S\Xi-\Xi S^\dagger=-\id/m$, so $M$ is semisimple if and only if $S$ is. The argument uses only $G^2=-\id$,
$G^\dagger=G$ and $[b,G]=0$.
\item $\chi_S(\lambda)=(\lambda^2+\kappa^2)\bigl((\lambda-i\kappa)^2+w\cdot w\bigr)\bigl((\lambda+i\kappa)^2+\overline{w\cdot w}\bigr)$.
\item In the frame in which $b$ is constant, every solution is bounded if and only if $e^{-b\tau}$ is bounded, that is, if and only if $w=0$ or
$w\cdot w>0$. The frequencies are then $\kappa$ and $|\kappa\pm s|$ with $s=\sqrt{w\cdot w}$. For $w\cdot w=0\neq w$ the zero-momentum flow grows
like $\tau^2$; for $w\cdot w\notin[0,\infty)$ it grows exponentially. For a gauge-admissible $b$, $w=\mathbf a$ is real and every solution is bounded.
\item For the rotation $b=\beta(R_{12}-R_{45})=-\beta L_3$ the spectrum of $M$ is $\pm i\kappa$, $\pm i(\kappa+\beta)$, $\pm i(\kappa-\beta)$, each
double. For the boost $b=\beta K_3$ the exponents are $\pm\beta\pm i\kappa$, together with $\pm i\kappa$ from the factor $\lambda^2+\kappa^2$ of (ii): complex
rather than real, so the orbits spiral outward at the binding frequency. The growth is the boost factor $e^{-\beta K_3\tau}$ of \cref{bf:larmor}, and $K_3$ is not a gauge generator (\cref{cor:noboost}).
\end{enumerate}
\end{theorem}

\begin{proof}
(i) With $[b,G]=0$: $-(b+\kappa G)\Xi-\Xi(-b+\kappa G)=\frac1{2\kappa m}(bG-Gb+2\kappa G^2)=-\frac1m$. (ii) On $(V,G)\cong\C^3$, $S$ is
$z\mapsto i\kappa z+w\times z$, and $(w\times)^3=-(w\cdot w)(w\times)$ gives $w\times$ the eigenvalues $0,\pm is$; the real characteristic polynomial
of a complex-linear map is its complex one times the conjugate. (iii) By \cref{bf:larmor}, $X=e^{-b\tau}Y$ with $Y$ a solution at $b=0$, bounded by
(i) and (ii); conversely the zero-momentum solutions are $e^{-b\tau}e^{-\kappa G\tau}X_0$ with $e^{\kappa G\tau}$ bounded. And $e^{-b\tau}$ is bounded
exactly when $w\times$ is semisimple with imaginary eigenvalues: $s$ real, and not $w\cdot w=0\neq w$, where $N=w\times$ has
$N^2\bar w=|w|^2w\neq0=N^3$. For real $w$, $w\times$ is real antisymmetric. (iv) Here $s=\pm\beta$, respectively $\pm i\beta$; in the first case $S$ is
semisimple and (i) applies.
\end{proof}

The theorem is fixed, but its content deflates. Criterion~(iii) is a property of the frame in which $b$ is constant, not of the binding. In the
co-moving frame of \cref{bf:larmor} the binding survives every commuting connection, as the remark after \cref{prop:tear} anticipates (growth on
one worldline proves nothing), and for admissible $b$ the motion in the frame in which $b$ is constant is bounded as well. In that frame, not in the co-moving one, a rotation lifts the degeneracy of the single
frequency across the three pairs, noted at the end of \cref{sec:neutral}, into the triplet $\kappa-\beta$, $\kappa$, $\kappa+\beta$. Read as a
frame vector through the identification of the frame with $\lG$ (\cref{thm:brackets}, with $L_i=e_i$ and $K_i=e_{i+3}$), the connection
has the invariant $q=w\cdot w$; on this reading the lab-frame motion stays bounded exactly for a connection of magnetic type, $\bE_b\cdot\bB_b=0$ and
$|\bB_b|>|\bE_b|$, while the binding itself survives every commuting connection, and an admissible one has $\bE_b=0$. The Zeeman pattern labels the
levels of the binding frequency by their $\lG$ charge. On one worldline it has no gauge-invariant content, and it is observable only against a
reference that $b$ does not rotate, as the Zeeman effect is observed against the laboratory; that is a reading, but a sharp one. A commuting $b$
supplies a spectrum and no invariant scale, since $\beta/\kappa$ is free.

% --- source: tex/s17c_bfield#5 (Connections anticommuting with G, and screening)
\subsection{\texorpdfstring{Connections anticommuting with $G$, and screening}{Connections anticommuting with G, and screening}}
\label{sec:bf-anticommuting}

\begin{theorem}[Connections anticommuting with $G$]
\label{bf:pG}
Let $b\in\mathfrak p_G$ with $b\neq0$.
\begin{enumerate}[label=(\roman*),leftmargin=2.2em]
\item $(\pm b+\kappa G)^2=b^2-\kappa^2$, and $M^2=\operatorname{diag}(b^2-\kappa^2,\,b^2-\kappa^2)+N$ with
$N=\bigl(\begin{smallmatrix}0&-2b/m\\0&0\end{smallmatrix}\bigr)\neq0$. Hence $M$ is never semisimple, and some solutions grow at least linearly.
\item For the admissible counter-rotation $b=\beta(R_{12}+R_{45})$, which is $(D,-D)$ with $\mathbf d=-\beta\mathbf e_3$,
$\det(\lambda-M)=(\lambda^2+\kappa^2)^2(\lambda^2+\kappa^2+\beta^2)^4$. The exponents $\pm i\sqrt{\kappa^2+\beta^2}$ have algebraic multiplicity four and
geometric multiplicity two, and $\pm i\kappa$ are semisimple. The bounded solutions are exactly those with $\pi_1=\pi_2=\pi_4=\pi_5=0$, that is, with
no momentum on the support of $b$. The zero-momentum orbits are bounded and turn at the rate $\sqrt{\kappa^2+\beta^2}$ on the support.
\item The grading root as a connection satisfies $(\nu\cK+\kappa G)^2=(\nu^2-\kappa^2)\id$ on all of $V$.
\end{enumerate}
\end{theorem}

\begin{proof}
(i) $\{b,G\}=0$ gives both squares, and $M^2$ has off-diagonal block $(S^\dagger-S)/m=-2b/m$. With $b^2-\kappa^2=Q_s+Q_n$ (Jordan--Chevalley),
$M^2=\operatorname{diag}(Q_s,Q_s)+\bigl(\operatorname{diag}(Q_n,Q_n)+N\bigr)$ with commuting summands ($Q_n$ is a polynomial in $b$), so by uniqueness of the
decomposition $M^2$, and hence $M$, is semisimple only if $\operatorname{diag}(Q_n,Q_n)+N=0$, which is impossible for an off-diagonal $N\neq0$. A
non-semisimple generator of a Hamiltonian flow has a Jordan block at some exponent with non-negative real part. (ii) On the plane
$\langle e_3,e_6\rangle$, $b$ vanishes and $S=\kappa G$, semisimple with exponents $\pm i\kappa$, each double, by \cref{bf:lG}(i). On the support
$\langle e_1,e_2,e_4,e_5\rangle$, $b^2=-\beta^2$, so by (i) $S^2=(S^\dagger)^2=-\omega^2$ with $\omega^2=\kappa^2+\beta^2$ and, on the corresponding
eight-dimensional block of phase space, $M^2+\omega^2=N$ with $N^2=0$ and $\rank N=\rank b=4$. Hence the exponents there are $\pm i\omega$, each of
algebraic multiplicity four, $\ker(M^2+\omega^2)=\ker N$ has dimension four and splits into $\ker(M\mp i\omega)$ of dimension two each, and the bounded
data on the block are $\ker N=\{b\pi=0\}$. The zero-momentum flow is $e^{-S\tau}$ with $S^2=-\omega^2$ on the support. (iii) $\cK^2=\id$ and
$\{\cK,G\}=0$.
\end{proof}

Part~(i) is an obstruction: no purely anticommuting connection, admissible or not, is compatible with bounded normal modes. In particular the
Jordan blocks found after \cref{prop:tear} for a pairing-plane boost $\nu\sigma_j$ below the threshold occur for every $\nu\neq0$, below, at and above
it. Parts~(ii) and~(iii) are fixed. The rate changes to $\sqrt{\kappa^2+|\mathbf d|^2}$ for an admissible element and to $\sqrt{\kappa^2-\nu^2}$ for a
pairing boost (a metric connection, not a gauge field), and the momentum sector behaves as under the positive form of \cref{thm:virial}:
exponents carried by Jordan blocks, imaginary for a compact element and for a boost below the threshold.

\begin{proposition}[What a gauge transformation removes]
\label{bf:remove}
Let $b=b_l+b_p\in\lG\oplus\mathfrak p_G$ be constant. Some $g(\tau)\in O(3,3)$ carries $(b,G)$ to $(0,G')$ with $G'$ constant if and only if $[b,G]=0$.
For $b\in\mathfrak p_G$, $e^{b\tau}Ge^{-b\tau}=e^{2b\tau}G$. The gauge-covariant remainder is $\nabla_{\!\tau}G=[b,G]=2b_pG$: along a worldline $b_l$ is
pure gauge and $b_p$ is not.
\end{proposition}

\begin{proof}
$b^g=0$ forces $g=e^{-b\tau}g_0$, and then $G'=g_0^{-1}e^{b\tau}Ge^{-b\tau}g_0$ has derivative $g_0^{-1}e^{b\tau}[b,G]e^{-b\tau}g_0$. If $b$ anticommutes
with $G$, then $Gf(b)=f(-b)G$ for every power series $f$.
\end{proof}

This is fixed. A connection on a line has no curvature, since the non-abelian term $[b_\mu,b_\nu]$ needs two directions (inherited); once the
binding root is present, the invariant is $\nabla_{\!\tau}G$, which is exactly the anticommuting part. One might expect boundedness to force
$b\in\lG$. It does not, and the admissible case shows it.

\begin{theorem}[Compact screening]
\label{bf:screen}
Let $b$ be admissible, as in \eqref{bf:eq:compact}, and $\kappa>0$.
\begin{enumerate}[label=(\roman*),leftmargin=2.2em]
\item If every solution is bounded, then $\mathbf a\cdot\mathbf d=0$: the two triples rotate at equal rates, $|\mathbf a+\mathbf d|=|\mathbf a-\mathbf d|$.
\item Let $\mathbf a\cdot\mathbf d=0$ and $\mathbf a\neq0$. Then $\spec S=\{\pm i\omega_1,\pm i\omega_2,\pm i\omega_3\}$, where the $\omega_j$ are the
roots of
\[
f(\omega)=\omega^3+\kappa\omega^2-\bigl(\kappa^2+|\mathbf a|^2+|\mathbf d|^2\bigr)\,\omega-\kappa\bigl(\kappa^2+|\mathbf d|^2-|\mathbf a|^2\bigr),
\]
and if $|\mathbf a|^2\neq\kappa^2+|\mathbf d|^2$ every solution is bounded, although $[b,G]=2b_pG\neq0$ when $\mathbf d\neq0$.
\item If $\mathbf a\cdot\mathbf d=0$, $\mathbf d\neq0$ and $|\mathbf a|^2=\kappa^2+|\mathbf d|^2$, the vector
$X_0=\bigl(\mathbf n,(\mathbf a+\mathbf d)\times\mathbf n/\kappa\bigr)$ with $\mathbf n=\mathbf a\times\mathbf d$ is static, $SX_0=0$, with
$\eta(X_0,X_0)=-2|\mathbf d|^2|\mathbf n|^2/\kappa^2\neq0$, and some solutions drift linearly.
\end{enumerate}
For $\mathbf a=0$ and $\mathbf d\neq0$, \cref{bf:pG} applies.
\end{theorem}

Bounded motion is not robust under compact perturbations: the necessary condition $\mathbf a\cdot\mathbf d=0$ of (i) has codimension one in the
six-dimensional space of compact connections. For $\kappa=1$, $\mathbf a=(0,0,2)$ and $\mathbf d=(1,0,\epsilon)$ one has $\tr\bigl(\eta(S^\dagger)^2\bigr)=-8\,\mathbf a\cdot\mathbf d=-16\epsilon$,
and $M$ is not semisimple at $\epsilon=\tfrac1{10}$ although $S$ is real antisymmetric. A physical model has to explain what restricts the connection to
this set.

\begin{proof}
$S=b+\kappa G$ is real antisymmetric, since $b$ and $G$ are; so $S$ and $S^\dagger=-\eta S\eta=-b+\kappa G$ are normal with imaginary spectrum, and
by \cref{bf:criterion} the motion is bounded if and only if $P_\mu\eta P_\mu=0$ for the orthogonal spectral projectors $P_\mu$ of $S^\dagger$.
(i) Then $\Tr\bigl(\eta(S^\dagger)^2\bigr)=\sum_\mu\mu^2\Tr(P_\mu\eta P_\mu)=0$. But $\{b,G\}$ is block-off-diagonal, $\Tr\eta=0$ and
$\Tr[\mathbf v]_\times^2=-2|\mathbf v|^2$, so $\Tr\bigl(\eta(S^\dagger)^2\bigr)=\Tr(A+D)^2-\Tr(A-D)^2=-8\,\mathbf a\cdot\mathbf d$.
(ii) Let $\rho$ be the half-turn about $\mathbf a$ and $\hat\rho=\operatorname{diag}(\rho,\rho)$. Since $\rho\mathbf a=\mathbf a$ and $\rho\mathbf d=-\mathbf d$,
$\hat\rho b\hat\rho=\operatorname{diag}(A-D,A+D)=G^{-1}bG$, so $G':=G\hat\rho$ commutes with $b$ and $G$, hence with $S$ and $S^\dagger$. Moreover
$G'^2=-\id$ and $G'^\dagger=G'$, so, as for $G$ in \cref{prop:howmuch}, each eigenspace $\mathcal V_\pm$ of $G'$ (eigenvalue $\pm i$) is totally
$h_\eta$-isotropic: $-i\,h_\eta(y,y')=h_\eta(G'y,y')=h_\eta(y,G'y')=i\,h_\eta(y,y')$ on $\mathcal V_+$. In the $(\bB,\bE)$ blocks
$\mathcal V_+=\{(\mathbf x,-i\rho\mathbf x)\}$, and $S$ acts there as $\mathbf x\mapsto(\mathbf a+\mathbf d)\times\mathbf x+i\kappa\rho\mathbf x$. After a
common rotation of both triples, which commutes with $G$ and $\eta$, to $\mathbf a=|\mathbf a|\mathbf e_3$ and $\mathbf d=|\mathbf d|\mathbf e_1$, one finds
$\det\bigl(i\omega-[\mathbf a+\mathbf d]_\times-i\kappa\rho\bigr)=-i f(\omega)$. So $S$ has the eigenvalues $i\omega_j$ on $\mathcal V_+$ and $-i\omega_j$ on
$\mathcal V_-=\overline{\mathcal V_+}$; since $\eta$ exchanges $\mathcal V_\pm$, $S^\dagger$ has the same eigenvalues there. Every eigenspace of $S^\dagger$
splits along $\mathcal V_\pm$, and lies in one of them unless $f(\omega)=f(-\omega)=0$ for some $\omega$. But
$\half[f(\omega)+f(-\omega)]=\kappa\bigl(\omega^2-\kappa^2-|\mathbf d|^2+|\mathbf a|^2\bigr)$ and
$\half[f(\omega)-f(-\omega)]=\omega\bigl(\omega^2-\kappa^2-|\mathbf a|^2-|\mathbf d|^2\bigr)$, so a common root forces $\mathbf a=0$, or $\omega=0$ and
$|\mathbf a|^2=\kappa^2+|\mathbf d|^2$. (iii) With $\mathbf u=\mathbf a+\mathbf d$ and $\mathbf v=\mathbf a-\mathbf d$,
$SX_0=\bigl(\mathbf u\times\mathbf n-\mathbf u\times\mathbf n,\ \mathbf v\times(\mathbf u\times\mathbf n)/\kappa+\kappa\mathbf n\bigr)=0$, because
$\mathbf v\cdot\mathbf n=0$ and $\mathbf u\cdot\mathbf v=|\mathbf a|^2-|\mathbf d|^2=\kappa^2$; and $\eta(X_0,X_0)=|\mathbf n|^2-|\mathbf u|^2|\mathbf n|^2/\kappa^2$.
So $\mathcal E_0(S^\dagger)=\eta\,\mathcal E_0(S)$ is not $\eta$-null, and \cref{bf:criterion} fails at $\mu=0$.
\end{proof}

The theorem is fixed, and it is the case that matters for a gauge field. Its mechanism is a second binding root $G'=G\hat\rho$ commuting with $S$. The exceptional ratio $|\mathbf a|^2=\kappa^2+|\mathbf d|^2$ has
measure zero. At $\mathbf d=0$ the theorem gives the Zeeman lines of \cref{bf:lG}. For $\mathbf d\neq0$ the screened frequencies cannot be removed
along a worldline (\cref{bf:remove}), but they are $\kappa$ times functions of $|\mathbf a|/\kappa$ and $|\mathbf d|/\kappa$: screening supplies no
scale.

\begin{proposition}[Gyroscopic screening by a boost]
\label{bf:gyro}
Let $b=\alpha L_3+b_\Sigma$, where $b_\Sigma\in\mathfrak p_G$ is the boost $z\mapsto i\Sigma\bar z$ with
$\Sigma=\bigl(\begin{smallmatrix}y_1&y_2&0\\y_2&-y_1&0\\0&0&0\end{smallmatrix}\bigr)$, and let $r^2=y_1^2+y_2^2>0$. Then $[b,G]\neq0$ and
\[
\chi_S(\lambda)=(\lambda^2+\kappa^2)\bigl(\lambda^2+(\kappa-\alpha)^2-r^2\bigr)\bigl(\lambda^2+(\kappa+\alpha)^2-r^2\bigr).
\]
If $\alpha\neq0$ and $r<\min(|\kappa-\alpha|,|\kappa+\alpha|)$, every solution is bounded; for $\alpha=0$ and for $|\kappa-\alpha|<r<|\kappa+\alpha|$ some
are not. Numerically, near bounded points at which the five conditions of the proof are independent, the set of $b\in\so{3,3}$ with all motion
bounded appears to be a $10$-dimensional submanifold, and at the tested points of $\lG$ and of this family the conditions are degenerate. Both
statements rest on finite-difference Jacobian ranks, and rank loss of a defining map does not by itself make its zero set singular (as for
$x^2=0$); they are numerical observations until an exact rank and local-geometry argument is given.
\end{proposition}

\begin{proof}
On $\langle e_1,e_2,e_4,e_5\rangle$, $K_3=GL_3$ is an $\eta$-antisymmetric involution commuting with $G$, $L_3$ and $b_\Sigma$, so its eigenplanes
$\mathcal U_\pm$ are totally $\eta$-null and invariant. There $L_3=\mp G$, so $S=(\kappa\mp\alpha)G+b_\Sigma$ and $S^\dagger=(\kappa\pm\alpha)G-b_\Sigma$, each
squaring to a multiple of $\id$ ($b_\Sigma^2=r^2$, $\{b_\Sigma,G\}=0$), negative under the stated bound on $r$. For $\alpha\neq0$ the restrictions of
$S^\dagger$ to $\mathcal U_+$ and to $\mathcal U_-$ have no common eigenvalue, so every eigenspace of $S^\dagger$ lies in $\mathcal U_+\oplus\langle e_3,e_6\rangle$ or in
$\mathcal U_-\oplus\langle e_3,e_6\rangle$ (complexified). Since $\mathcal U_\pm$ are totally $\eta$-null and $\eta$-orthogonal to the $b$-free plane
$\langle e_3,e_6\rangle$, on which the eigenspaces of $\kappa G$ are isotropic lines, $h_\eta$ vanishes on every eigenspace, and \cref{bf:criterion}
applies. For $\alpha=0$, $b\in\mathfrak p_G$ and
\cref{bf:pG} applies; for $|\kappa-\alpha|<r<|\kappa+\alpha|$, $\chi_S$ has a real root. For the dimension statement: at simple frequencies
\cref{bf:criterion}(ii) imposes five real conditions on $b$ (two for imaginary spectrum after the trace relation, three for isotropy). Their
finite-difference Jacobian has rank five at four least-squares solutions and at an exact compact screened point, and rank four, with residuals of
second order along the kernel directions, at three points of this family and at two elliptic points of $\lG$.
\end{proof}

The family is fixed; the dimension statements are numerical. This is the gyroscopic stabilisation of a saddle
\cite{ThomsonTait1879,Kirillov2013}: the rotation detunes the two light-like planes of $K_3$ in which the boost acts, so that
$\eta(\pi,\pi)=2\eta(\pi_+,\pi_-)$, with $\pi_\pm$ the components of $\pi$ in the null planes $\mathcal U_\pm$ of the proof, averages to zero. Since $b_\Sigma$ is a boost, the family is not a gauge field. Together with \cref{bf:screen} it
shows that an anticommuting part need not cause secular growth, and that boundedness does not single out $\lG$; gauge removability does
(\cref{bf:remove}).

\begin{proposition}[A connection alone does not bind]
\label{bf:alone}
\begin{enumerate}[label=(\roman*),leftmargin=2.2em]
\item Let $\kappa=0$ and let $b\in\so{3,3}$ be any constant connection, $b=\nu\cK$ included. Then $X=e^{-b\tau}Y$ gives
$\mathcal L=\tfrac m2\eta(\dot Y,\dot Y)$ and $\ddot Y=0$. The generator $M=\bigl(\begin{smallmatrix}-b&\id/m\\0&-b\end{smallmatrix}\bigr)$ has
$\det(\lambda-M)=\chi_{-b}(\lambda)^2$ and is never semisimple: the motion is free drift, and there is no binding frequency.
\item The rescaling $(\tau,\kappa,b,m)\mapsto(\vartheta\tau,\kappa/\vartheta,b/\vartheta,\vartheta m)$ maps $M$ to $M/\vartheta$, solutions to solutions and
the action $\int\mathcal L\,d\tau$ to itself, and leaves $\eta$, $G$, $\lG$, $\mathfrak p_G$ and $\so4_\Delta$ of \cref{thm:PSbinding} unchanged. No value of $\kappa$, hence of
$\ell=c/\kappa$, is selected, with or without a constant $b$.
\item $M(m)=TM(1)T^{-1}$ with $T=\operatorname{diag}(\id,m\id)$: the spectrum of $M$ does not involve $m$, and the rescaling of (ii) keeps $m\kappa$,
not $m/\kappa$, fixed.
\end{enumerate}
\end{proposition}

\begin{proof}
(i) $\dot X+bX=e^{-b\tau}\dot Y$. Write $-\operatorname{diag}(b,b)=S'+N'$ (Jordan--Chevalley). Then $M=S'+(N'+N)$ with
$N=\bigl(\begin{smallmatrix}0&\id/m\\0&0\end{smallmatrix}\bigr)$, and $N'+N\neq0$ because $N$ is off-diagonal and $N'$ is block-diagonal. (ii) and (iii) are
direct.
\end{proof}

Part~(i) is inherited, since every connection on a line is pure gauge, and its consequence is an obstruction: a Yang--Mills potential on a
worldline, with no self-adjoint shift, produces free motion. Part~(ii) is dimensional analysis, inherited; it restates that $\kappa$ is not
fixed (end of \cref{sec:neutral}) and that no ratio among the inertia constants and $m$ is fixed by the Clifford structure (after
\cref{prop:several}). Part~(iii) is fixed and elementary: nothing in the interior dynamics ties $m$ to $\kappa$, and $mc^2=\hbar\kappa$, that is
$\ell=\hbar/(mc)$, is a matching added from outside, never a result.

% --- source: tex/s17c_bfield#6 (The split virial identity and the vacuum)
\subsection{The split virial identity and the vacuum}
\label{sec:bf-vacuum}

The split virial identity of \cref{thm:splitvirial} holds unchanged with the connection: for every constant $S$, and so for every
$b\in\so{3,3}$, $\tfrac d{d\tau}\eta(X,\pi)=\eta(\pi/m-SX,\pi)+\eta(X,S^\dagger\pi)=\eta(\pi,\pi)/m$ along every solution. The quantum statement
needs the signature of the energy on each frequency.

\begin{theorem}[Krein signatures and the zero-point energy]
\label{bf:krein}
Suppose every solution is bounded, and let $\omega>0$ have multiplicity $k_\omega$ in $\spec(-iS)$. On the eigenspace $\mathcal E_{i\omega}$ of the
generator, the energy form $v\mapsto\bar v^{\!\top}\mathcal Hv$ has signature $(k_\omega,k_\omega)$. Hence, with positive-norm quantisation of the normal
modes, the zero-point energy $\sum_j(\pm)\half\hbar\omega_j$ cancels frequency by frequency, and every level has a lower one.
\end{theorem}

\begin{proof}
Work in the canonical variables of \cref{bf:M}. On $\mathcal E_{i\omega}$, $\mathcal Hv=-i\omega\mathbb Jv$, so the energy form equals $\omega$ times the
Krein form $-i\bar v^{\!\top}\mathbb Jv$, which is nondegenerate there because eigenspaces of distinct imaginary eigenvalues are orthogonal for it
\cite{YakubovichStarzhinskii75}. The zero-momentum part of $\mathcal E_{i\omega}$ is $\mathcal E_{i\omega}(-S)$, of dimension $k_\omega$, and $H$ vanishes on
it (\cref{bf:M}). Since the generator is semisimple, $\dim\mathcal E_{i\omega}=2k_\omega$ by \cref{bf:M}. A nondegenerate Hermitian form with a
half-dimensional isotropic subspace has signature $(k,k)$, and a mode of energy sign $\pm$ quantises to $\pm\hbar\omega(n+\half)$.
\end{proof}

The theorem is fixed, although structural: it holds for any nondegenerate $\eta$. It extends the vanishing of the zero-point energy on
$\R^{3,3}$ at $b=0$ (\cref{thm:subfs}) and the absence of a ground state (\cref{mb:noground}) to every connection that keeps the motion bounded: the
cancellation, forced by the perfect-square form \eqref{bf:eq:L}, survives the splitting line by line, and so does the absence of a ground state.
On the coordinate subfields of \cref{thm:subfs}, where the zero-point energy is $\pm\tfrac32\hbar\kappa$ or $\pm\half\hbar\kappa$, the theorem does not
apply, since $G$ does not preserve the subfield and the zero-momentum flow leaves it (\cref{thm:subfs}).

\begin{corollary}[No Fock vacuum with Jordan blocks]
\label{bf:nofock}
Suppose $M$ is not semisimple at a frequency $\omega\neq0$, as for the counter-rotation of \cref{bf:pG}(ii) and for a pairing boost below the threshold. Then no linear symplectic change of
variables writes the Hamiltonian on that block as a sum of signed oscillators, and there is no Fock vacuum annihilated by normal-mode lowering
operators.
\end{corollary}

\begin{proof}
A sum of signed oscillators generates a semisimple flow, and symplectic changes of variables act on the generator by similarity.
\end{proof}

This is an obstruction. A model block with the same Jordan structure is $H=\omega(q_1p_2-q_2p_1)+\half(q_1^2+q_2^2)$, whose two terms
Poisson-commute \cite{Williamson36}; that its quantisation has continuous spectrum is inherited and is not re-derived here.

% --- source: tex/s17c_bfield#7 (Three constituents: the binding-compatible algebra so(4)Delt)
\subsection{\texorpdfstring{Three constituents: the binding-compatible algebra $\so4_\Delta$}{Three constituents: the binding-compatible algebra so(4)Delta}}
\label{sec:bf-three}

At three constituents we use the generators $E_A$ of \eqref{eq:tower}, the numbering of \cref{sec:closure} (the temporal directions, their
partners $i'=i+3$, the weak directions $W=\{7,13,14,15\}$, the colour directions $\{10,11,12,16,17,18\}$ and the span $E$), $G_3=G\oplus G\oplus G$,
and $\mathrm{rot}(i,j)$ acting by $\tfrac12\ad$; $R_{AB}$ is defined as at one constituent. By \cref{thm:PSbinding}, the part of the Pati--Salam
algebra of $P_*$ that commutes with $G_3$ is $\so4_\Delta=\operatorname{span}\{\mathrm{rot}(i,j):i<j\in W\}\cong\su2\oplus\su2$, compact, and neither
Pati--Salam factor meets $\Cent(G_3)$. The spectra below use an explicit splitting.

\begin{lemma}[The two factors of $\so4_\Delta$]
\label{bf:su2}
The elements $T^\pm_1=\half\bigl(\mathrm{rot}(13,14)\pm\mathrm{rot}(7,15)\bigr)$, $T^\pm_2=\half\bigl(\mathrm{rot}(7,14)\mp\mathrm{rot}(13,15)\bigr)$ and
$T^\pm_3=\half\bigl(\mathrm{rot}(7,13)\pm\mathrm{rot}(14,15)\bigr)$ satisfy $[T^+_i,T^-_j]=0$ and $[T^\pm_1,T^\pm_2]=-T^\pm_3$ cyclically, and the trace
form is negative definite on their span, so $\so4_\Delta=\su2_+\oplus\su2_-$. On the eighteen vectors each $T^\pm_3$ has the eigenvalues $\pm\tfrac i2$,
four times each, and $0$, ten times.
\end{lemma}

\begin{proof}
Each $\mathrm{rot}(i,j)$ with $i,j\in W$ rotates the plane $(e_i,e_j)$ and its partner plane $(e_{i'},e_{j'})$ alike, so on $\operatorname{span}W$, and
identically on the partner span $\{10,16,17,18\}$, the six elements act as the standard generators of $\so4$, and the $T^\pm_i$ are their
self-dual and anti-self-dual combinations; the brackets follow by direct computation. Each $T^\pm_3$ is half a sum of two commuting plane rotations
on each of the two four-dimensional spans and vanishes on the remaining ten directions.
\end{proof}

The lemma is fixed; it makes the isomorphism $\so4_\Delta\cong\su2\oplus\su2$ of \cref{thm:PSbinding} explicit. Each $\su2$ factor locks a weak rotation of $W$ to the colour rotation of the partners $\{10,16,17,18\}$. Reading $\so4_\Delta$ as Yang and Mills'
isospin algebra, doubled, is a reading, as is the description after \cref{thm:survive}.

\begin{proposition}[Spectra on the eighteen-dimensional frame]
\label{bf:so4spec}
Let $\kappa=5$ and $m=1$.
\begin{enumerate}[label=(\roman*),leftmargin=2.2em]
\item For $b=\beta T^+_3$ with $\beta=2$, $\chi_S=(\lambda^2+25)^5(\lambda^2+16)^2(\lambda^2+36)^2$: the frequency $\kappa$ on the five complex singlet
directions, and $\kappa\mp\beta/2$ on the two $\su2_+$ doublets, of weights $\pm\half$.
\item For $b=\beta_+T^+_3+\beta_-T^-_3$ with $(\beta_+,\beta_-)=(2,4)$, the lines are $\kappa\pm\half(\beta_+\pm\beta_-)=8,6,4,2$, together with $\kappa$.
\item In both cases $\Xi=-G_3/(2\kappa m)$ solves the Sylvester equation, $M$ is semisimple, the Krein signatures are $(k,k)$, and the zero-point
energy vanishes.
\item The weak-only rotation $R_{7,13}$ and the colour-only rotation $R_{10,16}$ give a non-semisimple $M$, hence secular growth.
\item $b=R_{7,13}\pm R_{10,11}\in\mathfrak{ps}_W$ has $[b,G_3]\neq0$, and every solution is bounded, for every $\kappa>0$. At $\kappa=5$ the frequencies
are $\kappa$ (six pairs) and $|\omega_j|\approx4.319,\,5.050,\,5.731$, with $\omega_j$ the roots of $\omega^3+5\omega^2-26\omega-125$.
\end{enumerate}
\end{proposition}

\begin{proof}
(i)--(iii) The connection lies in $\so4_\Delta\subset\Cent(G_3)$ and is real antisymmetric, so \cref{bf:lG}(i), whose argument uses only
$G_3^2=-\id$, $G_3^\dagger=G_3$ and $[b,G_3]=0$, gives semisimplicity and boundedness, and \cref{bf:krein} gives the signatures and the zero-point
energy. Since $b$ commutes with $G_3$, $S=b+\kappa G_3$ is complex linear on $(\R^{18},G_3)\cong\C^9$ and acts as $i\kappa+b$. By the proof of \cref{bf:su2}, on the four
complex directions of $W$ the commuting $T^+_3,T^-_3$ have the joint eigenvalues $\pm\tfrac i2(1,1)$ and $\pm\tfrac i2(1,-1)$, and both vanish on the other five; the
characteristic polynomials follow as in the proof of \cref{bf:lG}(ii). (iv) and (v) $G_3$ maps the six-plane $\langle e_7,e_8,e_{13};e_{10},e_{11},e_{16}\rangle$ as
$G$ maps $\langle e_1,e_2,e_3;e_4,e_5,e_6\rangle$, and each $b$ named preserves that six-plane and vanishes on its $\eta$-orthogonal complement, where
the motion is the binding at $b=0$. On the six-plane $b$ is admissible. $R_{7,13}$ rotates the temporal triple alone and $R_{10,16}$ the spatial
triple alone, so $\mathbf a\cdot\mathbf d=\pm|\mathbf a|^2\neq0$, and \cref{bf:screen}(i) applies; since $S$ is real antisymmetric, the unbounded
solutions grow polynomially. For $R_{7,13}\pm R_{10,11}$ the rotation vectors of the two triples are the unit vectors $\mathbf e_2$ and $\pm\mathbf e_3$:
equal rates and perpendicular axes, so $\mathbf a=\half(\mathbf e_2\pm\mathbf e_3)\perp\mathbf d=\half(\mathbf e_2\mp\mathbf e_3)$ and
$|\mathbf a|^2=|\mathbf d|^2=\half\neq\kappa^2+\half$. \Cref{bf:screen}(ii) applies, with $f$ as stated at $\kappa=5$; the three frequencies are its
roots, evaluated numerically.
\end{proof}

This is fixed. Every constant $\so4_\Delta$-valued connection is compact and commutes with $G_3$, so it keeps the interior bounded; the splitting
follows the $\su2$ weights, $\beta/\kappa$ is free, and by \cref{bf:larmor} the pattern can be removed along one worldline. The single-factor
rotations tear the binding, but boundedness does not single out the locked diagonal: the connection of part~(v) is bounded with $[b,G_3]\neq0$,
while the anti-locked $R_{7,13}+R_{10,16}$, which has $\mathbf a=0$ on its six-plane, is not. What singles out $\so4_\Delta$ is commutation with
$G_3$ (\cref{thm:PSbinding}), that is, removability along a worldline (\cref{bf:remove}).

% --- source: tex/s17c_bfield#8 (Quantum Yang--Mills facts, and the channels of the vacuum)
\subsection{Quantum Yang--Mills facts, and the channels of the vacuum}
\label{sec:bf-quantum}

Classically the connection splits, screens or tears the binding but supplies no scale. The only place left for it to generate one is the
quantum theory, and the following inherited facts are needed to say what that would involve.

\begin{proposition}[Inherited facts]
\label{bf:qym}
\begin{enumerate}[label=(\roman*),leftmargin=2.2em]
\item For $\su4$ with $\tr(T^AT^B)=\half\delta^{AB}$, $\sum_AT^AZT^A=\half\bigl(\tr(Z)\id-\tfrac14Z\bigr)$ for every $4\times4$ matrix $Z$: the factor is
$\tfrac{15}8=C_F$ on $\id$ and $-\tfrac18$ on traceless $Z$.
\item The channel factors $T_1\cdot T_2=\half[C_2(\text{pair})-C_2(R_1)-C_2(R_2)]$ are $+\tfrac38,-\tfrac58$ on $\mathbf{10},\mathbf6\subset\mathbf4\otimes\mathbf4$,
$-\tfrac{15}8,+\tfrac18$ on $\mathbf1,\mathbf{15}\subset\bar{\mathbf4}\otimes\mathbf4$, and $+\tfrac14,-\tfrac34$ on $\mathbf3,\mathbf1\subset\mathbf2\otimes\mathbf2$ of
$\su2$, with $C_2=\tfrac{15}8,\tfrac92,\tfrac52,4$ on $\mathbf4,\mathbf{10},\mathbf6,\mathbf{15}$. Hence the Dirac bilinear
$\overline{(\mathbf4,\mathbf2,\mathbf1)}\otimes(\mathbf4,\mathbf1,\mathbf2)=(\mathbf1,\mathbf2,\mathbf2)\oplus(\mathbf{15},\mathbf2,\mathbf2)$ has $\su4$ factors
$-\tfrac{15}8$ and $+\tfrac18$, and a Lorentz-scalar Majorana bilinear of $(\mathbf4,\mathbf1,\mathbf2)$, symmetric in its internal indices, lies in
$\mathrm{Sym}^2(\C^4\otimes\C^2)=(\mathbf{10},\mathbf1,\mathbf3)\oplus(\mathbf6,\mathbf1,\mathbf1)$, of dimensions $30+6$, with factors
$(+\tfrac38,+\tfrac14)$ and $(-\tfrac58,-\tfrac34)$.
\item In Euclidean momentum space
\[
\int_{k^2\le k_{\max}^2}\frac{d^4k}{(2\pi)^4}\frac1{k^2+M_{\rm dyn}^2}=\frac1{16\pi^2}\bigl[k_{\max}^2-M_{\rm dyn}^2\ln(1+k_{\max}^2/M_{\rm dyn}^2)\bigr].
\]
The function
$f_{\rm gap}(x)=1-x^2\ln(1+x^{-2})$ decreases strictly from $1$ to $0$, so the contact gap equation $1=\lambda_{\rm c}f_{\rm gap}(x)$ has a root
$x=M_{\rm dyn}/k_{\max}>0$, with $M_{\rm dyn}$ the dynamical mass, $k_{\max}$ the cutoff and $\lambda_{\rm c}$ the dimensionless contact coupling, if and only if $\lambda_{\rm c}>1$.
\item The ladder critical coupling $\pi/(3C_F)$ \cite{FukudaKugo76,Miransky93} is $8\pi/45\approx0.5585$ at $C_F=\tfrac{15}8$. For $SU(2)$ with $N_f$
Dirac doublets the one-loop coefficient is $\beta_0=(22-2N_f)/3$, equal to $6$ at $N_f=2$ \cite{GrossWilczek73,Politzer73}, and
$\Lambda=\mu\,e^{-8\pi^2/(\beta_0g^2(\mu))}$ is independent of $\mu$. The breaking $SU(4)\to Sp(4)$ of $N_f=2$ leaves $15-10=5$ Goldstone
directions \cite{Peskin80}.
\end{enumerate}
\end{proposition}

All of this is inherited \cite{NambuJonaLasinio61,Slansky}; the strict decrease in (iii) follows from $\ln(1+u)>u/(1+u)$. The use of (ii) is a reading.
Under the most-attractive-channel heuristic of one-boson exchange \cite{RabyDimopoulosSusskind80}, in which a negative factor means attraction, and
with a common coupling for $\su4$ and $\su2_R$, which this manuscript adopts only as an assumption (\cref{ass:field} and the remark after
\cref{prop:coupled3}), the exchange favours the Dirac singlet channel $(\mathbf1,\mathbf2,\mathbf2)$ and the colour-breaking $(\mathbf6,\mathbf1,\mathbf1)$.
It disfavours the $(\mathbf{15},\mathbf2,\mathbf2)$ that the Georgi--Jarlskog factor $-3$ needs (\cref{prop:EW}; \cite{GeorgiJarlskog79}), and the
$(\mathbf{10},\mathbf1,\mathbf3)$ that carries the Standard-Model vacuum of \cref{sec:selection}. That vacuum is a minimum of a scalar potential, not a
gauge-exchange condensate, so the heuristic addresses a different, dynamical scenario, in which a gauge-generated vacuum would point away from the
one selected there. In the selected vacuum the order is moreover the reverse of a mass generated by the gauge field: the Pati--Salam vector bosons
outside the Standard Model acquire their masses by eating nine Goldstone directions (\cref{thm:nomajoron}).

% --- source: tex/s17c_bfield#9 (What the connection part supplies)
\subsection{What the connection part supplies}
\label{sec:bf-synthesis}

The results of this section combine as follows.
\begin{enumerate}[label=(\alph*),leftmargin=2.2em]
\item The binding frequency comes from the self-adjoint part of the shift, never from a connection. The rate $\kappa$, with
$(\kappa G)^2=-c^2/\ell^2$, sits in the tensor part (\cref{bf:covariant}, fixed), and a connection alone gives free drift (\cref{bf:alone}, an
obstruction).
\item A commuting connection splits the binding frequency without supplying a scale: the lines $\kappa$ and $\kappa\pm\sqrt{w\cdot w}$ are bounded for
admissible $b$ and removable along any worldline, with $\beta/\kappa$ free (\cref{bf:lG,bf:larmor}, fixed; the Zeeman description is a reading). At
three constituents $\so4_\Delta$ gives lines labelled by $\su2$ weights and is not the only bounded part of the Pati--Salam algebra
(\cref{bf:so4spec}), and the zero-point cancellation survives line by line (\cref{bf:krein}).
\item A purely anticommuting connection tears the binding (\cref{bf:pG,bf:nofock}, obstructions), with rates $\sqrt{\kappa^2+|\mathbf d|^2}$ for an
admissible element or $\sqrt{\kappa^2-\nu^2}$ for a pairing boost, carried by Jordan blocks. A mixed admissible connection keeps every orbit bounded
exactly at equal rates of the two triples about axes that are not opposite, outside the ratio $|\mathbf a|^2=\kappa^2+|\mathbf d|^2$
(\cref{bf:screen}); a boost can be screened gyroscopically (\cref{bf:gyro}); neither gives a scale.
\item The only role left for the connection in generating a scale is quantum, and it is open. The arena leaves $\ell$ free (\cref{bf:alone}, and
\cref{sec:restclosure} for a dynamical rate), and classical four-dimensional Yang--Mills theory has no dimensionful parameter; a scale arises only
by dimensional transmutation, $\Lambda_b=\mu\,e^{-8\pi^2/(\beta_0g^2(\mu))}$ \cite{ColemanWeinberg73,GrossWilczek73}, for a Yang--Mills theory whose
potential is $b$.
\end{enumerate}

\begin{conjecture}[A transmuted binding length]
\label{bf:transmute}
$\ell=\hbar c/\Lambda_b$ for the scale $\Lambda_b$ of a Yang--Mills theory whose potential is the connection part of the shift; that is,
$\hbar\kappa=\Lambda_b$.
\end{conjecture}

The conjecture is open: nothing in the arena derives it. Reading the rest energy $mc^2$ as $\hbar\kappa$ is a separate identification and not part of
the conjecture, since the spectrum of $M$ does not involve $m$ and the rescaling of \cref{bf:alone} keeps $m\kappa$, not $m/\kappa$, fixed. A
mechanism would need $\kappa$ to be dynamical, as the rotor rate of \cref{sec:closure} is, and coupled to a quantum gauge sector with a fixed
$\beta_0$ and a matching coupling; even then, by \cref{bf:qym}(ii), the most attractive channels are not those of the vacuum selected in
\cref{sec:selection}. A classification of the compact connections that keep the binding bounded at three constituents, beyond the commuting ones and the equal-rate
family embedded as in \cref{bf:so4spec}(v), is also open, as is a description of the bounded set in $\so{3,3}$ beyond the numerical statement of
\cref{bf:gyro}.

Read in the variables of $E=mc^2$, and only through the added identification $mc^2=\hbar\kappa$, the results of this section complete the reading
of \cref{sec:energy-readings}, in which the self-adjoint summand $\kappa G$ carries $c$ and alone binds; what follows is a reading of the fixed
results above. The Yang--Mills potential is the antisymmetric, transport summand of the same shift, and it is compact. Classically it generates neither $m$ nor $\kappa$: when it commutes with $G$ it splits the interior rate $\kappa$
into charge-labelled lines, a purely anticommuting potential tears the binding, and a mixed one may or may not (\cref{bf:screen}). Whether
its quantum scale is the length $\ell$ that the arena leaves free is open.

\section{Gauge transport of the two energies}
\label{sec:w:gaugetransportofthetwoen}

% --- source: tex/s14b_energy#11 (What transport by a compact gauge field does to the two ener)
\subsubsection{What transport by a compact gauge field does to the two energies}
\label{sec:energy-transport}

An admissible Yang--Mills generator is compact: by \cref{thm:sig} and \cref{cor:noboost} gauge generators lie in $\so3_T\oplus\so3_S\subset\so{3,3}$,
the rotations of the two triples. Write such a generator as $b=(A+D,A-D)$, with $A+D$ acting on the temporal and $A-D$ on the spatial triple. Its
diagonal part $(A,A)$ lies in $\lG$ and its anti-diagonal part $(D,-D)$ in $\mathfrak p_G$. Boost-valued connections are legitimate metric
connections (\cref{prop:tear}) but not admissible gauge potentials.

\begin{proposition}[What transport preserves]
\label{en:transport}
\begin{enumerate}[label=(\roman*),leftmargin=2.2em]
\item An element $Y\in\so{3,3}$ preserves $\bE\cdot\bB$ for all $X$, equivalently $q$, equivalently $E_0$, iff $[Y,G]=0$, that is iff $Y\in\lG$.
\item A compact $b=(A+D,A-D)$ preserves $E_6$ and $\Re q=|\bB|^2-|\bE|^2$, and it preserves $E_0$ for all $X$ iff $D=0$. The anti-diagonal
rotation through $\tfrac\pi4$ about $e_3$ carries the null field $(1,i,0)$ to the field $\tfrac{1+i}{\sqrt2}(1,1,0)$ at rest, $E_0$ going from $0$ to $1$
at $E_6=1$.
\item At three constituents the part of the Pati--Salam algebra of $P_*$ commuting with $G_3$, the algebra $\so4_\Delta\cong\su2\oplus\su2$ of
\cref{thm:PSbinding}, consists of paired rotations, so it preserves the $\so{9,\C}$ invariants $\tfrac12|Z|^2=\tfrac12\sum_k|z^{(k)}|^2$ and
$q_3(Z)=\sum_kz^{(k)}\cdot z^{(k)}$, hence $\tfrac12|q_3|$.
\item It preserves neither the rest energies of the constituents nor that of their summed density: the rotation through $\tfrac\pi4$ generated by
$\tfrac12\mathrm{rot}(7,13)\in\so4_\Delta$ maps $(z^{(1)},z^{(2)},z^{(3)})=(0,(1,i,0),(1,0,0))$ to $(0,(0,i,0),(\sqrt2,0,0))$, keeping
$\tfrac12|Z|^2=\tfrac32$ and $q_3=1$, while $\sum_kE_0^{(k)}$ goes from $\tfrac12$ to $\tfrac32$ and
$\bigl(\sum_kE_6^{(k)};\sum_k\mathbf S_{\rm t}^{(k)}\bigr)$ goes from $(\tfrac32;0,0,1)$ to $(\tfrac32;0,0,0)$, its Lorentz length from $\tfrac{\sqrt5}2$
to $\tfrac32$.
\end{enumerate}
\end{proposition}

\begin{proof}
(i) For $Y\in\so{3,3}$, $Y^{\!\top}\eta G+\eta GY=\eta[G,Y]$, so $Y$ preserves the form $\eta G$, that is $\bE\cdot\bB$, iff $[Y,G]=0$. Every $Y\in\so{3,3}$
preserves $\Re q$; if it also preserves $|q|$, then $\Im q\cdot\delta(\Im q)=0$, so $\delta(\Im q)=0$ because polynomials form an integral domain. (ii)
$b$ rotates each triple orthogonally, so it preserves $|\bB|^2$ and $|\bE|^2$; $(A,A)$ commutes with $G$ and $G(D,-D)G^{-1}=(-D,D)$, so by (i) $E_0$ is
preserved for all $X$ iff $D=0$. In the example $\bB=(1,0,0)\mapsto\tfrac1{\sqrt2}(1,1,0)$ and $\bE=(0,1,0)\mapsto\tfrac1{\sqrt2}(1,1,0)$. (iii) By
\cref{thm:so9C} the paired rotations are the elements $\bigl(\begin{smallmatrix}R&0\\0&R\end{smallmatrix}\bigr)$ with $R$ real antisymmetric, which act on
$Z\in\C^9$ by a real orthogonal matrix and preserve $|Z|^2$ and $q_3$; by \cref{thm:PSbinding} $\so4_\Delta$ is spanned by such elements. (iv)
$\mathrm{rot}(7,13)$ rotates the first components of $z^{(2)}$ and $z^{(3)}$ into each other by a real rotation, and the stated values follow by
direct evaluation, with $\mathbf S_{\rm t}=(0,0,1)$ for $(1,i,0)$.
\end{proof}

Parts (i) and (ii) are inherited linear algebra, and they cut both ways. Pointwise, the anti-diagonal part of an admissible $b$ does turn a null
field into one at rest at fixed undivided energy, while the diagonal part, a rotation in $\lG$, changes no invariant. But the anti-diagonal part is
the part that fails to commute with $G$: along a worldline it cannot be removed by a rotating frame, unlike the diagonal part (\cref{bf:remove}; Larmor's theorem
\cite{Larmor1897}), and in the interior it does not commute with the shift $\kappa G$. There, alone, it produces Jordan blocks and secular growth
(\cref{bf:pG}), and a mixed compact connection keeps every orbit bounded exactly when it rotates the two triples at equal rates about axes that
are not opposite, outside one exceptional ratio of the rates to $\kappa$ at which some solutions drift (\cref{bf:screen}); these are statements
about the interior transport, not about the pointwise invariants, and they are not used here. Nothing in the proposition selects an amount of rest energy. Parts (iii) and (iv) are an
obstruction for the $\so{9,\C}$ invariants only: no transport valued in the paired rotations changes $\tfrac12|Z|^2$ or $q_3$. They are no obstruction
for the rest energy of the whole in the diagonal Lorentz reading, which the example changes; $\mathrm{rot}(7,13)$ does not commute with the diagonal
rotations, and which of the two quantities is physical depends on the Lorentz action at three constituents, which the algebra leaves open
(\cref{tab:choices}; by \cref{thm:nocommon} the closure symmetry meets the interior Lorentz algebra of the first constituent in $0$). Moreover
$\so4_\Delta$ is the part of the Pati--Salam algebra commuting with $G_3$, not the only part whose transport keeps orbits bounded (\cref{bf:so4spec}(v)). Here $b$ enters
only as a constant transport element or a holonomy, with no field strength and no energy of its own, so where a gauge field belongs in the
differentiation of the energy is, at this level, a matter of reading.

% --- source: rest_energy/d7_between#0 ((head))
% Track d7 -- The b-field between manifolds.

% --- source: rest_energy/d7_between#1 (The b-field between manifolds)
\section{\texorpdfstring{The $b$-field between manifolds}{The b-field between manifolds}}
\label{d7:sec:main}

The author's first remark is that the $b$-particle is a deal breaker when two manifolds or more are in question. \Cref{d5:thm:bfield2}
tested it for two binding frames and one field, the copy rotation. This section asks the question for any number of manifolds and any
field that can act between them, and answers it in five steps. First, which fields exist depends on how the Lorentz algebra acts: a common
binding action admits the complex orthogonal copy algebra $\so{n,\C}$, a common interior action only its compact part, the copy rotations $\so n$, and a spectator action decouples the Lorentz-carrying manifold from every
binding-compatible field (fixed, partly restated from \cref{sec:closure}; \cref{d7:sec:exist}). Second, the closure itself admits such a field: with a
common binding action the rotor may act by any element of $\uu n$, which is a copy rotation plus a matrix of rates (constructed;
\cref{d7:sec:rotor}). Third, every binding-compatible field turns the frame flow into a Schr\"odinger flow on $\C^{3n}$, unitary for the
undivided energy exactly when the field is compact. The rates are gauge invariants, the field redistributes the frequencies at fixed sum, and
it is removable exactly when it preserves the rate eigenspaces (fixed; \cref{d7:sec:schr,d7:sec:boost}). Fourth, three manifolds carry an
invariant that two do not: the sign of $B_{12}B_{23}B_{31}$, odd under interior time reversal, the analogue of Jarlskog's invariant (fixed;
\cref{d7:sec:three}). Fifth, on two and three manifolds under a common Lorentz action the copy isospin is integral: frames carry isospin one, never one-half. The Yang--Mills doublet acting on two copies
of the undivided energy is a connection in one of its three directions and a rate term in the other two, and on three manifolds doublets exist only under the
spectator action, where unequal internal rates break the $\so4_\Delta$ of \cref{thm:PSbinding} to a diagonal $\so3$ (fixed; \cref{d7:sec:isospin}). The verdict
is in \cref{d7:sec:assess}. Identifiers d7-$n$ refer to the \texttt{PASS} lines of \texttt{d7\_bfield\_between.py} (\cref{d7:tab:claims}).

\paragraph{Notation.} As in \cref{d5:sec:main}, with $n$ copies: $V_n=\R^n\otimes\R^6$ in the copy-major basis, $\eta_n=\id_n\otimes\eta$,
$G_n=\id_n\otimes G$, and $z=(z_1,\dots,z_n)\in\C^{3n}$ with $z_k=\bB_k+i\bE_k$, on which $G_n$ acts as $i$. $E_{kl}$ are the matrix units on the copy
index. The \emph{copy rotation} $C^{kl}=(E_{lk}-E_{kl})\otimes\id_6$ maps copy $k$ to copy $l$, and the \emph{copy boost} is $(E_{lk}-E_{kl})\otimes G$. A
\emph{copy field} is $b=B\otimes\id_6$ with $B\in\so n$ real antisymmetric, and $\plaq :=B_{12}B_{23}B_{31}$ at $n=3$. The binding term with
rates $\kappa_k\ge0$ is $K\otimes G$ with $K=\operatorname{diag}(\kappa_1,\dots,\kappa_n)$, and $S=K\otimes G+b$, $\dot X=-SX+\pi/m$, $\dot\pi=S^\dagger\pi$, with
$M$ as in \cref{d5:thm:bfield2}. A field is \emph{binding-compatible} if it commutes with $G_n$. Rates are rates, and reading one as a rest
energy needs $mc^2=\hbar\kappa$ (\cref{cl:norel}).

% --- source: rest_energy/d7_between#2 (Which fields can act between manifolds)
\subsection{Which fields can act between manifolds}
\label{d7:sec:exist}

\begin{proposition}[Inter-manifold fields and the Lorentz action]
\label{d7:prop:exist}
\begin{enumerate}[label=(\roman*),leftmargin=2.2em]
\item \emph{Common binding action} (the diagonal $\lG$). The elements of $\so{3n,3n}$ that commute with it and with $G_n$ are
$\{B\otimes\id_6+B'\otimes G:\ B,B'\in\so n\}\cong\so{n,\C}$, of dimension $0,2,6,12$ for $n=1,\dots,4$. The compact part consists of the copy
rotations, of dimension $0,1,3,6$; the copy boosts $B'\otimes G$ are noncompact.
\item \emph{Common interior action} (the diagonal $\lint$, the action $\Delta$ of \cref{sec:actions}). The elements of
$\so{3n,3n}$ commuting with it and with $G_n$ are the copy rotations alone, of dimension $0,1,3$ for $n=1,2,3$; no copy boost commutes with $\lint$.
At $n=2$ the copy rotation is $C^{12}=-M_{12}+C_{\rm surp}$, where $M_{12}$ of \cref{sec:diagonal}, the spin lift of $E_{12}-E_{21}$, acts on vectors as $-\varepsilon\otimes P_{\rm int}$. It is the copy rotation of the interior
four-planes, the copy algebra of \cref{thm:diagcent}, and $C_{\rm surp}=\varepsilon\otimes P_{\rm surp}$ that of the surplus pairs. Both
commute with the diagonal $\lint$, and neither commutes with $G_2$.
\item \emph{Spectator action} ($\lG^{(1)}$ or $\lint^{(1)}$ on copy~$1$ alone). The elements of $\so{3n,3n}$ commuting with it and with $G_n$ form
$\so{3n-3,\C}$ of copies $2,\dots,n$, of dimension $6$ and $30$ at $n=2,3$. Every one of them vanishes on copy~$1$ and has no component into it.
\end{enumerate}
\end{proposition}

Part (i) restates the third row of \cref{prop:commutantsn} (fixed there) and \cref{d5:prop:so6C}(iii) at $n=2$. For $\lint^{(1)}$, part (iii) restates
\cref{thm:nocommon} at $n=3$ (fixed there) and \cref{d5:prop:so6C}(iv) at $n=2$. For $\lG^{(1)}$ it is fixed here (d7-1d), and its proof is Schur's lemma: $V$ is an irreducible,
nontrivial module of $\lG$, and the other copies are trivial modules of it. For $\lint^{(1)}$ Schur's lemma alone does not suffice. $V$ splits into the interior four-plane
$\langle e_1,e_4,e_5,e_6\rangle$, which is irreducible and nontrivial, and the surplus pair $\langle e_2,e_3\rangle$, on which $\lint$ acts trivially. The commutant of $\lint^{(1)}$ alone, of
dimension $28$ and $91$ at $n=2,3$ (the first row of \cref{prop:commutantsn}), therefore acts on the surplus pair of copy~$1$. Commutation with $G_n$ removes
these elements, because $G$ maps the surplus pair into the interior four-plane, as in the proof of \cref{thm:PSbinding}. Part (ii) and the decomposition of
$C^{12}$ are fixed (d7-1a to d7-1d). Two consequences matter below. A binding-compatible, Lorentz-invariant field can couple the manifold that
carries the Lorentz algebra to the others only when that algebra acts on all of them in the same way. Under the spectator action copy~$1$ is out
of reach, and the inter-manifold fields act among the internal manifolds only. And the copy algebra $\mathfrak k_{\rm cop}$ of \cref{ch:ch05_tower} does not respect
the binding: the binding-compatible copy rotation $C^{kl}$ is $-M_{kl}$ plus the surplus copy rotation.

% --- source: rest_energy/d7_between#3 (The closure admits an inter-manifold field)
\subsection{The closure admits an inter-manifold field}
\label{d7:sec:rotor}

\cref{thm:closureunique} shows that, on one constituent, the only equivariant closure is $\R\ltimes_GV$: the rotor acts by $G$. After it the
paper extends the result to several constituents ``factor by factor, since equivariant maps between distinct constituents vanish'', with a common
rate. That is correct for a Lorentz algebra acting on each constituent separately, once the common rate is imposed. A common action admits more.

\begin{proposition}[Rotor generators on $n$ binding frames]
\label{d7:prop:rotor}
The endomorphisms of $V_n$ commuting with the diagonal $\lG$ form $\mathfrak{gl}(n,\C)=\mathfrak{gl}(n)\otimes\operatorname{span}\{\id,G\}$, of dimension
$2n^2$. Those that rotate the frame covector without dilating it, the condition of \cref{thm:closureunique}, form
\[
\uu n=\so n\otimes\id_6\ \oplus\ \operatorname{Sym}(n)\otimes G ,
\]
of dimension $n^2$. The first summand is $\eta_n$-antisymmetric (copy rotations) and the second is $\eta_n$-self-adjoint (rate matrices). The paper's
rotor $G_n$ spans the centre. For $\mathfrak s_n=\R\ltimes_AV_n$ with $A=K\otimes G+B\otimes\id_6\in\uu n$ and the rotor clamped at $\theta=\Omega\tau$, the substitution $X=e^{-\theta A}x$ of \cref{prop:groupred} gives $e^{-\theta A}\dot x=\dot X+\Omega AX$, so the shift is $S=\Omega A$.
\end{proposition}

At $n=2$ the algebra is the energy symmetry $\uu2$ of \cref{d5:prop:energysym}(i). The algebra $\uu n$, its dimension $n^2$ and its centre $G_n$
restate row four of \cref{prop:commutantsn} and the remark after it (fixed there). The commutant $\mathfrak{gl}(n,\C)$ and the
decomposition into copy rotations and rate matrices are fixed (d7-2a); the reading as rotor generators is new.

The semidirect product $\R\ltimes_AV_n$ is a Lie algebra for every $A$, so the closure exists. It is not the only closure with rotor $A$, because
the half of \cref{thm:closureunique} that makes the frame abelian does not carry over to a common action. An equivariant bracket on $V_n$ is
$[X,Y]_k=\sum_{i,j}\Gamma^k_{ij}\,\beta_0(X_i,Y_j)$ with complex $\Gamma^k_{ij}=\Gamma^k_{ji}$ ($G$ acting as $i$), a space of real dimension $12$ at $n=2$. In a
copy basis in which $A$ acts as $i\operatorname{diag}(\omega_1,\dots,\omega_n)$, the Jacobi identity with the rotor reads $\Gamma^k_{ij}(\omega_k-\omega_i-\omega_j)=0$. The
frame is therefore abelian for $A=G_n$ and whenever no $\omega_k$ equals $\omega_i+\omega_j$. For a resonant $A$ it need not be: for
$A=\operatorname{diag}(1,2)\otimes G$ the bracket $[X,Y]=(0,c\,\beta_0(X_1,Y_1))$, $c\in\C$, satisfies every Jacobi identity (fixed; d7-2b). This chapter uses only the abelian-frame closure, to which the shift $S=\Omega A$ refers.

Which rotor acts is a choice. \cref{thm:closureunique} selects $G_n$ only when the action is taken per constituent and a common rate is
imposed. Per constituent alone it admits $K\otimes G$ with any real diagonal $K$: unequal rates, but no copy field, since equivariant maps between
distinct constituents vanish. Under the common action a common rate alone still admits $\kappa G_n+B\otimes\id_6$: a copy field, but one that commutes
with the binding and is removable (\cref{d7:thm:schr}(v)). The two parts of a non-central $A$, the $\eta_n$-antisymmetric copy field
$B\otimes\id_6$ and the $\eta_n$-self-adjoint rate matrix $K\otimes G$, are independent; their ratios are fixed by $A$ and their scale by the rotor.
The common action is the one place found where the closure admits unequal rates together with a copy field, the combination that
\cref{d5:thm:bfield2} needs. That such an $A$ is admissible is constructed: nothing in the algebra selects it. The paper's one-rotor closure, with
$A=G_n$, gives every constituent the same rate (\cref{prop:several}).

% --- source: rest_energy/d7_between#4 (The frame flow is a Schr\"odinger flow)
\subsection{The frame flow is a Schr\"odinger flow}
\label{d7:sec:schr}

\begin{theorem}[Binding-compatible fields between $n$ manifolds]
\label{d7:thm:schr}
Let $b\in\Cent(G_n)\cap\so{3n,3n}$. By \cref{d5:prop:so6C}(i), $b$ acts on $z$ as $R+iQ$ with $R,Q\in\so{3n}$ real: $R$ is
its paired-rotation part and $Q$ its paired-boost part. Put $\tilde K=K\otimes\id_3$ and $H:=\tilde K-iR+Q$.
\begin{enumerate}[label=(\roman*),leftmargin=2.2em]
\item $S$ acts on $z$ as $iH$, so the free frame flow is $z(\tau)=e^{-iH\tau}z(0)$.
\item The flow preserves $E_6=\tfrac12|z|^2$ for every $K$ and every initial value iff $Q=0$, that is, iff $b$ is compact; then $H$ is Hermitian. The
compact fields, the only ones admissible as Yang--Mills fields (\cref{cor:noboost}), are exactly those whose frame flow is unitary for the
undivided energy.
\item The rates are gauge invariants. Under any $\eta_n$-isometric change of frame $X=U(\tau)Y$, the $\eta_n$-self-adjoint part of $S$ is conjugated,
so its spectrum $\{\pm i\kappa_k\}$ is unchanged. No field, compact or not, changes a rate.
\item For compact $b$: $\tr H=3\sum_k\kappa_k$ and $\tr H^2=3\sum_k\kappa_k^2+\|R\|^2$. The field redistributes the frequencies of the frame at fixed
sum and widens their spread: their variance grows by $\|R\|^2/(3n)$, the largest frequency does not decrease and the smallest does not increase.
\item $b$ is removable along the worldline iff $[\tilde K,R+iQ]=0$, that is, iff it preserves every rate eigenspace. At $n=3$ the compact
binding-compatible fields, of dimension $36$, have a removable part of dimension $36$, $18$ and $9$ for rates $(\kappa,\kappa,\kappa)$, $(\kappa,\kappa,\kappa')$
and three distinct rates. Within the copy $\so3$ the dimensions are $3$, $1$ and $0$.
\item (Boundedness) Let $b$ be compact. Every orbit of $(X,\pi)$ is bounded iff, for every $\lambda$ with $\lambda,-\lambda\in\spec H$ (including $\lambda=0$),
the eigenspaces $\mathcal E_\lambda(H)$ and $\mathcal E_{-\lambda}(H)$ are orthogonal for the bilinear form $v\cdot w=\sum_av_aw_a$ on $\C^{3n}$.
\end{enumerate}
\end{theorem}

\begin{proof}
(i) On the sector $G_n=i$, $K\otimes G$ acts as $i\tilde K$, a paired rotation as $R$ and a paired boost as $iQ$ (\cref{bf:lp,d5:prop:so6C}); $\dot X=-SX$ becomes $\dot z=-iHz$. (ii) $E_6=\tfrac12|X|^2$, so it is conserved for all $X$ iff $S$ is Euclidean
antisymmetric. $K\otimes G$ and the paired rotations are, and the paired boosts are Euclidean symmetric. (iii) is
\cref{bf:covariant}(ii) on $V_n$. The self-adjoint part transforms as a tensor, $U^{-1}(K\otimes G)U$, while $U^{-1}bU+U^{-1}\dot U$ stays in
$\so{3n,3n}$. (iv) $\tr(\tilde KR)=0$ because $\tilde K$ is symmetric and $R$ antisymmetric, and $\tr(-iR)^2=\|R\|^2$. The variance follows from the two sum rules. For a real unit eigenvector $v$ of $\tilde K$, $v^{\top}Rv=0$, so $\bar v^{\top}Hv$ is its rate, and the largest and smallest eigenvalues of $H$ bound it from above and below. (v) is
\cref{bf:larmor}: $e^{-b\tau}$ intertwines the systems iff $b$ commutes with $K\otimes G$. (vi) is \cref{bf:criterion} with
$S^{\!\top}=-S$ (both summands are Euclidean antisymmetric). Then $\mathcal E_\mu(S^\dagger)=\eta_n\mathcal E_{-\mu}(S)$, and the condition reads
$\eta_n(\mathcal E_\nu(S),\mathcal E_{-\nu}(S))=0$ for every $\nu$. The sectors $G_n=\pm i$ are $\eta_n$-orthogonal, and on the sector $G_n=i$,
$\eta_n=2\,v\cdot w$ in the coordinates $z$. The eigenspaces of $S$ there are those of $iH$, and on the other sector their conjugates.
\end{proof}

Parts (i), (ii), (iv) and (v) are fixed but elementary: the arena's version of the fact that a compact gauge field acts unitarily
(d7-3a, d7-3b, d7-3d, d7-3e). The widening in (iv) concerns the variance and the range, not each gap. At rates $(1,2,4)$ with $(B_{12},B_{23},B_{31})=(\tfrac12,\tfrac13,\tfrac15)$ the upper gap falls from $2$ to $1.916$, so the field does not repel every pair of levels (x7-3d). Part (iii) is \cref{bf:covariant}(ii) restated for $n$ copies, checked exactly on a Cayley curve (d7-3c). Part (vi) is
a corollary of \cref{bf:criterion}. On five cases at $n=3$, exact semisimplicity of $M$, the criterion and the growth of $\|e^{Mt}\|$ agree
(d7-3f). Bounded: rates $(1,2,4)$ with $(B_{12},B_{23},B_{31})=(\tfrac12,\tfrac13,\tfrac15)$; the seesaw $(0,3,5)$ with $(1,\tfrac12,\tfrac13)$, a
frame without rate confined by two couplings; equal rates $(1,1,1)$ with $B_{12}=1$. Unbounded: $(1,2,3)$ with $(1,1,1)$, where $\spec H=\{0,2,4\}$ and the
zero mode has $v\cdot v\neq0$, a resonance created by the field between manifolds with different rates; and zero rates with $(1,1,1)$, where
the field produces only drift. Part (vi) generalises \cref{d5:thm:bfield2}(v): at $n=2$, for the field of \cref{d5:thm:bfield2}, $\spec H$ is the pair of lines $\omega_\pm$ of \cref{d5:thm:bfield2}(ii), with $\omega_++\omega_-=\kappa_1+\kappa_2$, so for $\kappa_1+\kappa_2>0$ the only coincidence
is $\omega_-=0$, at $\beta^2=\kappa_1\kappa_2$. At $\kappa_1=\kappa_2=0$ and $\beta\neq0$ it is $\omega_+=-\omega_-=|\beta|$, with eigenspaces that are not
orthogonal for the bilinear form, so at zero rates the orbits are unbounded for every $\beta$, as in \cref{d5:thm:bfield2}(v).

% --- source: rest_energy/d7_between#5 (Copy boosts: metric connections, not gauge fields)
\subsection{Copy boosts: metric connections, not gauge fields}
\label{d7:sec:boost}

\begin{proposition}[The copy boost between two manifolds]
\label{d7:prop:boost}
Let $n=2$ and $b=\beta'\varepsilon\otimes G$, so that $H=K+\beta'\varepsilon$ on the copy index (\cref{d7:thm:schr}(i)), real and not symmetric. Its eigenvalues are
$\tfrac12(\kappa_1+\kappa_2)\pm\sqrt{\tfrac14(\kappa_1-\kappa_2)^2-\beta'^2}$.
\begin{enumerate}[label=(\roman*),leftmargin=2.2em]
\item Below the exceptional point, $|\beta'|<\tfrac12|\kappa_1-\kappa_2|$, and provided $\kappa_1\kappa_2+\beta'^2>0$, the spectrum is real, $M$ is semisimple with
imaginary spectrum, and every orbit is bounded. The flow then preserves a positive form $z^\dagger\mathcal Qz$ with $H^\dagger\mathcal Q=\mathcal QH$, which for
$\beta'\neq0$ is not $E_6$. (At $(\kappa_1,\kappa_2,\beta')=(0,1,0)$ the first copy is free and $M$ is not semisimple.)
\item At $|\beta'|=\tfrac12|\kappa_1-\kappa_2|$ with $\beta'\neq0$, $M$ is not semisimple (at $\beta'=0$ and $\kappa_1=\kappa_2>0$ the copies are uncoupled
oscillators and $M$ is semisimple; at $\kappa_1=\kappa_2=\beta'=0$ the motion is free and $M$ has nilpotent Jordan blocks); above it, $M$ has eigenvalues with real part $\pm\sqrt{\beta'^2-\tfrac14(\kappa_1-\kappa_2)^2}$.
\item At equal rates the threshold is zero: every $\beta'\neq0$ grows exponentially, although $b$ commutes with $K\otimes G$ and is removable, by a
boosted frame $e^{-b\tau}$ that is itself unbounded.
\end{enumerate}
\end{proposition}

This is fixed, checked exactly at rates $(1,4)$ with $\beta'=\tfrac12,1,\tfrac32,2$, and at $(2,2)$ with $\beta'=\tfrac13$ (d7-4a, d7-4b); the boundary cases at $\beta'=0$, rates $(1,1)$ and $(0,0)$, are d7-4c. Its threshold is the
copy analogue of the homogeneous threshold of the pairing boosts of \cref{prop:tear}, with the rate difference
$\tfrac12|\kappa_1-\kappa_2|$ in the place of $\kappa$. The analogy stops there. A pairing boost anticommutes with $G$ and gives Jordan blocks for
every $\nu\neq0$, below the threshold as well (\cref{bf:pG}(iii)), whereas the copy boost commutes with $G_2$ and keeps every orbit bounded
below its threshold. At equal rates it is the copy analogue of the commuting boosts $K_i$ of \cref{bf:lG}(iv), with exponents
$\pm\beta'\pm i\kappa$, removed by a boosted frame that is itself unbounded. The spectral pattern is inherited: the
non-Hermitian dimer with an exceptional point \cite{Kato66,BenderBoettcher98}, and the conserved metric of a pseudo-Hermitian Hamiltonian
\cite{Mostafazadeh02}. A copy boost is a metric connection, legitimate in the interior Lagrangian, but not a gauge field (\cref{cor:noboost}). Where it keeps the motion bounded it does so at the price of the undivided energy, which it replaces by another positive form.

% --- source: rest_energy/d7_between#6 (Three manifolds: a time-odd invariant)
\subsection{Three manifolds: a time-odd invariant}
\label{d7:sec:three}

\begin{theorem}[The copy field between three manifolds]
\label{d7:thm:three}
Let $n=3$, distinct rates, and $b=B\otimes\id_6$, so $H=(K-iB)\otimes\id_3$.
\begin{enumerate}[label=(\roman*),leftmargin=2.2em]
\item $\det(x-K+iB)=x^3-(\textstyle\sum\kappa)x^2+(\sum_{k<l}\kappa_k\kappa_l-\sum B^2)x-(\kappa_1\kappa_2\kappa_3-\kappa_3B_{12}^2-\kappa_1B_{23}^2-\kappa_2B_{31}^2)$: the spectrum
depends on the couplings only through their squares.
\item The plaquette $H_{12}H_{23}H_{31}=iJ$ is purely imaginary, and $\det[H,K]=2iJ(\kappa_1-\kappa_2)(\kappa_1-\kappa_3)(\kappa_2-\kappa_3)$.
\item The endomorphisms commuting with $G_3$, the diagonal $\lG$ and $K\otimes G$ are $\sum_k(a_k+b_kG)$ on copy $k$, and the $\eta_3$-isometries among them are
exactly the eight copy reflections $\operatorname{diag}(\pm1,\pm1,\pm1)\otimes\id_6$. Under them $B_{kl}\mapsto s_ks_lB_{kl}$, so the invariants of the field
are $|B_{12}|,|B_{23}|,|B_{31}|$ and the sign of $\plaq $.
\item With $P=\id_3\otimes\operatorname{diag}(\id_3,-\id_3)$ ($\bE\mapsto-\bE$) and $T=-P$ ($\bB\mapsto-\bB$), parity maps $(K,B)$ to $(-K,B)$, and interior
time reversal, $X(\tau)\mapsto TX(-\tau)$, maps $(K,B)$ to $(K,-B)$. At $n=2$ a copy reflection undoes $B\mapsto-B$; at $n=3$ with $\plaq \neq0$ none does, since
$\plaq \mapsto-\plaq $. For general $n$, $(K,B)$ and $(K,-B)$ are related by a copy reflection iff the graph of the nonzero couplings is bipartite.
\item With energy started in copy~$1$, the copy energies $P_k(\tau)=|z_k(\tau)|^2/|z(0)|^2$ obey $P_k(\tau;-B)=P_k(-\tau;B)$. $P_1$ is the same for $B$ and
$-B$, and $P_2,P_3$ are not. At rates $(1,2,4)$ with $B_{12}=B_{23}=B_{31}=1$ and $\tau=\tfrac12$, $P_2=0.293$ for $B$ against $0.116$ for $-B$: the sign
of $\plaq $ is the sense in which the energy circulates around the three manifolds.
\end{enumerate}
\end{theorem}

\begin{proof}
(i), (ii) Symbolic expansion (d7-5a, d7-5b). For a $3\times3$ Hermitian matrix, $\det H$ contains $2\Re(H_{12}H_{23}H_{31})$, which vanishes here.
(iii) The commutant is exact (d7-5c). $(a+bG)^{\!\top}\eta(a+bG)=(a^2-b^2)\eta+2ab\,\eta G$ because $G^{\!\top}\eta=\eta G$ and $G^2=-\id$, so $ab=0$
and $a^2-b^2=1$. (iv) $PGP^{-1}=-G$, $P$ commutes with $B\otimes\id_6$, and $\dot X=-SX$ goes to $\dot X'=-S'X'$ with $S'=-TST^{-1}$. A copy
reflection $s$ maps $B$ to $-B$ iff $s_ks_l=-1$ on every edge, which is a two-colouring (d7-5d, d7-5e). (v)
$|(e^{-i\bar H\tau})_{k1}|=|(e^{iH\tau})_{k1}|$, and $P_1$ is even in $\tau$ for Hermitian $H$ (d7-5f).
\end{proof}

The statements are fixed. Their pattern is inherited three times over. With two families there is no rephasing-invariant phase and with three
there is one, the Kobayashi--Maskawa phase \cite{KobayashiMaskawa73}, whose invariant measure is Jarlskog's commutator determinant
\cite{Jarlskog85}, the form of (ii). A purely imaginary hopping around a triangle is a flux $\pm\pi/2$ through the plaquette, the Peierls phase
\cite{Peierls33} of a lattice gauge field. And time-reversal-odd hopping makes a quantum walk chiral \cite{Zimboras13}. What the arena adds is
the identification. The copy field is an $\eta$-compatible connection on the three-site ``lattice'' of manifolds in copy space, with $\plaq $ its
holonomy around the triangle, the copy-space analogue of a Wilson loop; it is not the Wilson loop of a spacetime connection, and its sign is not
by itself a physical CP phase (\cref{e4:thm:cycles}). Its phase is fixed at $\pm\pi/2$ because the field is real, $H_{kl}=-iB_{kl}$. Like a Wilson loop it is unchanged by per-copy duality rotations, which multiply each $H_{kl}$ by a phase and are in any case not $\eta$-isometries
(Proposition~\ref{cl:norel}(ii)). On one manifold the binding is odd under parity and even under interior time reversal. At two manifolds the
copy field is even under both, up to relabelling. From three on, a field whose couplings close a triangle is odd under interior time reversal,
and no relabelling removes it.

% --- source: rest_energy/d7_between#7 (Isospin one, not one-half)
\subsection{Isospin one, not one-half}
\label{d7:sec:isospin}

Yang and Mills' $b$-field acts on the isotopic doublet of nucleons \cite{YangMills54}. The question is whether an inter-manifold field can.

\begin{proposition}[Copy isospin]
\label{d7:prop:isospin}
\begin{enumerate}[label=(\roman*),leftmargin=2.2em]
\item On the tower module, the spin lift of the copy rotation of a set $D$ of directions is $-\tfrac12\sum_{a\in D}\eta_{aa}e^{(k)}_ae^{(l)}_a$, a sum of $|D|$
commuting terms of square $-\tfrac14$, so $\exp(2\pi\,\cdot)=(-1)^{|D|}$. At $n=2$: for the six binding directions the eigenvalues are $i\{-3,\dots,3\}$ and
$\exp(2\pi C)=+\id$. For the four interior directions ($M_{12}$) they are $i\{-2,\dots,2\}$, $+\id$. For the three temporal directions they are
$i\{\pm\tfrac12,\pm\tfrac32\}$, $-\id$, and that rotation does not commute with $G_2$. A binding-compatible copy rotation pairs each temporal direction with
its partner, so $|D|$ is even.
\item At $n=3$ the three binding copy rotations close on $\so3$ on the $512$-dimensional module, with Casimir $j(j+1)=0,2,6,12$ on $40,216,200,56$
dimensions ($j=0,\dots,3$), and $\exp(2\pi C^{12})=+\id$. The frames form the vector, $j=1$.
\item The doublet $\uu2$ acting on $(z_1,z_2)$ preserves $E_6\oplus E_6$ and commutes with $G_2$ and the diagonal $\lG$. In real form $i\id\mapsto\id\otimes G$,
$i\sigma_x\mapsto\sigma_x\otimes G$, $i\sigma_y\mapsto-\varepsilon\otimes\id_6$ and $i\sigma_z\mapsto\sigma_z\otimes G$. Only $i\sigma_y$, the copy rotation, is
$\eta_2$-antisymmetric. The other three are $\eta_2$-self-adjoint, terms of the same type as the binding. For $S=\kappa G_2+b_{\rm YM}$,
$b_{\rm YM}=\sum_a\beta_a\,(i\sigma_a)$, the lines are $\kappa\pm|\boldsymbol\beta|$, and the rates, the spectrum of the self-adjoint part, are
$\kappa\pm(\beta_x^2+\beta_z^2)^{1/2}$.
\end{enumerate}
\end{proposition}

This is fixed (d7-6a to d7-6c). Part (i) restates \cref{prop:orthospin} (``the copy algebras act through orthogonal groups'') and
\cref{d5:prop:spin}(i), and adds that binding-compatibility forces $|D|$ to be even, so every binding-compatible copy rotation lifts single-valuedly. The criterion is the parity of $|D|$, not the binding: $M_{12}$ is not binding-compatible and also lifts single-valuedly. Under a common Lorentz action the inter-manifold field
is therefore an $SO(n)$ gauge field. At $n=3$ it acts with integer isospin only: the frames are triplets, as in the model of Georgi and Glashow \cite{GeorgiGlashow72}, the tower module carries $j=0,\dots,3$, and it never acts on
doublets. Part (iii) restates \cref{d5:prop:energysym}(i) in the Yang--Mills reading and adds the spectra; it is the precise form of
``Yang--Mills between two copies''. The doublet group exists as a symmetry of the undivided energy
(the copy $\uu2$), but $\eta_2$ cuts it to its circle. The two directions it removes are not connections at all:
they belong to the symmetric part, where the binding $c^2$ lives in the dictionary of \cref{s0:sec:question}. A doublet field that splits the
rates between two manifolds does so only through those two directions. In the arena that is a modification of the symmetric part, not a
$b$-field.

\begin{proposition}[The isospin $\so4_\Delta$ between unequal constituents]
\label{d7:prop:so4unequal}
At $n=3$, with the numbering of \cref{sec:closure} and $T^\pm_i$ as in \cref{bf:su2}, let $K=\operatorname{diag}(\kappa_1,\kappa_2,\kappa_3)$ with
$\kappa_2\neq\kappa_3$.
\begin{enumerate}[label=(\roman*),leftmargin=2.2em]
\item The part of $\so4_\Delta$ commuting with $K\otimes G$ is $\operatorname{span}\{\mathrm{rot}(13,14),\mathrm{rot}(13,15),\mathrm{rot}(14,15)\}
=\operatorname{span}\{T^+_1+T^-_1,\ T^+_2-T^-_2,\ T^+_3-T^-_3\}$, the graph of an isomorphism $\su2_+\to\su2_-$: a diagonal $\so3$. The three directions
$\mathrm{rot}(7,13),\mathrm{rot}(7,14),\mathrm{rot}(7,15)$, which join constituent two to constituent three, are not removable.
\item For $\kappa=(5,3,5)$ and $b=2T^+_3$,
\[
\det(\lambda-S)=(\lambda^2+9)^2(\lambda^2+16)(\lambda^2+25)^3(\lambda^2+36)(\lambda^4+36\lambda^2+196),
\]
and $M$ is semisimple. The quartic carries the lines $4\pm\sqrt2$, the two-level mixing of $z_7$ and $z_{13}$ with rates $3,5$ and coupling $1$
(\cref{d5:thm:bfield2}(ii)). At equal rates $5$ the same field gives $(\lambda^2+25)^5(\lambda^2+16)^2(\lambda^2+36)^2$, the removable Zeeman pattern
of \cref{bf:so4spec}(i).
\end{enumerate}
\end{proposition}

This is fixed (d7-7a, d7-7b). It completes \cref{bf:so4spec}, which used equal rates. The paper's binding-compatible non-abelian algebra is
itself an inter-manifold field: three of its six directions, $\mathrm{rot}(7,13)$, $\mathrm{rot}(7,14)$ and $\mathrm{rot}(7,15)$, join constituent two to constituent three, and the other three rotate constituent three alone. It carries doublets ($(\tfrac12,\tfrac12)$ on
$W$) because the spectator action leaves constituents two and three free of the Lorentz algebra (\cref{d7:prop:exist}(iii)). With equal rates it is
removable along a worldline, as \cref{sec:bf-three} found. With unequal rates it breaks to a diagonal $\so3$, and its three remaining directions
become physical. The pattern $\su2\oplus\su2\to\su2_{\rm diag}$ is that of chiral symmetry breaking and of the custodial symmetry of the Higgs
sector \cite{SikivieEtAl80}. That is an interpretive parallel only: here the breaking is by a constant rate matrix, not by a vacuum.

% --- source: rest_energy/d7_between#8 (Assessment)
\subsection{Assessment}
\label{d7:sec:assess}

\begin{remark}[The deal breaker, settled]
\label{d7:rem:verdict}
\begin{enumerate}[label=(\alph*),leftmargin=2.2em]
\item \emph{True, in this sense} (fixed). A compact binding-compatible field that couples manifolds of different rates cannot be removed: by \cref{d7:thm:schr}(v) it is removable exactly when it preserves every rate eigenspace. It shifts and
splits the lines, it can confine a frame that has no rate of its own, it can create a resonance, and from three manifolds on it is non-abelian and
carries a time-odd invariant. What does not extend from one manifold is the removability of a field that commutes with the binding root
(\cref{bf:larmor}). What does extend is the no-go of \cref{bf:alone}: a connection alone does not bind, and at zero rates it produces
only drift (\cref{d5:thm:bfield2,d7:thm:schr,d7:thm:three}).
\item \emph{False, in this sense} (fixed). The field never creates or changes a rate. Rates are the spectrum of the self-adjoint part, a tensor under
every change of frame. The field only redistributes the frequencies, the lines of $\spec H$, at fixed sum. With equal rates every binding-compatible field, abelian or the $\so4_\Delta$ of \cref{thm:PSbinding}, is removable along a worldline. With zero rates it produces only drift (\cref{d7:thm:schr}(iii)--(vi)).
\item \emph{Where unequal rates come from.} Not from the closure of \cref{sec:closure}: it imposes a common rate, so its single rotor acts through $G_n$ at one rate, and within that closure as it stands
the inter-manifold field is removable (fixed for that closure). Under a common binding action the closure admits a rotor acting by a non-central element
of $\uu n$. The copy field and the rate differences can then both enter, as two independent parts of one generator, and either can occur
without the other (constructed: admitted, not selected; \cref{d7:prop:rotor}).
\item \emph{Isospin one-half} (fixed). On two and three manifolds under a common Lorentz action the copy isospin is integral, and the doublet field on two copies is a connection
in one of its three directions only. On three manifolds doublets need the spectator action, where the Lorentz-carrying manifold is decoupled
(\cref{d7:prop:exist,d7:prop:isospin,d7:prop:so4unequal}). This does not extend to four manifolds. There the binding-compatible copy rotations under a common
action form $\so4\cong\su2\oplus\su2$, the copy index is its vector $(\tfrac12,\tfrac12)$, and doublets exist without a spectator (d7-6d).
\item \emph{The arrow reversed} (interpretive). The non-removable part of the copy field is $[K\otimes G,B\otimes\id_6]=[K,B]\otimes G$, with
$[K,B]_{kl}=(\kappa_k-\kappa_l)B_{kl}$ and $\|[K\otimes G,b]\|^2=6\sum_{k\neq l}(\kappa_k-\kappa_l)^2B_{kl}^2$ (d7-8, fixed). Formally this is the gauge-boson
mass term $g^2\|[T,\langle\Phi\rangle]\|^2$ of the Englert--Brout--Higgs--Kibble mechanism \cite{EnglertBrout64,Higgs64,Kibble67}, with the rate matrix
in the role of an adjoint vacuum value. Read with $mc^2=\hbar\kappa$, so that a mass is a rate, the arena reverses the arrow of the hypothesis. It is not the $b$-field that differentiates the rates;
the differentiation of the rates, which breaks the copy symmetry, is what gives the $b$-field an effect. The parallel is not dynamical: here $b$ is a
constant connection without kinetic term, and the rates are inputs. Read instead from the lines, the spectrum of $H$, a mass is shifted and split by the field, and a frame without rate of its own is confined at a nonzero line, the seesaw (\cref{d5:thm:bfield2}(vi)). That effect is non-removable only when the rates differ (\cref{d7:thm:schr}(v)), so in this reading too the field acts through the rate differences and does not create them. In the closure of \cref{d7:prop:rotor} the copy field and the rate differences are independent parts of one generator, so the order is logical, not causal.
\end{enumerate}
\end{remark}

\begin{table}[t]
\centering
\footnotesize
\begin{tabular}{@{}>{\raggedright\arraybackslash}p{8.1cm}>{\raggedright\arraybackslash}p{5.0cm}>{\raggedright\arraybackslash}p{2.0cm}@{}}
\toprule
statement & class & checks\\
\midrule
\cref{d7:prop:exist}(i), (iii): $\so{n,\C}$ under a common binding action; copy~$1$ decoupled under the spectator action & fixed (restated: Prop.~13.14; Thm~17.20 for $\lint^{(1)}$) & d7-1a, d7-1d\\
\cref{d7:prop:exist}(ii): $\so n$ under the interior action; $C^{12}=-M_{12}+C_{\rm surp}$ & fixed & d7-1b, d7-1c\\
\cref{d7:prop:rotor}: rotor generators $\uu n=\so n\oplus\operatorname{Sym}(n)\otimes G$; abelian frame only for non-resonant $A$ & fixed (restated: Prop.~13.14); non-central rotor constructed & d7-2a, d7-2b\\
\cref{d7:thm:schr}(i), (ii), (iv), (v): Schr\"odinger form, unitarity iff compact, sum rules, removability & fixed (elementary) & d7-3a, d7-3b, d7-3d, d7-3e\\
\cref{d7:thm:schr}(iii): rates are gauge invariants & fixed (restated: \cref{bf:covariant}) & d7-3c\\
\cref{d7:thm:schr}(vi): bilinear boundedness criterion; five cases at $n=3$ & fixed (corollary of \cref{bf:criterion}) & d7-3f\\
\cref{d7:prop:boost}: copy boost, exceptional point, conserved positive form & fixed; pattern inherited & d7-4a, d7-4b\\
\cref{d7:thm:three}: spectrum, plaquette, Jarlskog determinant, reflections, $P$ and $T$, bipartite criterion, circulation & fixed; parallels inherited & d7-5a to d7-5f\\
\cref{d7:prop:isospin}: integral copy isospin at $n=2,3$; the doublet is one connection and two rate terms; doublets at $n=4$ & fixed (restated: Prop.~10.9, \cref{d5:prop:energysym}) & d7-6a to d7-6d\\
\cref{d7:prop:so4unequal}: $\so4_\Delta\to$ diagonal $\so3$ at unequal rates & fixed; custodial parallel interpretive & d7-7a, d7-7b\\
\cref{d7:rem:verdict}: the deal breaker; the commutator; the arrow reversed & fixed; (c) constructed; (e) interpretive & d7-8\\
\bottomrule
\end{tabular}
\caption{Claims of \cref{d7:sec:main}, their classes, and the checks printed by \texttt{d7\_bfield\_between.py}.}
\label{d7:tab:claims}
\end{table}

% --- source: emergence/e4_bfield_ladder#0 ((head))
% Track e4 -- The b-field between the constituents of a rung, up the ladder.
% New references: sections/bib_e4.tex.

% --- source: emergence/e4_bfield_ladder#1 (The b-field up the ladder)
\section{\texorpdfstring{The $b$-field up the ladder}{The b-field up the ladder}}
\label{e4:sec:main}

\Cref{d7:sec:main} studied the field that acts between the constituents of one rung, the copy field $b=B\otimes\id_6$, and settled what it
does to rates and lines on two and three constituents. This section follows that field up the ladder, $n=2,3,4,\dots$, and reads it in
the Fock picture of \cref{e1:sec:fock,e2:sec:quanta}. The field always acts between the constituents of one rung $\Cl(3n,3n)$; going up the
ladder means increasing $n$.

The answers, with their classes, are these. (1) Every compact binding-compatible field is the second quantisation of a one-quantum
Hamiltonian. It conserves the number $N$ of quanta, moves quanta between the constituents, and its many-quantum lines are sums of
one-quantum lines. The copy boosts create pairs, one quantum in each of the two constituents they join (fixed; \cref{e4:sec:fock}). (2) On $n$
constituents with distinct rates and all couplings nonzero, the invariants of the copy field are the $|B_{kl}|$ and the signs of the
products of $B$ around the cycles of the complete graph. Exactly $(n-1)(n-2)/2$ of these signs are independent, and all of them can be taken
time-odd. The Jarlskog-type determinant $\det[H,K]$ has time parity $(-1)^n$, so it detects time reversal only at odd $n$ (fixed;
\cref{e4:sec:cycles}). (3) At $n=4$ the copy group is $SO(4)$, and each of its two $\su2$ factors has half-integer isospin exactly on the states
with an odd number of quanta. The copy inversion acts as $(-1)^N$. Above $n=4$ the same holds for the $\so4$ of any four constituents, with the
quanta in those four. At no rung does a spinor of the copy algebra occur (fixed; \cref{e4:sec:isospin}). (4) For rates of multiplicities $n_1+\dots+n_r=n$ the removable part of the copy $\so n$ is $\bigoplus_i\so{n_i}$, and
$\sum_{i<j}n_in_j$ directions are physical (fixed; the Higgs-like reading is interpretive; \cref{e4:sec:rates}). (5) For a rotor with rates
$(\kappa_k)$, a frame bracket between copies $i,j$ and $k$ is admitted only if $\kappa_k=\kappa_i+\kappa_j$. A harmonic ladder of rates
$\kappa_k=k\kappa_1$ therefore admits non-abelian, nilpotent frames. From $n=4$ on the Jacobi identity among the frame brackets is a
quadratic condition (fixed; the reading is constructed; \cref{e4:sec:rotor}). \Cref{e4:sec:answer} says what the field can and cannot do for the emergence of particles.
Identifiers e4-$n$ refer to the \texttt{PASS} lines of \texttt{e4\_bfield\_ladder.py} (\cref{e4:tab:claims}).

\paragraph{Notation.} As in \cref{d7:sec:main}: $V_n=\R^n\otimes\R^6$, $\eta_n$, $G_n=\id_n\otimes G$, $z=(z_1,\dots,z_n)\in\C^{3n}$ with
$z_k=\bB_k+i\bE_k$, copy rotations $c^{kl}=(E_{lk}-E_{kl})\otimes\id_6$, copy boosts $(E_{lk}-E_{kl})\otimes G$, copy field $b=B\otimes\id_6$ with
$B\in\so n$, rates $K=\operatorname{diag}(\kappa_1,\dots,\kappa_n)$, $\tilde K=K\otimes\id_3$, and, for a compact binding-compatible $b$ acting on $z$ as the
real antisymmetric $R$, the frame Hamiltonian $H=\tilde K-iR$ of \cref{d7:thm:schr}. $\sigma$ is the spin lift to
$S_n=(\R^8)^{\otimes n}$, with $[\sigma(X),v]=Xv$ for $v\in V_n$, and $\hat C^{kl}=\sigma(c^{kl})=-C^{kl}$, where $C^{kl}=\half\sum_a\eta_{aa}E_{(k,a)}E_{(l,a)}$ are the spin-level copy rotations of \cref{e3:sec:main} and $C^{12}=C_b$. The $3n$ modes are $p=(k,i)$, $k\le n$, $i\le3$, with
$a^\pm_{k,i}=\half(E_{(k,i)}\pm E_{(k,i+3)})$ (\cref{d5:prop:fock}). $N_{k,i}=a^+_{k,i}a^-_{k,i}$, $N^{(k)}=\sum_iN_{k,i}$ is the number in
constituent $k$, $N_{[i]}=\sum_kN_{k,i}$ the number in direction $i$, and $N=\sum_kN^{(k)}$. For a complex $3n\times3n$ matrix $M$ we write
$\mathcal N(M)=\sum_{p,q}M_{pq}\,a^+_pa^-_q$. The first basis vector of the realisation is the Fock vacuum $|0\rangle$ (e4-1a).

% --- source: emergence/e4_bfield_ladder#2 (Quanta between constituents: transfer, not creation)
\subsection{Quanta between constituents: transfer, not creation}
\label{e4:sec:fock}

\Cref{e1:thm:Nrule} shows, at every rung, that the compact binding-compatible fields are exactly the binding-compatible fields that
conserve $N$, and \cref{e2:prop:hop,e2:prop:copyboost} give the hop and the pair at $n=2$. The theorem below adds the second-quantised form,
which says how the quanta move, and the bookkeeping by constituent and by direction at $n=3$.

\begin{theorem}[The field between constituents, second quantised]
\label{e4:thm:fock}
\begin{enumerate}[label=(\roman*),leftmargin=2.2em]
\item Let $X$ be a paired rotation, acting on $z$ as the real antisymmetric $R\in\so{3n}$. Then $\sigma(X)=\mathcal N(R)$. Hence $\sigma(X)$ commutes
with $N$, and $\sigma(X)\,a^+_q|0\rangle=\sum_pR_{pq}\,a^+_p|0\rangle$.
\item For the copy field $b=B\otimes\id_6$ and rates $K$, the operator $\mathsf H=\sum_k\kappa_kN^{(k)}-i\sigma(b)$ on $S_n\otimes\C$ equals
$\mathcal N(H)$, $H=(K-iB)\otimes\id_3$. It is Hermitian and commutes with $N$. Its block on $\{N=1\}$, in the basis $a^+_p|0\rangle$, is $H$, and its
eigenvalues on $\{N=m\}$ are the sums of $m$ eigenvalues of $H$ with distinct indices.
\item At $n=3$, $\hat C^{kl}=\sum_i\bigl(a^+_{l,i}a^-_{k,i}-a^+_{k,i}a^-_{l,i}\bigr)$. It commutes with $N$, with $N^{(k)}+N^{(l)}$, with $N^{(m)}$ for
$m\notin\{k,l\}$ and with every $N_{[i]}$, and not with $N^{(k)}$. The three lifts close on $\so3$ as on vectors, $[\hat C^{12},\hat C^{13}]=\hat C^{23}$.
\item At $n=3$, $\sigma\bigl((E_{lk}-E_{kl})\otimes G\bigr)=-\sum_i\bigl(a^+_{k,i}a^+_{l,i}+a^-_{l,i}a^-_{k,i}\bigr)$. It joins only sectors whose $N$ differs
by $2$, in both directions, and it commutes with $N^{(k)}-N^{(l)}$ and with $N^{(m)}$, $m\notin\{k,l\}$: every pair it creates has one quantum in
constituent~$k$ and one in constituent~$l$.
\item At $n=3$ the quarter turn $U=e^{\pi\hat C^{12}/2}=\tfrac18\prod_a(\id+2t_a)$, with $t_a=-\half\eta_{aa}E_{(1,a)}E_{(2,a)}$ the six commuting terms of
$\hat C^{12}$, is orthogonal, commutes with $N$ and $N^{(3)}$, and satisfies $UN^{(1)}U^{-1}=N^{(2)}$. For $\kappa_1\neq\kappa_2$ it changes the rate
operator $\sum_k\kappa_kN^{(k)}$ and keeps its spectrum.
\end{enumerate}
\end{theorem}

\begin{proof}
(i) A paired rotation acts in the same way on the temporal directions and on their partners, so $X(2a^\pm_q)=\sum_pR_{pq}\,2a^\pm_p$, and
$[\sigma(X),a^\pm_q]=\sum_pR_{pq}a^\pm_p$. The CAR give $[\mathcal N(R),a^+_q]=\sum_pR_{pq}a^+_p$ and $[\mathcal N(R),a^-_q]=-\sum_pR_{qp}a^-_p$,
which is the same because $R^{\!\top}=-R$. So $\sigma(X)-\mathcal N(R)$ commutes with every generator, and $S_n$ is irreducible, so it is a multiple
of $\id$. Both operators are traceless: $\sigma(X)$ is a sum of products of two anticommuting generators, and $\Tr(a^+_pa^-_q)=0$ for $p\neq q$, while
$R_{pp}=0$. Hence they are equal. Each term $a^+_pa^-_q$ conserves $N$, and $a^-_q|0\rangle=0$ gives the action on one quantum. (ii)
$N^{(k)}=\mathcal N(E_{kk}\otimes\id_3)$, and by (i) $\sigma(b)=\mathcal N(B\otimes\id_3)$, so $\mathsf H=\mathcal N(H)$. $N^{(k)}$ is real diagonal and
$\sigma(b)$ real antisymmetric, so $\mathsf H$ is Hermitian. For $H=WDW^\dagger$ with $W$ unitary, $b^\pm_r=\sum_pW_{pr}^{(\pm)}a^\pm_p$ ($W^{(+)}=W$,
$W^{(-)}=\bar W$) satisfy the CAR and $\mathsf H=\sum_rd_r\,b^+_rb^-_r$; its eigenvectors are $b^+_{r_1}\cdots b^+_{r_m}|0\rangle$ with $r_1<\dots<r_m$. (iii) is
(i) for $R=(E_{lk}-E_{kl})\otimes\id_3$. Each term moves one quantum between $(k,i)$ and $(l,i)$. (iv) A copy boost anticommutes with $K_3$, so its
lift is a combination of $a^+a^+$ and $a^-a^-$ terms (\cref{e2:lem:N}(ii)). It maps the plane of mode $(k,i)$ to that of mode $(l,i)$, which fixes
the modes paired, and the coefficients are computed. (v) $t_a^2=-\tfrac14$, so
$e^{\theta\hat C^{12}}=\prod_a(\cos\tfrac\theta2+2\sin\tfrac\theta2\,t_a)$. At $\theta=\tfrac\pi2$, $\Ad U$ maps $E_{(1,a)}$ to $\pm E_{(2,a)}$ and $E_{(2,a)}$ to
$\mp E_{(1,a)}$, and fixes $E_{(3,a)}$. $N^{(k)}$ is a function of the products $E_{(k,i)}E_{(k,i+3)}$, so $UN^{(1)}U^{-1}=N^{(2)}$, and conjugation keeps
spectra.
\end{proof}

The theorem is fixed. Part (i) is the standard realisation of $\so{3n}$ on the exterior algebra by one-body operators
\cite{Chevalley54,Fock32}, inherited, and its realisation in the tower is checked exactly for all $15$ and $36$ basis elements at $n=2,3$
(e4-1b) and for the copy rotations at $n=4$ (e4-3a). Part (ii) is checked exactly on $\{N=1\}$ at $n=3$, rates $(1,2,4)$ and
$(B_{12},B_{23},B_{31})=(\tfrac12,\tfrac13,\tfrac15)$, and the additivity on $\{N=2\}$ and $\{N=3\}$ to $10^{-9}$ (e4-1e). Parts (iii) to (v) are exact
(e4-1c, e4-1d, e4-1f); in (iv) a random paired boost of $\so{9,\C}$ also changes $N$ by exactly $\pm2$. The external $i$ in $\mathsf H$ is constructed,
as in \cref{e2:prop:hop}(iii). Under $z\mapsto\sum_pz_pa^+_p|0\rangle$ the one-quantum sector of $e^{-i\mathsf Ht}$ is the frame flow $z(t)=e^{-iHt}z(0)$ of
\cref{d7:thm:schr}(i).

Quadratic fermion dynamics of this kind is standard \cite{Bravyi05}. Nonnegative rates do not make $H$ positive: for the bounded seesaw example,
rates $(0,3,5)$ and couplings $(1,\tfrac12,\tfrac13)$, $H$ has the eigenvalue $-0.326115$, so the empty state is not the lowest state of
$\mathsf H=\mathcal N(H)$, which is bounded below with its negative modes filled. This exterior-algebra model is distinct from the canonical bosonic
system of \cref{bf:krein} and is not its quantisation; the removability of $b$ refers to removing it while keeping the prescribed constant binding
tensor.

\begin{corollary}[Transfer, not creation]
\label{e4:cor:transfer}
Let $b$ be a compact element of $\Cent(G_n)\cap\so{3n,3n}$, the class of gauge-admissible fields between the constituents (\cref{d7:thm:schr}(ii); \cref{cor:noboost}).
\begin{enumerate}[label=(\roman*),leftmargin=2.2em]
\item $e^{t\sigma(b)}$ conserves $N$ and acts on $\{N=m\}$ as the $m$-th exterior power of $e^{tR}$. It moves quanta between modes and
constituents, and never creates or annihilates one.
\item The binding-compatible elements that change $N$ are the paired boosts, among them the copy boosts. Each of their terms creates or
annihilates one pair, and a copy boost puts one quantum in each of the two constituents it joins.
\item Let $\hat U$ lift a compact $\eta_n$-isometry commuting with $G_n$, which acts on $z$ as a unitary $W$. A change of frame by $\hat U$ takes the
rate operator $\mathcal N(\tilde K)$ to $\hat U\mathcal N(\tilde K)\hat U^{-1}=\mathcal N(W\tilde KW^{-1})$. This operator commutes with $N$ and has the same
spectrum on every sector $\{N=m\}$, the sums of $m$ rates. The numbers $N^{(k)}$ are not kept.
\end{enumerate}
\end{corollary}

\begin{proof}
(i) The proof of \cref{e4:thm:fock}(i) holds at every $n$, and $\mathcal N(R)$ is the derivation of the exterior algebra induced by $R$, so its
exponential is $\Lambda(e^{tR})$. (ii) is \cref{e1:thm:Nrule}(ii), with the pairing of \cref{e4:thm:fock}(iv) at $n=3$ and of
\cref{e2:prop:copyboost}(iii) at $n=2$. (iii) $\Ad\hat U$ acts on vectors as the isometry, hence on the $a^+_p$ as $W$, which gives the formula. A compact isometry commuting with $G_n$
commutes with $K_n=-\eta_nG_n$, so $\hat U$ commutes with $N$, and conjugation keeps spectra. This is the Fock form of \cref{d7:thm:schr}(iii), and \cref{e4:thm:fock}(v) exhibits a frame change that moves $N^{(1)}$ to $N^{(2)}$.
\end{proof}

The corollary is fixed. It answers the first question precisely. The gauge-admissible field between constituents transfers quanta and never
creates them. It does not change a rate either (\cref{d7:thm:schr}(iii)). Which constituent carries a quantum depends on the frame
when the copy rotation is gauged, as \cref{e2:prop:hop} found at $n=2$. How many quanta there are, and the multiset of their rates, does not.
Pair creation across constituents needs a copy boost, which is a metric connection and not a gauge field (\cref{d7:prop:boost}).
Because $\mathsf H$ is one-body, its many-quantum lines carry no interaction energy. The field mixes quanta; it does not bind them. This is a
statement about the number-preserving one-body class $\mathcal N(H)$ (the second quantisation $d\Gamma(H)$); it says nothing about extensions
of the framework with two-body terms, which are not constructed here.

% --- source: emergence/e4_bfield_ladder#3 (Time-odd invariants on n constituents)
\subsection{\texorpdfstring{Time-odd invariants on $n$ constituents}{Time-odd invariants on n constituents}}
\label{e4:sec:cycles}

\cref{d7:thm:three} found on three constituents one invariant that two do not carry, the sign of $\plaq=B_{12}B_{23}B_{31}$, odd under interior
time reversal. For a cycle $\gamma=(k_1,\dots,k_L)$ of distinct copy indices write $\Pi_\gamma B=B_{k_1k_2}B_{k_2k_3}\cdots B_{k_Lk_1}$; reversing the
orientation multiplies it by $(-1)^L$.

\begin{theorem}[The copy field on $n$ constituents]
\label{e4:thm:cycles}
Let $n\ge3$, let the rates be distinct, and let $b=B\otimes\id_6$ with every $B_{kl}\neq0$.
\begin{enumerate}[label=(\roman*),leftmargin=2.2em]
\item The endomorphisms of $V_n$ commuting with $G_n$, the diagonal $\lG$ and $K\otimes G$ are $\sum_k(a_k+b_kG)$ on copy $k$, and the $\eta_n$-isometries
among them are the $2^n$ copy reflections $\operatorname{diag}(s_1,\dots,s_n)\otimes\id_6$, $s_k=\pm1$, under which $B_{kl}\mapsto s_ks_lB_{kl}$.
\item Two fields with the same $|B_{kl}|$ are related by a copy reflection iff $\Pi_\gamma B$ has the same sign for both on every cycle $\gamma$. The
signs of the $c_n=\tfrac12(n-1)(n-2)$ triangles $B_{1k}B_{kl}B_{l1}$, $2\le k<l\le n$, are independent and determine all the others. The
invariants of the field are the $\tfrac12n(n-1)$ magnitudes and these $c_n$ signs.
\item Interior time reversal, $(K,B)\mapsto(K,-B)$ (\cref{d7:thm:three}(iv)), multiplies $\Pi_\gamma B$ by $(-1)^L$. It flips all $c_n$ triangle signs,
and no copy reflection undoes it.
\item The triangle invariants are spectral. With $P_k$ the spectral projectors of $K$, polynomials in $K$, and $H=K-iB$ on the copy index,
$\tr(P_kHP_lHP_mH)=iB_{kl}B_{lm}B_{mk}$ for every triangle $(k,l,m)$.
\end{enumerate}
\end{theorem}

\begin{proof}
(i) The commutant of $G_n$ and the diagonal $\lG$ is $\{A\otimes\id+A'\otimes G\}$ with $A,A'\in\mathfrak{gl}(n)$ (\cref{d7:prop:rotor}).
$[A\otimes\id+A'\otimes G,K\otimes G]=[A,K]\otimes G-[A',K]\otimes\id$, so $A$ and $A'$ commute with $K$ and are diagonal. The isometry condition per copy is
$a_kb_k=0$, $a_k^2-b_k^2=1$, as in the proof of \cref{d7:thm:three}(iii). (ii) Encode the signs of $B_{kl}$, $k<l$, as $\sigma\in\Z_2^E$, $E$ the
edges of the complete graph $K_n$. A reflection adds $\delta s$, $(\delta s)_{kl}=s_k+s_l$, so the orbits are the cosets of $\operatorname{im}\delta$.
$\ker\delta$ is the constants because $K_n$ is connected, so $\rank\delta=n-1$ and there are $2^{|E|-n+1}=2^{c_n}$ orbits. The functions of $\sigma$ that are
constant on cosets are the characters $\sigma\mapsto\sum_{e\in\gamma}\sigma_e$ with $\gamma$ in the annihilator of $\operatorname{im}\delta$, the edge sets of even
degree at every vertex, which form the cycle space \cite{Diestel17}. Its dimension is $|E|-n+1=c_n$, and the fundamental cycles of the star at
vertex $1$, the triangles $(1,k,l)$, form a basis. A character pairs perfectly with the cosets, so the triangle signs separate the orbits. The
oriented product $\Pi_\gamma B$ differs from the unoriented one by a sign fixed by $\gamma$. (iii) $-B$ adds the all-ones vector to $\sigma$, and its
character on $\gamma$ is $L\bmod2$. A triangle is odd, so the all-ones vector is not in $\operatorname{im}\delta$. (iv) $P_kHP_l=H_{kl}E_{kl}$ and
$H_{kl}=-iB_{kl}$ for $k\neq l$, and $P_k=\prod_{j\neq k}(K-\kappa_j)/(\kappa_k-\kappa_j)$ for distinct rates.
\end{proof}

The theorem is fixed. Part (i) is exact at $n=4$, rates $(1,2,3,5)$ (e4-2a), and at $n=3$ it is \cref{d7:thm:three}(iii). Parts (ii) and (iii)
are checked exhaustively on all sign patterns at $n=3,4,5,6$, with $c_n=1,3,6,10$ and $\rank\delta=n-1$ over $\Z_2$ (e4-2b, e4-2c). Part (iv) is
symbolic at $n=4$, where the product of the four triangle invariants is a product of squares, because every edge of $K_4$ lies on two
triangles (e4-2d). The count $c_n$ is the cycle rank of $K_n$, inherited from graph theory. It is also the number of rephasing-invariant phases of $n$
families of quarks \cite{KobayashiMaskawa73,BrancoLavouraSilva99} (inherited). Here the phases are fixed at $\pm\pi/2$ because the field is real
(\cref{d7:sec:three}), so they reduce to signs. On $n$ constituents the copy field carries $c_n$ time-odd signs, the holonomies of the
constant copy connection around the triangles of the complete graph of constituents. These are discrete holonomies in copy space, not Wilson
loops of a spacetime gauge field, and ``time-odd'' refers to the interior time reversal $(K,B)\mapsto(K,-B)$. The comparison with quark mixing
is a count of phases only: no mixing matrix between mass and interaction states, no CP operation on a state space and no gauge connection on
spacetime is specified, so the signs are not shown to be physical CP violation.

\begin{remark}[Couplings that vanish]
\label{e4:rem:graphs}
For a general coupling graph $\Gamma_B$, with $|E'|$ nonzero couplings and $c(\Gamma_B)$ connected components, the same argument gives
$2^{|E'|-n+c(\Gamma_B)}$ classes of sign patterns, so the number of independent cycle signs is the cycle rank $|E'|-n+c(\Gamma_B)$ (exhaustive over all
$64$ graphs at $n=4$, rank over $\Z_2$ for all graphs at $n=5,6$; e4-2f). A chain of constituents carries no sign. A ring of $n$ constituents
carries one sign, odd under time reversal iff $n$ is odd, in agreement with the bipartite criterion of \cref{d7:thm:three}(iv). This is fixed.
\end{remark}

\begin{proposition}[The Jarlskog-type determinant up the ladder]
\label{e4:prop:jarlskog}
On the copy index $[H,K]=iD$ with $D_{kl}=(\kappa_k-\kappa_l)B_{kl}$, a real symmetric matrix with zero diagonal. Hence $\det[H,K]=i^n\det D$, and
$\det D$ is the sum over derangements $\pi$ of $(-1)^{n-\#\text{cycles}(\pi)}\prod_kD_{k\pi(k)}$. Under $B\mapsto-B$ it changes by $(-1)^n$.
\begin{enumerate}[label=(\roman*),leftmargin=2.2em]
\item $n=3$: $\det[H,K]=2iJ(\kappa_1-\kappa_2)(\kappa_1-\kappa_3)(\kappa_2-\kappa_3)$.
\item $n=4$: $\det[H,K]=D_{12}^2D_{34}^2+D_{13}^2D_{24}^2+D_{14}^2D_{23}^2-2(D_{12}D_{23}D_{34}D_{41}+D_{12}D_{24}D_{43}D_{31}+D_{13}D_{32}D_{24}D_{41})$. It is even
under time reversal and contains no triangle.
\item $n=5$: $\det[H,K]$ is odd under time reversal; every cycle cover contains a $5$-cycle or a triangle.
\end{enumerate}
\end{proposition}

\begin{proof}
$[K-iB,K]=-i[B,K]$ and $[B,K]_{kl}=-(\kappa_k-\kappa_l)B_{kl}$. $D$ is a product of two antisymmetric factors entrywise, hence symmetric. The Leibniz
expansion of a matrix with zero diagonal runs over derangements, and $-D$ multiplies each term by $(-1)^n$. At $n=4$ the derangements are the three
double transpositions and the six $4$-cycles. At $n=5$ they are the $24$ five-cycles and the $20$ products of a transposition and a $3$-cycle.
\end{proof}

The proposition is fixed. Part (i) restates \cref{d7:thm:three}(ii), and the identities are symbolic at $n=3,4,5$, with the cycle-cover
sum checked at $n=4,5$ (e4-2e). The determinant measures cycle covers. A $2$-cycle contributes the magnitude $(\kappa_k-\kappa_l)^2B_{kl}^2$, and a
cycle of length $L\ge3$ contributes twice the product of the rate differences around it times $\Pi_\gamma B$. At $n=4$ the $4$-cycle products are
time-even: $B_{12}B_{23}B_{34}B_{41}=-B_{12}B_{23}B_{31}\cdot B_{13}B_{34}B_{41}/B_{13}^2$ is fixed by the product of two triangle signs. So
$\det[H,K]$ cannot tell $(K,B)$ from $(K,-B)$ at $n=4$, although the field is time-odd. The time-odd invariants there are the triangle traces of
\cref{e4:thm:cycles}(iv). Jarlskog's determinant \cite{Jarlskog85} measures time reversal on three constituents and, among the rungs, only at odd
$n$; the triangle traces measure it at every $n\ge3$.

% --- source: emergence/e4_bfield_ladder#4 (Copy isospin at four constituents)
\subsection{Copy isospin at four constituents}
\label{e4:sec:isospin}

On two and three constituents the copy isospin is integral (\cref{d7:prop:isospin}), and \cref{d7:rem:verdict}(d) notes that this does
not extend to four (d7-6d). We determine the copy isospin content of $S_4$, of dimension $4096$. Write
\[
I^\pm_1=\half(\hat C^{23}\pm\hat C^{14}),\qquad I^\pm_2=\half(\hat C^{31}\pm\hat C^{24}),\qquad I^\pm_3=\half(\hat C^{12}\pm\hat C^{34})
\]
(the self-dual and anti-self-dual combinations; $J^\pm_a$ in the script), and let $j_\pm(j_\pm+1)$ be the eigenvalues of the Casimirs
$-\sum_a(I^\pm_a)^2$.

\begin{theorem}[Copy isospin on four constituents]
\label{e4:thm:so4}
Let $n=4$.
\begin{enumerate}[label=(\roman*),leftmargin=2.2em]
\item Each $\hat C^{kl}$ equals $\sum_i(a^+_{l,i}a^-_{k,i}-a^+_{k,i}a^-_{l,i})$, commutes with every $N_{[i]}$, and is a sum of six commuting terms of square
$-\tfrac14$, so $e^{2\pi\hat C^{kl}}=+\id$. The six close on $\so4$, and the copy group acting on $S_4$ is $SO(4)$.
\item $[I^+_a,I^-_b]=0$ and $[I^\pm_1,I^\pm_2]=I^\pm_3$ cyclically, on vectors and on $S_4$. On vectors $4(I^\pm_a)^2=-\id$, so the copy index is the module
$(\half,\half)$. On $S_4$, $e^{2\pi I^+_3}=e^{2\pi I^-_3}=(-1)^N$. Hence $I^+_3=\half(\hat C^{12}+\hat C^{34})$ has eigenvalues $i\mu$ with $\mu$ half-integral
exactly on the sectors of odd $N$. In particular $e^{2\pi I^+_3}a^+_{1,1}|0\rangle=-a^+_{1,1}|0\rangle$.
\item As an $\so4$-module, $S_4\otimes\C\cong(\Lambda\C^4)^{\otimes3}$ with $\Lambda\C^4=2(0,0)\oplus2(\half,\half)\oplus(1,0)\oplus(0,1)$. The modules $(j_+,j_-)$ and their
multiplicities are those of \cref{e4:tab:so4}. Always $j_+-j_-\in\Z$. The modules with half-integral $j_+$ fill exactly the odd-$N$ sectors, $2048$
dimensions. The one-quantum sector is $(\half,\half)\otimes\R^3$.
\end{enumerate}
\end{theorem}

\begin{table}[htbp]
\centering
\footnotesize
\begin{tabular}{@{}lcccccccccc@{}}
\toprule
even $N$: $(j_+,j_-)$ & $(0,0)$ & $(1,0)$ & $(0,1)$ & $(1,1)$ & $(2,0)$ & $(0,2)$ & $(2,1)$ & $(1,2)$ & $(3,0)$ & $(0,3)$\\
multiplicity & $70$ & $84$ & $84$ & $90$ & $20$ & $20$ & $15$ & $15$ & $1$ & $1$\\
\midrule
odd $N$: $(j_+,j_-)$ & $(\half,\half)$ & $(\tfrac32,\half)$ & $(\half,\tfrac32)$ & $(\tfrac32,\tfrac32)$ & $(\tfrac52,\half)$ & $(\half,\tfrac52)$ & & & &\\
multiplicity & $140$ & $64$ & $64$ & $20$ & $6$ & $6$ & & & &\\
\bottomrule
\end{tabular}
\caption{The copy isospin content of $S_4\otimes\C$ under $\so4=\su2_+\oplus\su2_-$. Each half has dimension $2048$ (e4-3c).}
\label{e4:tab:so4}
\end{table}

\begin{proof}
(i) The one-body form is \cref{e4:thm:fock}(i), and the six terms are those of \cref{d7:prop:isospin}(i) with $|D|=6$. A module of $\so4$ on which
the $2\pi$ rotation in one plane acts as $\id$ integrates to $SO(4)$, not to its double cover. (ii) On vectors $2I^\pm_3=(c^{12}\pm c^{34})$ squares to
$-\id$, so $e^{2\pi I^\pm_3}=-\id$ on $V_4$. Its lift is the product of the twelve commuting factors $e^{2\pi u}$, where $u$ runs over the terms of
$I^\pm_3$, $u^2=-\tfrac1{16}$ and $e^{2\pi u}=4u$. That product is a multiple of the volume element of $\Cl(12,12)$, and the computation gives $(-1)^N$. On a
sector where $e^{2\pi I^+_3}=-\id$ the eigenvalues $i\mu$ of $I^+_3$ satisfy $e^{2\pi i\mu}=-1$. (iii) $S_4$ is the exterior algebra on the modes, and the
modes of direction $i$ span a copy $\C^4_i$ of the copy index, so $S_4\otimes\C=\Lambda(\C^4_1\oplus\C^4_2\oplus\C^4_3)=\bigotimes_i\Lambda\C^4_i$. By (i), $\so4$ acts
on each factor as the derivation induced by the vector module. $\Lambda^0=\Lambda^4=(0,0)$, $\Lambda^1\cong\Lambda^3=(\half,\half)$ and
$\Lambda^2=\Lambda^2_+\oplus\Lambda^2_-=(1,0)\oplus(0,1)$ \cite{Humphreys72}. The Clebsch--Gordan rule gives the table. Odd degree in $\Lambda\C^4$ means half-integral
$(j_+,j_-)$, and the parity of a product is the parity of the total degree, $(-1)^N$.
\end{proof}

The theorem is fixed. (i) and (ii) are exact on $S_4$, with sparse integer matrices (e4-3a, e4-3b), and $e^{2\pi\hat C^{12}}=+\id$ also on $S_5$, of dimension
$32768$ (e4-3a). (iii) is checked on the joint spectrum of the two Casimirs, block-diagonalised by the direction numbers $N_{[i]}$, which commute with
every copy rotation; the eigenvalues are integral to $10^{-8}$ and the multiplicities agree with the character (e4-3c). The decomposition rule is
inherited. The same rule at $n=3$, $\Lambda\C^3=2(j{=}0)\oplus2(j{=}1)$, gives the Casimir spectrum $j=0,1,2,3$ on $40,216,200,56$ dimensions of
\cref{d7:prop:isospin}(ii), recomputed here (e4-3d, restated).

\begin{corollary}[No spinor of the copy algebra at any rung]
\label{e4:cor:nospinor}
For every $n$, $e^{2\pi\hat C^{kl}}=+\id$ on $S_n$, and $S_n\otimes\C\cong(\Lambda\C^n)^{\otimes3}$ as a module of the copy $\so n$. The copy algebra acts through
$SO(n)$, and for $n\ge3$ no spinor module of $\so n$ occurs. At $n=4$ the half-integral isospin is that of the vector $(\half,\half)$ and of its odd
exterior powers. At $n=6$, where $\so6\cong\su4$, the fundamental $\mathbf 4$ does not occur.
\end{corollary}

\begin{proof}
The six terms of $\hat C^{kl}$ commute and square to $-\tfrac14$ at every $n$, and the proof of \cref{e4:thm:so4}(iii) holds with $\C^n$ in place of $\C^4$.
Exterior powers of the vector module have integral weights in the orthonormal basis, and the spinor modules have half-integral ones
\cite{LawsonMichelsohn89}.
\end{proof}

The corollary is fixed, checked on $S_4$ and $S_5$ (e4-3a). It makes the doublets of \cref{d7:rem:verdict}(d) precise. At $n=4$ a state carries
half-integral copy isospin under either $\su2$ factor iff it has an odd number of quanta; the copy inversion $-\id\in SO(4)$, the rotation through $\pi$
in the planes $(1,2)$ and $(3,4)$, acts as $(-1)^N$. This is possible without any double cover because $\so4$ is not simple and its vector is a
doublet of each factor. Up the ladder the copy algebra $\so n$ is simple again ($n\ge5$), so it has no $\su2$ ideal. Its subalgebra $\so4$ on any four
constituents still has doublets: on $S_5$, $e^{2\pi I^\pm_3}$ built on the constituents $(1,2,3,4)$ or $(2,3,4,5)$ is $(-1)^{N'}$, with $N'$ the number of
quanta in those four constituents (e4-3e). Half-integral copy isospin therefore first appears at the fourth rung and persists above it, under
the $\so4$ of four constituents. The doublets are frames and their odd products, not spinors of a spin cover.

% --- source: emergence/e4_bfield_ladder#5 (Rate patterns and the breaking of the copy symmetry)
\subsection{Rate patterns and the breaking of the copy symmetry}
\label{e4:sec:rates}

\begin{proposition}[Rate patterns]
\label{e4:prop:rates}
Let the rates take $r$ distinct values with multiplicities $n_1+\dots+n_r=n$.
\begin{enumerate}[label=(\roman*),leftmargin=2.2em]
\item The removable part of the copy $\so n$, the part commuting with $K\otimes G$, is $\bigoplus_i\so{n_i}$, block diagonal on the equal-rate
constituents, of dimension $\sum_i\half n_i(n_i-1)$. The non-removable directions number $\half n(n-1)-\sum_i\half n_i(n_i-1)=\sum_{i<j}n_in_j$. Within the
compact binding-compatible $\so{3n}$ the removable part is $\bigoplus_i\so{3n_i}$, and the non-removable directions number $9\sum_{i<j}n_in_j$.
\item The quadratic form $B\mapsto\|[K,B]\|^2$ on $\so n$ is diagonal in the basis $c^{kl}$ (on the copy index), with eigenvalue $2(\kappa_k-\kappa_l)^2$.
\end{enumerate}
\end{proposition}

\begin{proof}
(i) $[K,E_{lk}-E_{kl}]=(\kappa_l-\kappa_k)(E_{lk}+E_{kl})$, so $X\in\so n$ commutes with $K$ iff $X_{kl}=0$ whenever $\kappa_k\neq\kappa_l$, and likewise for $R$ and
$\tilde K$ on $\C^{3n}$. Removability is commutation with $K\otimes G$ (\cref{d7:thm:schr}(v)). The count uses $n^2=\sum n_i^2+2\sum_{i<j}n_in_j$. (ii) The
matrices $E_{lk}+E_{kl}$, $k<l$, are orthogonal, of norm squared $2$.
\end{proof}

The proposition is fixed, exact for every partition of $n=2,\dots,5$ in $\so n$ and of $n=2,3,4$ in $\so{3n}$ (e4-4a), and at rates $(1,1,2,5)$, where the
form is $\operatorname{diag}(0,2,32,2,32,18)$ on $c^{12},c^{13},c^{14},c^{23},c^{24},c^{34}$ (e4-4b). At $n=3$ it restates the dimensions $36,18,9$ and $3,1,0$ of
\cref{d7:thm:schr}(v). With the factor $6$ of $\id_6$, (ii) is the norm $6\sum_{k\neq l}(\kappa_k-\kappa_l)^2B_{kl}^2$ of \cref{d7:rem:verdict}(e).

\begin{remark}[The breaking pattern]
\label{e4:rem:higgs}
The rate matrix is a symmetric tensor of the copy $\so n$, and its stabiliser is $\bigoplus_i\so{n_i}$; its stabiliser in $O(n)$ is
$O(n_1)\times\dots\times O(n_r)$, which for distinct rates is the group of copy reflections of \cref{e4:thm:cycles}(i). Read as in \cref{d7:rem:verdict}(e), this is the pattern of a gauge symmetry broken by a symmetric-tensor vacuum value \cite{LiLF74}, with $\sum_{i<j}n_in_j$ massive directions
of masses proportional to $|\kappa_k-\kappa_l|$ \cite{EnglertBrout64,Higgs64,Kibble67}. Up the ladder, constituents with equal rates form a multiplet
whose internal copy rotations are removable, so which of them carries a quantum is a gauge choice. Constituents with different rates are
distinguishable, and the field between them is physical. The reading is interpretive, with the caveat of \cref{d7:sec:assess}: $b$ is a constant connection
without kinetic term, and the rates are inputs.
\end{remark}

% --- source: emergence/e4_bfield_ladder#6 (A ladder of rates admits a non-abelian frame)
\subsection{A ladder of rates admits a non-abelian frame}
\label{e4:sec:rotor}

\cref{d7:sec:rotor} shows that a common binding action lets the rotor act by a non-central $A\in\uu n$. It also shows that the frame then
admits equivariant brackets $[X,Y]_k=\sum_{i,j}\Gamma^k_{ij}\beta_0(X_i,Y_j)$, with $\beta_0(x,y)=z(x)\times z(y)$ the complex cross product and
$\Gamma^k_{ij}=\Gamma^k_{ji}\in\C$, subject to $\Gamma^k_{ij}(\omega_k-\omega_i-\omega_j)=0$, where $A$ acts on copy $k$ as $i\omega_k$. For $A=K\otimes G$ the
$\omega_k$ are the rates. With a copy field, $A$ acts on $z$ as $iH$, and the condition applies to the eigenvalues of $H$, the lines.

\begin{proposition}[Rotor-compatible frame brackets]
\label{e4:prop:rotor}
\begin{enumerate}[label=(\roman*),leftmargin=2.2em]
\item The real bilinear maps $V\times V\to V$ equivariant under $\lG$ form the two-dimensional space spanned by $\beta_0$ and $G\beta_0$, and all of them
are antisymmetric. The equivariant antisymmetric brackets on $V_n$ are the maps above, a space of real dimension $n^2(n+1)=12,36,80$ for $n=2,3,4$.
\item $A=K\otimes G$ is a derivation of the bracket iff $\Gamma^k_{ij}=0$ unless $\kappa_k=\kappa_i+\kappa_j$. The rotor-compatible brackets have real dimension
$0$, $0$, $2$ for rates $(1,1)$, $(1,3)$, $(1,2)$; $0$, $4$, $0$, $4$ for $(1,1,1)$, $(1,2,3)$, $(1,3,5)$, $(1,2,4)$; and $8$ for $(1,2,3,4)$. These count the
resonant coefficients; for $(1,2,3,4)$ the Jacobi identity ($ad=bc$, item (v)) leaves a bracket variety of real dimension $6$. For the
harmonic ladder $\kappa_k=k$ the dimension is $2\sum_{k\le n}\lfloor k/2\rfloor=2,4,8,12,18$ for $n=2,\dots,6$.
\item Rates $(1,2,3)$. The rotor-compatible brackets are $[X,Y]_2=a\,\beta_0(X_1,Y_1)$, $[X,Y]_3=b\,(\beta_0(X_1,Y_2)+\beta_0(X_2,Y_1))$, $(a,b)\in\C^2$. Every one of them
satisfies the Jacobi identity. For $ab\neq0$ the frame is a nilpotent Lie algebra $V_1\oplus V_2\oplus V_3$, graded by the rate and generated by $V_1$, with
lower central series of real dimensions $18,12,6,0$.
\item Rates $(1,2,4)$. The rotor-compatible brackets are $[X,Y]_2=a\,\beta_0(X_1,Y_1)$, $[X,Y]_3=b\,\beta_0(X_2,Y_2)$. For $X,Y\in V_1$ and $Z\in V_2$ the
Jacobiator is $ab\,\beta_0(\beta_0(X,Y),Z)\in V_3$, so the Jacobi identity holds iff $ab=0$.
\item Rates $(1,2,3,4)$, with $a=\Gamma^2_{11}$, $b=\Gamma^3_{12}$, $c=\Gamma^4_{13}$, $d=\Gamma^4_{22}$. For $X,Y\in V_1$ and $Z\in V_2$ the Jacobiator is
$(ad-bc)\,\beta_0(\beta_0(X,Y),Z)\in V_4$, and the Jacobi identity holds on all of $V_4$ iff $ad=bc$.
\end{enumerate}
\end{proposition}

\begin{proof}
(i) is an exact null-space computation on the $216$ coefficients of a bilinear map (e4-5a); $G$ acts as $i$, so $G\beta_0=i\beta_0$, and the maps are complex
bilinear. A bracket is a choice of $\Gamma^k_{ij}$ for each $k$ and each unordered pair $\{i,j\}$, $n\cdot\half n(n+1)$ complex numbers. (ii) On $z$,
$A$ acts as $i\kappa_k$ on copy $k$, and the bracket is complex bilinear, so $A[X,Y]_k-[AX,Y]_k-[X,AY]_k=i(\kappa_k-\kappa_i-\kappa_j)\Gamma^k_{ij}\beta_0(X_i,Y_j)$. For
$\kappa_k=k$ the pairs $i\le j$ with $i+j=k$ number $\lfloor k/2\rfloor$. (iii) The only composite is $[[X,Y]_2,Z]_3$, so the Jacobiator is
$ab\sum_{\rm cyc}\beta_0(\beta_0(X_1,Y_1),Z_1)$, which vanishes by the Jacobi identity of the complex cross product, the bracket of $\so{3,\C}$.
$\beta_0(V_1,V_1)=V_2$ and $\beta_0(V_1,V_2)=V_3$ for $ab\neq0$. The Jacobiator is trilinear and alternating, and inputs of rates $r_1,r_2,r_3$ give an output of rate $r_1+r_2+r_3$, so each of its components is fixed
by inputs whose rates add to the target rate; for $a\neq0=b$ the series is $18,6,0$ (a12-3). (iv) The only composite is $[[X,Y]_2,Z]_3=ab\,\beta_0(\beta_0(X_1,Y_1),Z_2)$, with inputs of rates $1,1,2$. For $X,Y\in V_1$ and $Z\in V_2$
the other two cyclic terms vanish, because $V_1\times V_2$ brackets to zero. For $ab=0$ there is no composite. (v) The $V_3$-component is as in (iii)
and vanishes. For $X,Y\in V_1$ and $Z\in V_2$, $[[X,Y],Z]_4=ad\,\beta_0(\beta_0(X,Y),Z)$, while
$[[Y,Z],X]_4+[[Z,X],Y]_4=bc\,(\beta_0(\beta_0(Y,Z),X)+\beta_0(\beta_0(Z,X),Y))=-bc\,\beta_0(\beta_0(X,Y),Z)$ by the Jacobi identity of $\beta_0$.
\end{proof}

The proposition is fixed. (i) restates the space of \cref{d7:sec:rotor} at $n=2$ (e4-5a). (ii) is exact, block by block, for the listed rates and
for the harmonic ladder up to $n=4$; for $n=5,6$ the count is the number of resonant blocks (e4-5b). At $n=2$ it restates the bracket of
$\operatorname{diag}(1,2)\otimes G$ of \cref{d7:sec:rotor}. (iii) to (v) are checked on random samples to $10^{-10}$ (e4-5c, e4-5d, e4-5e):
antisymmetry, $\lG$-equivariance, the rotor as a derivation and the Jacobi identity; at $(1,2,4)$ and $(1,2,3,4)$ the samples off the stated
conditions fail. That every rotor-compatible frame is nilpotent when all rates are positive is inherited: $A$ is then a nonsingular derivation, and a
Lie algebra with a nonsingular derivation is nilpotent \cite{Jacobson55}.

The graded brackets of (iii)--(v) are of current type, $\mathfrak g\otimes A$ with $\mathfrak g$ the complex cross product and $A$ a graded commutative
coefficient algebra without unit, $t_it_j=\sum_k\Gamma^k_{ij}t_k$; Lie algebras extended over commutative algebras without unit are treated by
Zusmanovich \cite{Zusmanovich10}. The truncated current algebra $\mathfrak g\otimes t\C[t]/(t^{n+1})$ is a particular case, not the general one: at
rates $(1,2,3)$ it is the case $ab\neq0$ of (iii), up to rescaling, while $a\neq0$, $b=0$ gives a two-step nilpotent frame with $V_3$ central and
decoupled, against the three-step lower central series of the truncation.

The reading is constructed. A harmonic ladder of rates admits a non-abelian frame, from $n=3$ on and already for $\operatorname{diag}(1,2)$ at $n=2$. Its
bracket joins two frames whose rates add into the frame of the summed rate, so the frame algebra is graded by rate. With three constituents
on the harmonic ladder the Lie frame brackets form the linear space $\C^2$. With four the Jacobi identity cuts out the cone $ad=bc$ in $\C^4$, and
on a longer ladder the same condition recurs, because the $V_4$-component of the Jacobiator involves only the blocks with target rate at most
$4$. The resonant pattern $(1,2,4)$ admits a non-abelian frame only with one of its two blocks, $ab=0$. Nothing selects such a rotor. The paper's closure takes $A=G_n$ and an abelian frame
(\cref{thm:closureunique} and \cref{prop:several}). The non-central rotor, its rate ladder and the bracket are admitted, not selected. The
resonance condition $\omega_k=\omega_i+\omega_j$ is that of three-wave mixing and sum-frequency generation \cite{Armstrong62}, a parallel only
(interpretive). Read with $mc^2=\hbar\kappa$, the bracket composes frames whose rest energies add (interpretive).

% --- source: emergence/e4_bfield_ladder#7 (Answer: what the field between constituents can and cannot d)
\subsection{Answer: what the field between constituents can and cannot do}
\label{e4:sec:answer}

\begin{remark}[What the field does up the ladder]
\label{e4:rem:answer}
\begin{enumerate}[label=(\alph*),leftmargin=2.2em]
\item \emph{It moves quanta and keeps their number} (fixed). A gauge-admissible field between the constituents acts as a one-body operator.
It hops quanta between constituents and mixes them, and its $N$-quantum lines are sums of one-quantum lines (\cref{e4:thm:fock} and
\cref{e4:cor:transfer}).
\item \emph{It does not take one to two} (fixed). It creates no quantum and changes no rate. Pair emergence across constituents needs a copy boost,
which is not a gauge field (\cref{e4:cor:transfer}(ii); \cref{d7:prop:boost}). Because it is one-body, it binds no quanta into composites;
interactions beyond this one-body class are not examined.
If two particles are to emerge from one, the process lies outside the compact field between constituents.
\item \emph{It carries time-odd invariants that grow with the rung} (fixed). From $n=3$ on it carries $\tfrac12(n-1)(n-2)$ independent triangle signs, each odd
under interior time reversal. The triangle traces measure them at every $n$, Jarlskog's determinant only at odd $n$
(\cref{e4:thm:cycles,e4:prop:jarlskog}).
\item \emph{It acts on doublets from the fourth rung on} (fixed). At $n=4$ the copy isospin is half-integral exactly on odd $N$, and above it under the
$\so4$ of any four constituents, on an odd number of quanta in them. At every rung the copy algebra acts through $SO(n)$, so no spinor of it
occurs (\cref{e4:thm:so4,e4:cor:nospinor}).
\item \emph{It is non-removable exactly between unequal rates} (fixed), with the prescribed binding tensor kept fixed. The non-removable directions number $\sum_{i<j}n_in_j$; equal-rate constituents stay
gauge-equivalent (\cref{e4:prop:rates}). The reading as symmetry breaking by the rate matrix is interpretive (\cref{e4:rem:higgs}).
\item \emph{It admits non-abelian frames on a harmonic ladder of rates} (constructed). The rotor must be non-central and its rates resonant, and
nothing in the closure selects either (\cref{e4:prop:rotor}).
\end{enumerate}
Up the ladder, then, the field between constituents distinguishes constituents by their rates, carries more time-odd structure, and from
$n=4$ on acts on doublets. It never changes the number of quanta that the rung carries. In the sense of \cref{e1:sec:fock}, a rung carries $N$
quanta, and the compact $b$-field keeps $N$ fixed. It decides where the quanta are, not how many there are. This reading is interpretive, drawn
from the fixed statements above.
\end{remark}

\begin{table}[htbp]
\centering
\footnotesize
\begin{tabular}{@{}>{\raggedright\arraybackslash}p{0.4cm}>{\raggedright\arraybackslash}p{2.6cm}>{\raggedright\arraybackslash}p{2.0cm}>{\raggedright\arraybackslash}p{1.8cm}>{\raggedright\arraybackslash}p{3.8cm}>{\raggedright\arraybackslash}p{2.3cm}@{}}
\toprule
$n$ & copy algebra, group on $S_n$ & non-removable directions, distinct rates & independent signs $c_n$ & half-integral copy isospin & rotor brackets, rates $1,\dots,n$: resonant coefficients\\
\midrule
$2$ & $\so2$, $SO(2)$ & $1$ & $0$ & no; charges $i\{-3,\dots,3\}$ & $2$\\
$3$ & $\so3$, $SO(3)$ & $3$ & $1$ & no; $j=0,\dots,3$ & $4$\\
$4$ & $\su2\oplus\su2$, $SO(4)$ & $6$ & $3$ & yes, exactly on odd $N$ & $8$, Jacobi: $ad=bc$\\
$5$ & $\so5$, $SO(5)$ & $10$ & $6$ & under the $\so4$ of four constituents; no spinor of $\so5$ & $12$\\
$6$ & $\so6\cong\su4$, $SO(6)$ & $15$ & $10$ & as for $n=5$; no $\mathbf 4$ of $\su4$ & $18$\\
\bottomrule
\end{tabular}
\caption{The field between constituents, rung by rung. The $n=2,3$ isospin entries restate \cref{d7:prop:isospin}. The others follow from
\cref{e4:thm:cycles,e4:thm:so4,e4:cor:nospinor,e4:prop:rates,e4:prop:rotor}; the rotor column gives the real dimension of the space of resonant
coefficients before the Jacobi identity is imposed (at $n=4$ the Jacobi identity cuts it to the variety $ad=bc$; at $n=5,6$ the Lie-bracket varieties are
not determined here). Non-removable: with the prescribed binding tensor kept fixed.}
\label{e4:tab:ladder}
\end{table}

\begin{table}[htbp]
\centering
\footnotesize
\begin{tabular}{@{}>{\raggedright\arraybackslash}p{8.0cm}>{\raggedright\arraybackslash}p{4.8cm}>{\raggedright\arraybackslash}p{2.0cm}@{}}
\toprule
statement & class & checks\\
\midrule
\cref{e4:thm:fock}(i): paired rotations are one-body operators $\mathcal N(R)$, conserve $N$ & fixed (structure inherited, realisation fixed) & e4-1a, e4-1b\\
\cref{e4:thm:fock}(ii): $\mathsf H=\mathcal N(H)$; $N=1$ block is $H$; additive lines & fixed; external $i$ constructed & e4-1e\\
\cref{e4:thm:fock}(iii), (iv): copy rotations and copy boosts at $n=3$ by constituent and direction & fixed & e4-1c, e4-1d\\
\cref{e4:thm:fock}(v): quarter turn moves $N^{(1)}$ to $N^{(2)}$, keeps the rate spectrum & fixed & e4-1f\\
\cref{e4:cor:transfer}: transfer, not creation; rates gauge invariant & fixed (partly restated: \cref{e1:thm:Nrule}; \cref{d7:thm:schr}(iii)) & e4-1b, e4-1d, e4-3a\\
\cref{e4:thm:cycles}(i): residual group is the copy reflections & fixed (restated at $n=3$: \cref{d7:thm:three}(iii)) & e4-2a\\
\cref{e4:thm:cycles}(ii), (iii): invariants are $|B_{kl}|$ and $c_n$ triangle signs; all time-odd & fixed (cycle space inherited) & e4-2b, e4-2c\\
\cref{e4:thm:cycles}(iv): triangle traces & fixed & e4-2d\\
\cref{e4:rem:graphs}: cycle rank of the coupling graph & fixed & e4-2f\\
\cref{e4:prop:jarlskog}: $\det[H,K]=i^n\det D$, parity $(-1)^n$, $n=3,4,5$ & fixed ($n=3$ restated: \cref{d7:thm:three}(ii)) & e4-2e, z1-4a\\
parallel with the $(n-1)(n-2)/2$ phases of $n$ quark families & count inherited; physical CP reading not established & ---\\
\cref{e4:thm:so4}(i), (ii): $SO(4)$, $(\half,\half)$, $e^{2\pi I^\pm_3}=(-1)^N$ & fixed & e4-3a, e4-3b\\
\cref{e4:thm:so4}(iii), \cref{e4:tab:so4}: $(\Lambda\C^4)^{\otimes3}$, doublets on odd $N$ & fixed (decomposition rule inherited) & e4-3c\\
$n=3$ Casimir spectrum $40,216,200,56$ & restated (\cref{d7:prop:isospin}(ii)) & e4-3d\\
\cref{e4:cor:nospinor}: no spinor of the copy algebra at any rung; doublets of the $\so4$ of four constituents at $n=5$ & fixed & e4-3a, e4-3e\\
\cref{e4:prop:rates}: removable $\bigoplus\so{n_i}$, $\sum_{i<j}n_in_j$ physical; mass form & fixed ($n=3$ restated: \cref{d7:thm:schr}(v)) & e4-4a, e4-4b\\
\cref{e4:rem:higgs}: breaking by the rate matrix & interpretive & ---\\
\cref{e4:prop:rotor}(i), (ii): equivariant brackets, resonance condition, dimensions & fixed ($n=2$ restated: \cref{d7:sec:rotor}) & e4-5a, e4-5b\\
\cref{e4:prop:rotor}(iii)--(v): Jacobi at $(1,2,3)$, $(1,2,4)$, $(1,2,3,4)$; nilpotency & fixed; nilpotency inherited & e4-5c, e4-5d, e4-5e\\
\cref{e4:sec:rotor}: a harmonic ladder admits a non-abelian frame & constructed; three-wave parallel interpretive & e4-5b, e4-5c\\
\cref{e4:rem:answer}, \cref{e4:tab:ladder}: what the field can and cannot do & (a)--(e) fixed, (f) constructed; summary interpretive & ---\\
\bottomrule
\end{tabular}
\caption{Claims of \cref{e4:sec:main}, their classes and their checks, which are \texttt{PASS} lines of \texttt{e4\_bfield\_ladder.py}.}
\label{e4:tab:claims}
\end{table}

% --- source: emergence/e5_closure_two#0 ((head))
% Track e5 -- The closure of the force law at two constituents (n = 2).
% New references: sections/bib_e5.tex.

% --- source: emergence/e5_closure_two#1 (The closure at two constituents)
\section{The closure at two constituents}
\label{e5:sec:main}

\cref{sec:closure} closes the force law $\dot p=\kappa G^{\!\top}p$ into the Lie--Poisson system of $\mathfrak s_n=\R\ltimes_{G_n}V_n$,
in which the rate is the angular velocity of a rotor. It treats one constituent, several constituents driven by one rotor (\cref{prop:several}), and three constituents (\crefrange{thm:so9C}{thm:inertia3}): the closure symmetry, the absence of a common Lorentz
algebra, the part of the Pati--Salam algebra that commutes with the binding root, what survives each projector, its chirality and
anomalies, and the freedom left in the inertia. The case $n=2$ is not treated there. Parts of it are in \cref{ch:ch05_tower}: the closure symmetry $\so{6,\C}$ (\cref{d5:prop:so6C}), the balance surface (\cref{d6:thm:balance}), the balanced pairs (\cref{d6:prop:pairs}) and the pair invariant $M_{\rm pair}$ (\cref{d6:thm:pair}). This section supplies the $n=2$ analogue of every item, with the five Lorentz actions at $n=2$ in place of the two
projectors, which exist only at $n=3$, and asks whether $n=2$ is a closure rung in its own right. The answers, with their classes, are these. (1) The closure symmetry $\so{6,\C}$ is the complex orthogonal algebra of
$q_{\rm tot}=z_1\cdot z_1+z_2\cdot z_2$. It mixes the two field strengths, keeps no quadratic invariant except $q_{\rm tot}$, and does not
preserve the pair invariant: on each of its orbits $M_{\rm pair}$ takes every value in $[\half|q_{\rm tot}|,\infty)$ (fixed;
\cref{e5:sec:symmetry}). (2) There is no common Lorentz algebra, and under $\lG^{(1)}$ the spinor module carries no Weyl spinor
(fixed, partly restated; \cref{e5:sec:nocommon}). (3) None of the compact parts of the five candidate gauge algebras of the Lorentz actions at $n=2$ contains the Pati--Salam algebra (obstructed within that tested class; other constructions at $n=2$ are not excluded). The binding root leaves an
$\su2$ of the spectator gauge algebra, compact, non-abelian and kinetically admissible, which acts on the Weyl block with integer
isospin and vector-like; of the diagonal connection algebra it leaves the copy rotation; of the mixed rung it leaves nothing (fixed;
\cref{e5:sec:gauge}). (4) The inertia freedom is $(2,0)$ under the closure symmetry, as at $n=1$ and $n=3$, and it is $(14,10)$ or $(18,18)$
under the binding-compatible gauge parts (fixed, partly restated; \cref{e5:sec:inertia}). (5) The balance surface is connected and of
codimension two. It is transversally neutral, since the imbalance is a constant of motion; the rate is stable, but the bounded frame orbits
are not, and they need $h(0)\neq0$. With equal inertias a massive balanced pair is never an exchange eigenstate of its momenta
(fixed; \cref{e5:sec:rotor}).
(6) \Cref{e5:tab:ladder} compares $n=1,2,3$, and \cref{e5:sec:ladder} answers the question. Identifiers e5-$n$ refer to the \texttt{PASS}
lines of \texttt{e5\_closure\_two.py} (\cref{e5:tab:claims}).

\paragraph{Notation.} As in \cref{d6:sec:main}. $V_2=V\oplus V$ has the basis $e^{(1)}_a$, $e^{(2)}_a$, the metric $\eta_2=\eta\oplus\eta$ and
the binding root $G_2=G\oplus G$. The temporal directions are $e^{(k)}_1,e^{(k)}_2,e^{(k)}_3$, each with its partner $e^{(k)}_{i+3}$, and the
paired rotations and boosts $\mathrm{rot}(i,j)$, $\mathrm{bst}(i,j)$ of temporal $i<j$ are those of \cref{sec:closurethree}. A vector $X$ of
$V_2$ gives $Z=(z_1,z_2)\in\C^6$, $z_k=z(X^{(k)})$, on which $G_2$ acts as $i$; $q_k=z_k\cdot z_k$, $q_{\rm tot}=q_1+q_2$, $Q_{kl}=z_k\cdot z_l$ is the
complex Gram matrix with singular values $\sigma_1\ge\sigma_2$, and $M_{\rm pair}=\half(\sigma_1+\sigma_2)$ is the rest-energy density of the pair, a local
invariant of the summed stress tensor and not an invariant mass (\cref{d6:thm:pair}). $C_v=\varepsilon\otimes\id_6$ is the copy rotation and $\varepsilon\otimes G$ the copy boost,
$\varepsilon=\bigl(\begin{smallmatrix}0&-1\\1&0\end{smallmatrix}\bigr)$. $\lG^{(k)}$ and $\lint^{(k)}$ are the frame and interior Lorentz
algebras of constituent $k$, and $\lG^\Delta$ is the diagonal $\lG$. On $S_2=\R^{64}$ the tower generators are $E_{(1,a)}=e_a\otimes\id$,
$E_{(2,a)}=\omega\otimes e_a$, the spin lift of $M\in\so{6,6}$ is $S(M)=\tfrac14\sum_{a,b}(M\eta_2)_{ab}E_aE_b$, and $\Gamma_{12}=E_{(1,1)}\cdots E_{(2,6)}$.
In the closure the per-constituent quantities are those of the momentum fields $\pi^{(k)}=\eta^{-1}p^{(k)}$, $w_k=w(p^{(k)})=q(z(\pi^{(k)}))$
(equation~\cref{eq:wdef}); Lorentz-class inertias are $N_k=(a_k\eta^{-1}+b_kG\eta^{-1})/m_k$, $c_k=(a_k+ib_k)/m_k$, $F(p)=\sum_kc_kw_k$ and
$R_{\rm eff}e^{i\varphi_{\rm eff}}=F(p(0))$. The \emph{kinetic condition} is that of \cref{prop:gauge1}, used in \cref{thm:survive}
with the kernels compressed by the projector: $X^{\!\top}K^\mu+K^\mu X=0$ for $\mu=0,\dots,3$, with the kernels $K^\mu=-\gamma^0\gamma^\mu$ of \cref{prop:kernels1}. At $\mu=0$ it says that $X$ is antisymmetric, and at $\mu=a$ that $X$ commutes with the boosts $\gamma^0\gamma^a$.

% --- source: emergence/e5_closure_two#2 (The closure symmetry and the invariants of a pair)
\subsection{The closure symmetry and the invariants of a pair}
\label{e5:sec:symmetry}

\begin{proposition}[The closure symmetry at two constituents]
\label{e5:prop:so6C}
$G_2^{\!\top}\eta_2G_2=-\eta_2$ and $G_2^2=-\id$. The centraliser of $G_2$ in $\so{6,6}$ has dimension $30$ and is spanned by the $15$ paired
rotations and the $15$ paired boosts. In the basis of the six temporal directions followed by their partners every element is
$X=\bigl(\begin{smallmatrix}R&S\\-S&R\end{smallmatrix}\bigr)$ with $R,S$ antisymmetric, and $X$ acts on $Z$ as the complex antisymmetric matrix
$R-iS$. The map $X\mapsto R-iS$ is a Lie isomorphism onto $\so{6,\C}$, the algebra of the complex orthogonal group of $q_{\rm tot}$, and the
compact part of the centraliser is $\so6$, the paired rotations.
\end{proposition}

This restates \cref{d5:prop:so6C}(i) (fixed there), the $n=2$ member of \cref{thm:so9C}; the bracket identity was checked
on all $435$ pairs of basis elements (e5-1a). By \cref{d5:prop:so6C}(ii), $\so{6,\C}=\lG^{(1)}\oplus\lG^{(2)}\oplus(\C^3\otimes\C^3)$,
and the cross part maps $z_1$ into $z_2$: the closure symmetry does not keep the two field strengths apart.

\begin{proposition}[Quadratic invariants of a pair]
\label{e5:prop:quad}
The holomorphic quadratic polynomials in $(z_1,z_2)$ invariant under
(a) $\so{6,\C}$ are the multiples of $q_{\rm tot}$;
(b) $\lG^{(1)}\oplus\lG^{(2)}$ form $\operatorname{span}_\C\{q_1,q_2\}$;
(c) $\lG^\Delta$ form $\operatorname{span}_\C\{q_1,q_2,z_1\cdot z_2\}$.
In each case the invariant real quadratic forms on $V_2$ have real dimension $2$, $4$, $6$ and are spanned by the real and imaginary parts of these.
\end{proposition}

\begin{proof}
$\lG\cong\so{3,\C}$ acts on $\C^3$ by the vector representation. By the first fundamental theorem for the orthogonal group \cite{Weyl39},
the polynomial invariants of $SO(3,\C)$ acting diagonally on $(z_1,z_2)$ are generated by the inner products $z_k\cdot z_l$ (and
$\det(z_1,z_2,\cdot)$, which needs three vectors), which gives (c). Under $\lG^{(1)}\oplus\lG^{(2)}$ a cross term $z_1\cdot Az_2$ is invariant
under $z_1\mapsto\Lambda z_1$ with $z_2$ fixed only if $A=0$, which gives (b). $\C^6$ is irreducible under $\so{6,\C}$, with $q_{\rm tot}$ its
only invariant form, which gives (a). A real invariant form that is not the real part of a holomorphic one would be Hermitian, $h(Xu,v)+h(u,Xv)=0$,
and applied to $X$ and $iX$ this forces $h(Xu,v)=0$ for all $X$, so $h=0$.
\end{proof}

Inherited (classical invariant theory), with the real dimensions computed as exact null spaces (e5-1b).

\begin{theorem}[The closure symmetry and the pair invariant]
\label{e5:thm:pair}
\begin{enumerate}[label=(\roman*),leftmargin=2.2em]
\item $q(z_1+z_2)=q_1+q_2+2\,z_1\cdot z_2$, $\det Q=q_1q_2-(z_1\cdot z_2)^2$, and
\[
M_{\rm pair}^2=\tfrac14\bigl(|q_1|^2+|q_2|^2+2|z_1\cdot z_2|^2\bigr)+\tfrac12\bigl|q_1q_2-(z_1\cdot z_2)^2\bigr| .
\]
So $M_{\rm pair}$, $z_1\cdot z_2$ and $q(z_1+z_2)$ are Lorentz invariants of the pair under $\lG^\Delta$. None of them is invariant under
$\lG^{(1)}\oplus\lG^{(2)}$: a boost of constituent~$2$ alone fixes $q_1,q_2$ and changes $z_1\cdot z_2$ and $M_{\rm pair}$. Under the closure
symmetry only $q_{\rm tot}=\tr Q$ is invariant.
\item $M_{\rm pair}$ is invariant under $\lG^\Delta$ and the copy rotation $C_v$, and the subalgebra of $\so{6,\C}$ that leaves it invariant
is $\lG^\Delta\oplus\R C_v$, of real dimension $7$.
\item The closure symmetry does not preserve $M_{\rm pair}$. The copy boost has the eigenvector $(e_1,ie_1)$ with eigenvalue $1$, and maps the
balanced pair $(e_1,ie_1)$, with $q_{\rm tot}=0$ and $M_{\rm pair}=1$, to $e^t(e_1,ie_1)$, with $M_{\rm pair}=e^{2t}$. The compact rotation
$R^{(1)}_\varphi\in SO(6)$ of \cref{d6:ex:R1} takes $((1,i,0),0)$ from $M_{\rm pair}=0$ to $M_{\rm pair}=|\sin\varphi|$.
\item For each $c\in\C$ the set $\{Z\neq0:\ q_{\rm tot}(Z)=c\}$ is a single orbit of $SO(6,\C)$, and on it $M_{\rm pair}$ takes every value in
$[\half|c|,\infty)$. The bound $M_{\rm pair}\ge\half|q_{\rm tot}|$ of \cref{d6:thm:pair}(iii) is attained at $(z,0)$, and the copy boost
carries $(z,0)$ to $(\cosh t\,z,\,i\sinh t\,z)$, with $q_{\rm tot}=q(z)$ and $M_{\rm pair}=\half|q(z)|\cosh2t$.
\end{enumerate}
\end{theorem}

\begin{proof}
(i) The first two identities are expansions, and with them the formula is $M_{\rm pair}^2=\tfrac14\|Q\|_F^2+\half|\det Q|$, which follows from
\cref{d6:thm:pair}(ii). A diagonal Lorentz transformation $\Lambda$ maps $Q=Z^{\!\top}Z$ to $Z^{\!\top}\Lambda^{\!\top}\Lambda Z=Q$; one acting on
constituent~$2$ alone maps $z_1\cdot z_2$ to $z_1\cdot\Lambda z_2$. The last statement is \cref{e5:prop:quad}(a). (ii) $C_v$ maps $Q$ to $RQR^{\!\top}$ with
$R$ real orthogonal, which keeps the singular values (\cref{d6:prop:gauge}(ii)). That nothing else in $\so{6,\C}$ keeps $M_{\rm pair}$ is
computed exactly: at $32$ pairs with Gaussian-integer entries and rational $|\det Q|$ the derivatives of $M_{\rm pair}^2$ along the $30$ real basis fields of $\so{6,\C}$ are rational numbers and form a matrix of exact rank $23$, so the stabiliser has dimension at most $7$; the seven directions of $\lG^\Delta\oplus\R C_v$ are independent and leave $M_{\rm pair}$ invariant identically (they act on $Q$ trivially or by $Q\mapsto RQR^{\!\top}$), so it is exactly $\lG^\Delta\oplus\R C_v$ (z1-5b). (iii) The copy
boost acts on the copy index as $\bigl(\begin{smallmatrix}0&-i\\i&0\end{smallmatrix}\bigr)$ (\cref{d5:prop:so6C}(iii)), and $M_{\rm pair}$
is homogeneous of degree two. At $(e_1,ie_1)$, $Q=vv^{\!\top}$ with $v=(1,i)$, so $\sigma_1=2$ and $\sigma_2=0$. The numbers for $R^{(1)}_\varphi$
restate \cref{d6:ex:R1}. (iv) Transitivity on the level sets of a nondegenerate quadratic form is Witt's theorem \cite{Witt37}. At $(z,0)$,
$Q=\operatorname{diag}(q(z),0)$. Since $(\varepsilon\otimes G)^2=\id$, $e^{t\,\varepsilon\otimes G}=\cosh t+\sinh t\,\varepsilon\otimes G$, and the image has
$Q=q(z)\,vv^{\!\top}$ with $v=(\cosh t,i\sinh t)$, so $\sigma_1=|q(z)|\,|v|^2=|q(z)|\cosh2t$ and $\sigma_2=0$. The orbit is connected and $M_{\rm pair}$
is continuous, so it takes every value between its infimum and its unbounded values; for $c=0$ use (iii).
\end{proof}

Items (i) and (iv) are fixed, with (i) partly restated from \cref{d6:thm:pair}; (ii) is fixed, with the stabiliser certified by an exact rational rank;
(iii) is fixed, with the second example restated. Checks: the identities symbolically, the formula against the singular values at $30$ random
pairs to $10^{-12}$ (e5-1e); invariance to $10^{-11}$, and the rank $23$ with a singular-value gap above $10^8$ (e5-1c), now replaced by the exact rank certificate z1-5b; the orbits to
$10^{-12}$ (e5-1d). So the closure symmetry mixes the two field strengths and does not preserve the pair invariant. The closure symmetry sees one field strength in $\C^6$ and its square $q_{\rm tot}$. To split it into two
field strengths that carry a pair invariant takes the seven-dimensional stabiliser $\lG^\Delta\oplus\R C_v$, a diagonal Lorentz action together
with the copy rotation, and the closure selects neither. With equal inertias its dynamics agrees: the frame energy is $\half\Re(c\,q_{\rm tot}(\pi))$
(\cref{prop:several}; \cref{e5:prop:rotorlaw}(i) with $c_1=c_2=c$), and its compact symmetry carries every balanced pair to one null constituent and an empty copy (\cref{d6:prop:gauge}(iii)). In that closure ``two constituents'' is a label that the dynamics does not see (interpretive).

% --- source: emergence/e5_closure_two#3 (No common Lorentz algebra)
\subsection{No common Lorentz algebra}
\label{e5:sec:nocommon}

\begin{proposition}[No common Lorentz algebra at two constituents]
\label{e5:prop:nocommon}
\begin{enumerate}[label=(\roman*),leftmargin=2.2em]
\item $\Cent(G_2)\cap\lint^{(1)}=0$. The elements of $\so{6,6}$ commuting with both $G_2$ and $\lint^{(1)}$ form $\lG^{(2)}\cong\so{3,\C}$, of
dimension $6$. At rung $n$ they form the $\so{3n-3,\C}$ of constituents $2,\dots,n$, of dimension $3(n-1)(3n-4)$: $0$, $6$ and $30$ at
$n=1,2,3$.
\item The bivectors of constituent~$1$ are $e_ae_b\otimes\id$, so $S_n\otimes\C$ carries $8^{n-1}$ copies of the content of $S\otimes\C$. Under
$\lG^{(1)}$, $S_2\otimes\C=16\,(\half,\half)$: no Weyl spinor ($128\,(\half,\half)$ at $n=3$). Under $\lint^{(1)}$,
$S_2\otimes\C=16\,(\half,0)\oplus16\,(0,\half)$.
\item Under $\lG^{(2)}$, $S_2\otimes\C=16\,(\half,\half)$, and under $\lG^\Delta$,
$S_2\otimes\C=4\bigl[(0,0)\oplus(1,0)\oplus(0,1)\oplus(1,1)\bigr]$: integer spin only.
\end{enumerate}
\end{proposition}

\begin{proof}
(i) An element of $\lint^{(1)}$ lives on constituent~$1$, where $\Cent(G)\cap\so{3,3}=\lG\cong\so{3,\C}$ (\cref{thm:GLorentz}). It fixes the
surplus pair, and an element of $\so{3,\C}$ that kills $z=(0,1,0)$ and $z=(0,0,1)$ is an antisymmetric matrix with two zero columns, hence zero. Under
$\lint^{(1)}$, $V_2$ is the interior four-plane of constituent~$1$, irreducible and nontrivial, plus eight Lorentz-trivial directions: the surplus
pair $F=\langle e^{(1)}_2,e^{(1)}_3\rangle$ and all of constituent~$2$. By Schur's lemma an equivariant $X\in\so{6,6}$ vanishes on the four-plane
and preserves the eight trivial directions. If $X$ also commutes with $G_2$, then for $v\in F$, $G_2v$ lies in the interior plane
$\langle e^{(1)}_5,e^{(1)}_6\rangle$, so $G_2Xv=XG_2v=0$ and $Xv=0$. Hence $X$ vanishes on constituent~$1$ and is an element of $\so{3,3}$ of
constituent~$2$ commuting with $G$, that is of $\lG^{(2)}$ (\cref{thm:GLorentz}). The same argument at rung $n$ leaves
$\Cent(G_{n-1})$ of constituents $2,\dots,n$, which is $\so{3n-3,\C}$ by \cref{d5:prop:so6C}(i). (ii) $E_{(1,a)}E_{(1,b)}=e_ae_b\otimes\id$,
and at one constituent $S\otimes\C=2\,(\half,\half)$ under $\lG$ and $2\,(\half,0)\oplus2\,(0,\half)$ under $\lint$ (\cref{sec:spin}).
(iii) The bivectors of constituent~$2$ are $\id\otimes e_ae_b$, so $\lG^\Delta$ acts on $(S\otimes\C)\otimes_\C(S\otimes\C)$, and
$(\half,\half)\otimes(\half,\half)=(0,0)\oplus(1,0)\oplus(0,1)\oplus(1,1)$.
\end{proof}

Part (i) restates \cref{d5:prop:so6C}(iv) and \cref{d7:prop:exist}(iii) (fixed there), the dimension $30$ restates \cref{thm:nocommon}, and the ladder formula is fixed (e5-2a). Parts (ii) and (iii) are fixed at $n=2$, restated at $n=3$, and the Clebsch--Gordan rule is
inherited (Casimir spectra to $10^{-9}$, computed on $S_2$ and on one constituent; e5-2b, e5-2c). No nonzero Lorentz generator leaves both the
closure and the spinor structure of matter invariant, at $n=2$ as at $n=1$ and $n=3$. What differs is what commutes with both: at $n=2$ it is a
whole second Lorentz algebra, the frame algebra of the spectator constituent. That is the material of the mixed rung, which pairs $\lint^{(1)}$ with
$\lG^{(2)}$ (\cref{e5:prop:mixed}).

% --- source: emergence/e5_closure_two#4 (Gauge algebras at two constituents and the binding root)
\subsection{Gauge algebras at two constituents and the binding root}
\label{e5:sec:gauge}

\begin{proposition}[No Pati--Salam algebra at two constituents]
\label{e5:prop:noPS}
The compact parts of the candidate gauge algebras $\mathfrak g(k_I,k_B,k_S)$ of the five Lorentz actions at $n=2$ (\cref{thm:algebra}; \cref{d5:tab:actions}) have dimensions $13$ (spectator, $\so5\oplus\so3$), $6$ (binding and spectator, $\so3\oplus\so3$), $7$ (diagonal interior,
$\so2\oplus\so4$), $1$ (mixed) and $1$ (binding on both). None contains the Pati--Salam algebra $\so4\oplus\so6$ \cite{PatiSalam}, of dimension $21$,
or its colour factor $\su4$, of dimension $15$.
\end{proposition}

Obstructed, by a dimension count; the five dimensions are re-verified (e5-3a, e5-3e, e5-3h, e5-3i). At $n=3$ the compact spectator algebra
$\so8\oplus\so6$, of dimension $43$, already contains $\mathfrak{ps}_U$ (\cref{sec:spectatorcompact}). The analogue at $n=2$ of \crefrange{thm:PSbinding}{thm:closurechiral} is therefore the intersection of each compact candidate with $\Cent(G_2)$.

\begin{proposition}[The spectator algebra and the binding root]
\label{e5:prop:spectator}
Let the Lorentz algebra act by $\lint^{(1)}$, with constituent~$2$ a spectator.
\begin{enumerate}[label=(\roman*),leftmargin=2.2em]
\item $\mathfrak g(1,0,1)=\Cent(\lint^{(1)})\cap\so{6,6}\cong\so{5,3}$ has dimension $28$. Its compact part is $\so5\oplus\so3$, with $\so5$ on the
positive Lorentz-trivial span $P=\langle e^{(1)}_2,e^{(1)}_3,e^{(2)}_1,e^{(2)}_2,e^{(2)}_3\rangle$ and $\so3$ on $N=\langle e^{(2)}_4,e^{(2)}_5,e^{(2)}_6\rangle$.
$\mathfrak g(1,0,1)\cap\Cent(G_2)=\lG^{(2)}$. Neither $\so5$ nor $\so3$ meets $\Cent(G_2)$, and with $E_2=\langle e^{(2)}_1,e^{(2)}_2,e^{(2)}_3\rangle$,
\[
(\so5\oplus\so3)\cap\Cent(G_2)=\bigl\{A+G_2AG_2^{-1}:\ A\in\so{E_2}\bigr\}\cong\su2 ,
\]
the three paired rotations of constituent~$2$: compact and non-abelian.
\item Among the spin lifts of the $30$ elements of $\Cent(G_2)$, those satisfying the kinetic condition with the kernels of constituent~$1$ are
exactly the paired rotations of constituent~$2$. The paired boosts of constituent~$2$ fail it.
\item Let $w$ be the product of the eight Lorentz-trivial generators. Then $w^2=-\id$, $w$ commutes with $\lint^{(1)}$ and with the spin lift of
$\so{5,3}$, and $\Gamma=w\,\Omega^{(1)}$, with $\Omega^{(1)}=E_{(1,1)}E_{(1,4)}E_{(1,5)}E_{(1,6)}$ the interior volume element of
constituent~1, is a chirality commuting with both. With $w$ as $i$, each chirality half is a Weyl block $(\half,0)\otimes(\mathbf4,\mathbf2)$, respectively $(0,\half)\otimes(\mathbf4,\mathbf2)$, of
$\lint^{(1)}\oplus\so5\oplus\so3$ (\cref{prop:spectator}). Under the binding-compatible $\su2$,
\[
L=(\half,0)\otimes\bigl(2\cdot\mathbf1\oplus2\cdot\mathbf3\bigr),\qquad R=(0,\half)\otimes\bigl(2\cdot\mathbf1\oplus2\cdot\mathbf3\bigr):
\]
integer isospin, and the same real representation on both halves, so the action is vector-like. With charges $h=-w\,S(X)$,
$\mathcal A(h)=\mathcal A(h^3)=0$, where $\mathcal A(f)=\tfrac14[\Tr(fP_L)-\Tr(fP_R)]$ as in \cref{thm:closurechiral}. There is no doublet, hence no
global $SU(2)$ anomaly \cite{Witten82}.
\end{enumerate}
\end{proposition}

\begin{proof}
(i) Let $A\in\so P$ and $B\in\so N$ with $[A+B,G_2]=0$, and follow the proof of \cref{thm:PSbinding}. $G_2$ maps $F=\langle e^{(1)}_2,e^{(1)}_3\rangle$
into the interior plane $\langle e^{(1)}_5,e^{(1)}_6\rangle$, on which $A$ and $B$ vanish, so $G_2Av=(A+B)G_2v=0$ and $Av=0$ for $v\in F$. For $v\in E_2$,
$Bv=0$ and $G_2v\in N$, so $G_2Av=BG_2v\in N$ and $Av\in G_2^{-1}N=E_2$. Hence $A\in\so{E_2}$ and $B=G_2AG_2^{-1}$ on $N$; conversely every such pair
commutes with $G_2$. Taking $B=0$ or $A=0$ shows that neither summand meets $\Cent(G_2)$. The rest is \cref{e5:prop:nocommon}(i). (ii) At $\mu=0$ the
condition is antisymmetry. The spin lift of a paired boost is a sum of boost bivectors, which are symmetric (\cref{thm:sig}). The spin lift
of a paired rotation of constituent~$2$ is antisymmetric and commutes with every bivector of constituent~$1$, in particular with $\gamma^0\gamma^a$.
A general element of $\Cent(G_2)$ was tested by an exact null space. (iii) $\so5\supset\so{E_2}\oplus\so F$ restricts the spinor $\mathbf4$ to two
isospin doublets, and $\so3$ acts on $\mathbf2$ as a doublet, so the diagonal $\su2$ sees $(\mathbf2\oplus\mathbf2)\otimes\mathbf2=2\cdot(\mathbf1\oplus\mathbf3)$.
Every representation of $\su2$ has weights symmetric under $m\mapsto-m$, so every odd trace of $h$ vanishes on each half
separately; integer isospin excludes doublets, and with them the global anomaly.
\end{proof}

Part (i) is fixed: it is the structure of \cref{thm:PSbinding} at $n=2$, with $U=P$ and $d=\dim E_2=3$. Part (ii) is fixed, and in
(iii) the Weyl blocks restate \cref{prop:spectator}, the branching is inherited, and the content and the anomaly traces are fixed. Checks:
exact Lie algebra (e5-3a), exact null space of the kinetic condition (e5-3d), Casimirs to $10^{-9}$ and integer traces (e5-3b, e5-3c); the odd traces vanish as well for the $\so3$ on $N$, which
acts by doublets (e5-3j). The binding
root therefore turns the isospin doublets of the spectator algebra, the ``doublets, no colour, real'' of \cref{sec:rungs}, into isospin
$0$ and $1$. No doublet survives at $n=2$, as for the copy isospin of \cref{d7:prop:isospin}. At $n=3$ the doublets survive in the
$(\half,\half)_{\pm1}$ of $\so4_\Delta$ on $P_*$ (\cref{thm:closurechiral}(ii)).

\begin{proposition}[The diagonal connection algebra and the binding root]
\label{e5:prop:diagonal}
Let the Lorentz algebra act by the diagonal $\lint$.
\begin{enumerate}[label=(\roman*),leftmargin=2.2em]
\item $\Cent(\text{diagonal }\lint)\cap\so{6,6}=\so2\oplus\so4$ has dimension $7$ and is compact. It meets $\Cent(G_2)$ in $\R C_v$, the copy
rotation of all six directions: the copy rotation of the interior four-planes plus that of the surplus pairs. This algebra is abelian and compact,
and $S(C_v)=-C_b$ with $C_b=\half\sum_a\eta_{aa}E_{(1,a)}E_{(2,a)}$.
\item There is no fermion under this action at $n=2$ (\cref{prop:orthospin}) and no copy-invariant gamma matrix (\cref{thm:copyparity}), so the kinetic condition has no canonical kernel. With the summed kernels $K^\mu_S=\sum_k\gamma^{(k)}_0\gamma^{(k)}_\mu$ it holds for
$C_b$. With the kernels $f_0f_\mu$ of the centre-of-mass generators $f_a=(E_{(1,a)}+E_{(2,a)})/\sqrt2$ (\cref{d5:prop:cmrel}) it fails.
\item $\Gamma_{12}^2=\id$, and $\Gamma_{12}$ commutes with the diagonal $\lint$ and with $C_b$. With copy charge $q$ defined by $C_b=iq$, the content
of $S_2\otimes\C$ is
\begin{center}
\small
\begin{tabular}{@{}ll@{}}
\toprule
$\Gamma_{12}=-1$ & $(0,0)_{\pm3}\oplus(0,0)_{\pm1}\oplus2\,(\half,\half)_{\pm1}\oplus(1,0)_{\pm1}\oplus(0,1)_{\pm1}$\\
$\Gamma_{12}=+1$ & $2\,(0,0)_{\pm2}\oplus(\half,\half)_{\pm2}\oplus2\,(\half,\half)_0\oplus2\,(1,0)_0\oplus2\,(0,1)_0$\\
\bottomrule
\end{tabular}
\end{center}
Spin is integral, $\Gamma_{12}=-1$ carries the odd and $\Gamma_{12}=+1$ the even copy charges, and each half is invariant under
$(q,(j_L,j_R))\mapsto(-q,(j_R,j_L))$: self-conjugate, hence vector-like. $\Tr(C_b^kP_\pm)=0$ for every odd $k$.
\end{enumerate}
\end{proposition}

\begin{proof}
(i) Under the diagonal $\lint$, $V_2=2V\oplus4\cdot\mathbf1$, and Schur's lemma gives $\so2\oplus\so4$, the $n=2$ case of \cref{thm:diagcent}. Within it, an element $aM+B$, with $M$ the copy rotation of the interior four-planes and $B$ in the surplus $\so4$, commutes with
$G_2$ exactly when $B$ is $a$ times the copy rotation of the surplus pairs, because $G_2$ maps each surplus pair into the interior plane of the same
constituent. This is \cref{d7:prop:exist}(ii), and the lift is \cref{d5:prop:cmrel}(ii). (ii) $C_b$ commutes with the diagonal boosts and is
antisymmetric, while $f_0f_\mu$ contains the cross terms $E_{(1,1)}E_{(2,\mu)}$ ($\gamma^0=e_1$), with which it does not commute. (iii) Computed.
\end{proof}

Part (i) restates \cref{d7:prop:exist}(ii) and \cref{d5:prop:cmrel} (fixed there; e5-3e), (ii) is fixed (exact; e5-3f), and (iii) is fixed
(eigenvalues to $10^{-9}$, traces exact; e5-3g). Under the diagonal action the binding root leaves an abelian, compact, Lorentz-invariant copy charge
on matter that has no fermions; the row ``$\so2\oplus\so4$, no fermions'' of \cref{sec:rungs} becomes ``$\so2$, no fermions''.

\begin{proposition}[The mixed rung and the binding root]
\label{e5:prop:mixed}
Let $\mathfrak l_{\rm mix}$ pair $\lint^{(1)}$ with $\lG^{(2)}$ generator by generator (\cref{prop:mixed}).
\begin{enumerate}[label=(\roman*),leftmargin=2.2em]
\item $\mathfrak l_{\rm mix}$ satisfies the $\so{1,3}$ brackets, its centraliser in $\so{6,6}$ is $\R J_1$, $J_1=e^{(1)}_2e^{(1)}_3$ acting on vectors
(\cref{prop:mixedcost}), and $\R J_1\cap\Cent(G_2)=0$. $G_2$ does not commute with $\mathfrak l_{\rm mix}$.
\item The $\mathfrak l_{\rm mix}$-equivariant endomorphisms of $V_2$ form a seven-dimensional algebra: the scalars on the four-vector, $\operatorname{span}\{\id,G\}$ on
the field strength of constituent~$2$, and $\mathfrak{gl}(2)$ on the two Lorentz-trivial directions. Its Euclidean-antisymmetric part, the generators
that rotate without dilating, is $\operatorname{span}\{G^{(2)},J_1\}$, with $G^{(2)}=0\oplus G$, and $[J_1,G^{(2)}]=0$.
\end{enumerate}
\end{proposition}

\begin{proof}
(i) $\R^{6,6}=V\oplus\Lambda^2V\oplus2\cdot\mathbf1$ under $\mathfrak l_{\rm mix}$ (\cref{prop:mixed}). $J_1$ rotates the surplus pair of
constituent~$1$, and $G_2$ maps that pair into the interior four-plane, which $J_1$ fixes, so $[J_1,G_2]\neq0$. No nonzero element of
$\lint^{(1)}$ commutes with $G_2$ (\cref{e5:prop:nocommon}(i)), so $G_2$ does not commute with $\mathfrak l_{\rm mix}$. (ii) Schur's lemma on the three isotypic pieces, with $\Hom_{\lG}(V,V)=\operatorname{span}\{\id,G\}$
(\cref{prop:lGhoms}).
\end{proof}

Part (i) restates \cref{prop:mixed} and~\cref{prop:mixedcost} and is fixed in what it adds; (ii) is fixed (exact; e5-3h). The binding root of the
closure is not covariant under the mixed action and leaves nothing of its algebra. The only covariant rotor generators that rotate without
dilating are the combinations of $G^{(2)}$ and $J_1$. With the rotor $G^{(2)}$, which binds constituent~$2$ alone, all of $\R J_1$ survives
(constructed: admitted, not selected). For completeness, under the two binding actions the binding root leaves $\lG^{(2)}$ of
$\mathfrak g(0,1,1)=\so{3,3}^{(2)}$, with compact part the $\su2$ of \cref{e5:prop:spectator}, and all of $\mathfrak g(0,2,0)=\so{2,\C}$, with compact part
$\R C_v$ (restating \cref{d5:prop:so6C}(iii) and \cref{d5:tab:actions}; e5-3i). Neither action carries a Dirac operator, so the kinetic condition does
not apply.

% --- source: emergence/e5_closure_two#5 (Inertia freedom at two constituents)
\subsection{Inertia freedom at two constituents}
\label{e5:sec:inertia}

\begin{theorem}[Inertia freedom at two constituents]
\label{e5:thm:inertia}
The real dimensions of the symmetric forms $N$ on $V_2^*$, invariant under a symmetry ($rN+Nr^{\!\top}=0$, \cref{def:inertiareq}), in the
Lorentz class $\{N,G_2\}=0$ and in the duality class $[N,G_2]=0$, are:
\begin{center}
\small
\begin{tabular}{@{}lcc@{}}
\toprule
symmetry & Lorentz class & duality class\\
\midrule
$\Cent(G_2)=\so{6,\C}$ & $2$ & $0$\\
$\lG^\Delta\oplus\R C_v$ (the stabiliser of $M_{\rm pair}$) & $2$ & $0$\\
$\lG^{(1)}\oplus\lG^{(2)}$ & $4$ & $0$\\
$\lG^\Delta$ & $6$ & $0$\\
$\lG^{(2)}$ & $14$ & $9$\\
$\su2$, the binding-compatible spectator part (\cref{e5:prop:spectator}) & $14$ & $10$\\
$\R C_v$, the binding-compatible diagonal part (\cref{e5:prop:diagonal}) & $18$ & $18$\\
none & $42$ & $36$\\
\bottomrule
\end{tabular}
\end{center}
Under the full closure symmetry $\so{3n,\C}$ the freedom is $(2,0)$ at $n=1,2,3$, spanned by $\eta_n^{-1}$ and $G_n\eta_n^{-1}$; with no symmetry it is
$(3n(3n+1),9n^2)$, that is $(12,9)$, $(42,36)$ and $(90,81)$.
\end{theorem}

\begin{proof}
Split $N$ into blocks between constituents. Under $\lG^{(k)}$ the invariant symmetric forms of one constituent are $\operatorname{span}\{\eta^{-1},G\eta^{-1}\}$,
of Lorentz class (\cref{prop:lGhoms} and \cref{thm:inertiareq}(ii),(iii)); an invariant cross block vanishes when only one constituent is acted on, and
is a complex multiple of the invariant bilinear form under $\lG^\Delta$. This gives $2\cdot2=4$ and $2\cdot3=6$, the real and imaginary parts of the
products $z_k\cdot z_l$ of \cref{e5:prop:quad}, and $\so{6,\C}$ keeps only $q_{\rm tot}$, which gives $2$. A copy rotation acts on the copy index, and the
copy-rotation-invariant complex symmetric $2\times2$ matrices are the multiples of $\id$, which gives $2$ for the stabiliser. Under $\lG^{(2)}$ constituent~$1$ is
free ($12$ and $9$ by \cref{thm:inertiareq}(i)), and its cross blocks vanish because $\so3$ moves every vector of $\R^3$. Under the compact $\su2$
alone the identity of constituent~$2$ is also invariant, which adds one duality form. Under $\R C_v$ alone, $N=\id\otimes A+\varepsilon\otimes B$ with $A$
symmetric and $B$ antisymmetric, which gives $12+6$ and $9+9$. With no symmetry the classes are the forms anticommuting and commuting with $G_n$.
\end{proof}

The first, third and fourth rows restate \cref{d6:sec:pairs} (``four under $\lG^{(1)}\oplus\lG^{(2)}$, six under $\lG^\Delta$ alone and two under
$\Cent(G_2)\cap\so{6,6}$''); the other rows, and the ladder, are fixed (exact ranks; e5-4a, e5-4b). Two observations follow. The seven-dimensional
stabiliser of the pair invariant fixes the inertia as much as the thirty-dimensional closure symmetry does: up to one complex scale. And, as at $n=3$, imposing
only what the binding root leaves of a gauge algebra enlarges the freedom to $14$ or $18$ real parameters in the Lorentz class, so the gauge parts do not
select the inertia.

% --- source: emergence/e5_closure_two#6 (Two constituents and one rotor)
\subsection{Two constituents and one rotor}
\label{e5:sec:rotor}

The closure on $\mathfrak s_2=\R\ltimes_{G_2}V_2$ with $H=h^2/(2I_0)+\sum_k\half p^{(k)\top}N_kp^{(k)}$ has $p^{(k)}=e^{-\theta G}p^{(k)}(0)$ for both
constituents, so $w_k=e^{2i\theta}w_k(0)$ and $I_0\ddot\theta=R_{\rm eff}\sin(2\theta+\varphi_{\rm eff})$. The rate is constant and nonzero iff
$R_{\rm eff}=0$ and $h(0)\neq0$, and then every frame moves on a circle of diameter $|N_kp^{(k)}(0)|/|\kappa|$. The balance set $\mathcal B=\{F=0\}$ has
codimension two away from $0$, its massive part is open and dense, and on it $|c_1|E_0(\pi^{(1)})=|c_2|E_0(\pi^{(2)})$ (\cref{d6:thm:balance}, restating
\cref{prop:several}). With equal inertias the pairs on the null subspaces $V_\pm=\{p^{(2)}=\mp Gp^{(1)}\}$ are balanced (\cref{d6:prop:pairs}). The following items are new.

\begin{proposition}[The rotor law with pair inertias]
\label{e5:prop:rotorlaw}
\begin{enumerate}[label=(\roman*),leftmargin=2.2em]
\item $\half p^{\!\top}Np=\Re\bigl((a+ib)w(p)\bigr)/(2m)$ for $N=(a\eta^{-1}+bG\eta^{-1})/m$, so block-diagonal Lorentz-class inertias give
$U(\theta)=\half\Re\bigl(e^{2i\theta}F(p(0))\bigr)$.
\item The Lorentz-class forms on $V_2^*$ invariant under $\lG^\Delta$ are exactly the real and imaginary parts of the three products $z_k\cdot z_l$ of the
momentum fields. For such an inertia $U=\half\Re\tr\bigl(C\,Q(\pi)\bigr)$ with $C$ complex symmetric, $Q(\pi(\tau))=e^{2i\theta}Q(\pi(0))$, and
\[
I_0\ddot\theta=R\sin(2\theta+\varphi),\qquad Re^{i\varphi}=\tr\bigl(C\,Q(\pi(0))\bigr).
\]
The balance condition is $\tr(CQ(\pi))=0$, and $C=\operatorname{diag}(c_1,c_2)$ recovers (i).
\end{enumerate}
\end{proposition}

\begin{proof}
(i) is \cref{thm:inertiareq}(ii) per constituent. (ii) The first statement is \cref{e5:thm:inertia} (row $\lG^\Delta$) with \cref{e5:prop:quad}(c). Every
$z(\pi^{(k)})$ turns as $e^{i\theta}$, so $Q\mapsto e^{2i\theta}Q$, $U(\theta)=\half R\cos(2\theta+\varphi)$, and $\dot h=-dU/d\theta$.
\end{proof}

Part (i) restates \cref{thm:inertiareq}(ii) (symbolic). Part (ii) restates the remark in \cref{d6:sec:pairs}, which states the form and the
balance condition; the rotor law with cross terms is fixed here. Checks: exact span, and the law along a flow with a random $C$ to $10^{-9}$ (e5-5a).

\begin{proposition}[The balance surface at two constituents]
\label{e5:prop:Bgeom}
Let $c_1,c_2\neq0$.
\begin{enumerate}[label=(\roman*),leftmargin=2.2em]
\item In the data $(w_1,w_2)\in\C^2$ the balance condition $c_1w_1+c_2w_2=0$ is a complex line, of real codimension two. In the momentum data
$\R^{12}$, $\mathcal B$ has codimension two away from $0$. If exactly one coefficient vanishes, say $c_2=0$, then $\mathcal B=\{w_1=0\}\times\R^6$
still has codimension two, but forces constituent~$1$ to be null; if both vanish, $\mathcal B$ is everything.
\item $\mathcal B\setminus0$ is connected. For $c_1=c_2$ it is the nonzero null cone $\{Z\cdot Z=0,\ Z\neq0\}$ of $Z=(z(\pi^{(1)}),z(\pi^{(2)}))\in\C^6$,
a single orbit of $SO(6,\C)$, diffeomorphic to $\R_{>0}\times V_2(\R^6)$ with $V_2(\R^6)$ the Stiefel manifold of orthonormal two-frames. For $c_1\neq c_2$
it is the image of that cone under $(z_1,z_2)\mapsto(z_1,\rho^{-1}z_2)$, $\rho^2=c_2/c_1$.
\item The all-null part of $\mathcal B$ has dimension $8$. The massive part, on which both constituents are massive, has dimension $10$, and $w_1$ maps it
submersively onto $\C^*$, with fibre over $s$ the product of the two complex spheres $\{w(p^{(1)})=s\}$ and $\{w(p^{(2)})=-c_1s/c_2\}$.
\end{enumerate}
\end{proposition}

\begin{proof}
(i) is \cref{d6:thm:balance}(iii); the degenerate cases are direct. (ii) $c_1q_1+c_2q_2=c_1(z_1,\rho z_2)\cdot(z_1,\rho z_2)$. A nonzero null vector of
$\C^6$ is $u+iv$ with $u,v\in\R^6$, $|u|=|v|>0$ and $u\perp v$; $SO(6,\C)$ is transitive on these by Witt's theorem \cite{Witt37}, and $V_2(\R^6)$ is
connected. (iii) The null set $\{w=0\}\subset\R^6$ has dimension four (\cref{d6:thm:balance}(iii)). On the massive part $w_1\neq0$ and $w_2=-c_1w_1/c_2\neq0$.
\end{proof}

Fixed, with (i) restated. Checks at $20$ random points each: $\operatorname{rank}dF=2$ for equal and unequal coefficients, orbit dimension $10$,
$\operatorname{rank}d(w_1,w_2)=4$ on the all-null part, and $\operatorname{rank}(dF,dw_1)=4$ on the massive part (e5-5f).

\begin{theorem}[Stability of the balance surface]
\label{e5:thm:stab}
Let $c_1,c_2\neq0$ and let $(\theta,h,p,x)$ solve the two-constituent closure with $\theta(0)=0$, $\kappa_0=h(0)/I_0$ and $\psi_0=\varphi_{\rm eff}$.
\begin{enumerate}[label=(\roman*),leftmargin=2.2em]
\item $R_{\rm eff}$ is a constant of motion for every datum. The balance surface and every level set $\{R_{\rm eff}=r\}$ are invariant, so the
surface is neither attracting nor repelling. A perturbation $\delta p$ of a balanced datum with nonzero momentum has imbalance $|dF(\delta p)|+O(|\delta p|^2)$, and
since $dF$ is onto there every small imbalance occurs (at $p=0$, $dF=0$).
\item Let $h(0)\neq0$. The rate is stable. With $\varepsilon=R_{\rm eff}/(I_0\kappa_0^2)$, $\kappa^2=\kappa_0^2[1+\varepsilon(\cos\psi_0-\cos\psi)]$ and
$|\kappa^2/\kappa_0^2-1|\le2\varepsilon$ for all $\tau$, and the rotor turns for all time when $\varepsilon<\half$.
\item Let $h(0)\neq0$. The bounded frame orbits are not stable. In the rotating regime, during each period of the rate frame $k$ is displaced by
$\lambda\,e^{-\varphi_{\rm eff}G}N_kp^{(k)}(0)$, with $\lambda$ of \cref{thm:freedrift}, and its mean velocity is
\[
\bar v_k=\tfrac\varepsilon4\bigl(1-\varepsilon\cos\psi_0\bigr)\,e^{-\varphi_{\rm eff}G}N_kp^{(k)}(0)+O(\varepsilon^3),
\]
of first order in the imbalance. A perturbation off the surface therefore moves the frames by a distance that grows linearly in $\tau$, at a rate linear
in the perturbation. A perturbation along the surface, or of $h(0)$, keeps them on circles.
\item If $h(0)=0$, a balanced datum is an equilibrium of the rotor, and each frame moves uniformly, $x_k=x_k(0)+\tau N_kp^{(k)}(0)$, without bound unless
$N_kp^{(k)}(0)=0$, in which case the frame is stationary.
Off the surface, unless $\psi_0\in\pi\Z$ (an equilibrium, at which the frames move uniformly), the rotor librates about $\psi=\pi$ with
$\Omega^2=2R_{\rm eff}/I_0$, amplitude $\xi_0\in(0,\pi)$ given by $\cos\xi_0=-\cos\psi_0$, and period $4\mathrm K(m)/\Omega$, $m=\sin^2(\xi_0/2)$,
which diverges as $R_{\rm eff}\to0$; $\mathrm K$ and $\mathrm E$ are the complete elliptic integrals in the parameter $m$ \cite{DLMF}. The
frames move with the mean velocity $\bigl(1-2\mathrm E(m)/\mathrm K(m)\bigr)e^{-\varphi_{\rm eff}G}N_kp^{(k)}(0)$. They drift, except at the one
amplitude with $2\mathrm E(m)=\mathrm K(m)$ ($m=0.8261$, $\xi_0=130.7^\circ$), at which every frame returns to its initial position after each
period.
\end{enumerate}
\end{theorem}

\begin{proof}
(i) $F(p(\tau))=e^{2i\theta}F(p(0))$ because every $w_k$ turns by $e^{2i\theta}$. (ii) Energy conservation, $H=\half I_0\kappa^2+\half R_{\rm eff}\cos\psi$, divided
by $\half I_0\kappa_0^2$; this is \cref{cor:multiscale} with $R_0/m$ replaced by $R_{\rm eff}$. (iii) Lorentz-class $N_k$ anticommute with $G$, so
$\dot x_k=e^{2\theta G}N_kp^{(k)}(0)$ with the common angle $\theta$, and the proof of \cref{thm:freedrift} applies to each frame. Expanding
$(1+\varepsilon(\cos\psi_0-\cos\psi))^{-1/2}$ to second order,
$\lambda=\frac{\pi\varepsilon}{4|\kappa_0|}\bigl(1-\tfrac32\varepsilon\cos\psi_0\bigr)+O(\varepsilon^3)$ and the period is
$\frac\pi{|\kappa_0|}\bigl(1-\half\varepsilon\cos\psi_0\bigr)+O(\varepsilon^2)$; their ratio gives $\bar v_k$. Along the surface $\varepsilon=0$, and a change
of $h(0)$ changes only the radius $|N_kp^{(k)}(0)|/(2|\kappa|)$ of the circles. (iv) At $h(0)=0$ and $R_{\rm eff}=0$, $\ddot\theta=0$ and $\theta\equiv0$, so
$\dot x_k=N_kp^{(k)}(0)$. Off the surface the rotor equation is the pendulum of \cref{cl:omega} with $R_0/m$ replaced by $R_{\rm eff}$. With
$\psi=\pi+\xi$, $\dot x_k=-e^{-\varphi_{\rm eff}G}(\cos\xi+G\sin\xi)N_kp^{(k)}(0)$. The substitution $\sin(\xi/2)=\sin(\xi_0/2)\sin t$ of that
proposition gives $d\tau=dt/(\Omega\sqrt{1-m\sin^2t})$, so over one period $\int\sin\xi\,d\tau=(2\sqrt m/\Omega)\int_0^{2\pi}\sin t\,dt=0$ and
$\int\cos\xi\,d\tau=(4/\Omega)\int_0^{\pi/2}(1-2m\sin^2t)(1-m\sin^2t)^{-1/2}dt=(4/\Omega)\bigl(2\mathrm E(m)-\mathrm K(m)\bigr)$; divide by the
period.
\end{proof}

Item (i) restates \cref{d6:thm:balance}(i); (ii) restates \cref{cor:multiscale} for two constituents; (iii) and (iv) are fixed. The
second-order term in (iii) is new. Item (iv) sharpens the closing sentence of \cref{prop:several}: on the balance surface the frames are
bounded only if $h(0)\neq0$, which is the qualifier ``constant and nonzero'' of \cref{d6:thm:balance}(i). Checks (DOP853, relative and absolute
tolerance $10^{-12}$; e5-5b to e5-5e): with equal inertias $a+ib=1+0.3i$, $m=I_0=1$, $h(0)=2$, and $p^{(2)}=-Gp^{(1)}$ drawn at random
($E_0(\pi^{(k)})=1.2156$ each), the rate stays constant to $10^{-11}$ over $\tau\in[0,60]$ and both frames lie on their circles (diameter $0.942$) to
$10^{-9}$. A generic balanced pair with unequal inertias ($m=1,2.5$; $a+ib=1+0.3i,\ 0.4-0.7i$) gives the same. Off the surface ($p^{(2)}$ perturbed by
$0.05$ along a fixed random direction, $R_{\rm eff}=0.391$, $\varepsilon=0.098$), $R_{\rm eff}$ and $w_1\overline{w_2}$ stay constant to $10^{-10}$, the
rotor obeys the pendulum law to $10^{-9}$, and the displacement per period is $\lambda e^{-\varphi_{\rm eff}G}N_kp^{(k)}(0)$ to relative $10^{-8}$. For
perturbations $10^{-3}$, $2\times10^{-3}$, $4\times10^{-3}$ the imbalance, the rate excursion and the mean frame velocity double with the perturbation (to
$2\%$). The mean velocity equals the quadrature value to $10^{-8}$, the first-order law to within $1.5\varepsilon$ and the second-order law to within
$\varepsilon^2$ (relative). With $h(0)=0$ the rotor librates with the period of (iv) to $10^{-8}$, and on the balance surface it stays at rest to
$10^{-12}$ while the frames move uniformly. Off the surface with $h(0)=0$ the displacement after one, two and three periods equals
the mean-velocity law of (iv) to relative $10^{-9}$ at four amplitudes, among them $\xi_0=130.7^\circ$, where it vanishes (e5-5i).

\begin{proposition}[Exchange of the constituents]
\label{e5:prop:exchange}
Let $\sigma(p^{(1)},p^{(2)})=(p^{(2)},p^{(1)})$.
\begin{enumerate}[label=(\roman*),leftmargin=2.2em]
\item $\sigma$ is an $\eta_2$-isometry commuting with $G_2$, so it preserves the Lie--Poisson bracket of $\mathfrak s_2$ and commutes with the rotor; it
preserves the equal-inertia Hamiltonian; as an element of $O(6,\C)$ it has determinant $-1$, so it lies outside the identity component of the closure
symmetry group. It exchanges $V_+$ and $V_-$.
\item With equal coefficients $R_{\rm eff}=2|c|\,|w_s|$ with $w_s=\half(w_1+w_2)$, the exchange-even part of $(w_1,w_2)$: a pair is balanced iff
$(w_1,w_2)$ is exchange-odd, $w_2=-w_1$.
\item The exchange eigenstates of the momenta, $p^{(2)}=\pm p^{(1)}$, both have $w_2=w_1$, so $F=(c_1+c_2)w_1$: unless $c_2=-c_1$ they are
balanced only when null. With equal inertias a massive balanced pair is never an exchange eigenstate of its momenta.
\end{enumerate}
\end{proposition}

\begin{proof}
(i) $\sigma$ permutes the copies; on $Z$ it is the permutation of three pairs of coordinates, of determinant $(-1)^3$. If $p^{(2)}=-Gp^{(1)}$, then
$\sigma p$ has second component $p^{(1)}=G(-Gp^{(1)})$. (ii) $R_{\rm eff}=|c|\,|w_1+w_2|$. (iii) $w$ is quadratic, so $F=c_1w_1+c_2w_1$.
\end{proof}

Fixed (e5-5g, e5-5j). Exchange parity of the duality invariants and exchange parity of the momenta are different things: the balanced pairs of $V_\pm$, a
magnetic and an electric constituent a quarter duality turn apart (\cref{d6:prop:pairs}(iii)), are exchanged with each other by $\sigma$ and are
not exchange eigenstates. With $c_1\neq c_2$, $\sigma$ is not a symmetry of the closure.

\begin{proposition}[What carries over from the rest closure]
\label{e5:prop:carry}
For one rotor on two constituents with Lorentz-class inertias:
\begin{enumerate}[label=(\roman*),leftmargin=2.2em]
\item (\cref{cl:frozen}.) A function of $p$ is a Casimir iff it is invariant under $p\mapsto e^{tG_2}p$. At $p\neq0$ the leaves are
two-dimensional, so there are $11$ local Casimirs, against $5$ at $n=1$, and the real and imaginary parts of $\bar z_az_b$, $a,b=1,\dots,6$, with $z\in\C^6$ the pair of
momentum fields, provide them. They include $E_6$, $E_0^2$ and $\mathbf S$ of each $\pi^{(k)}$ and the squared moduli $|Q_{kl}|^2$, which are smooth
everywhere; $E_0$ and $|Q_{kl}|$ are conserved invariant functions, smooth where they do not vanish (as for $E_0$ in \cref{cl:frozen}; a12-1). No Hamiltonian moves rest energy
between the constituents (\cref{d6:cor:nopair}).
\item (\cref{thm:pendulum}--4.8.) The pendulum, its small oscillations, the energy budget and the multiscale criterion hold with $R_0/m$ replaced
by $R_{\rm eff}$, which depends on the data only through $F(p(0))$. What changes is \cref{cl:flat}(i): a constant nonzero rate no longer
needs a null momentum field.
\item (\cref{cl:norel} and~\cref{cl:halfturn}.) $\Gamma_{12}^2=\id$, and $\Gamma_{12}$ anticommutes with the twelve generators, so it implements
$e^{\pi G_2}=-\id$ on $S_2$. Since $(e^{\theta G_2})^{\!\top}\eta_2e^{\theta G_2}=\cos2\theta\,\eta_2+\sin2\theta\,\eta_2G_2$, the rotor lifts to
$\mathrm{Pin}(6,6)$ only at $\theta\in\pi\Z$.
\item (\cref{sec:resttworot}.) The rotor turns both momentum fields at one rate, while the graded rest Hamiltonian of the tower has the single
frequency $(m_1^2+m_2^2)^{1/2}$ and the ungraded one the two frequencies $m_1$, $m_2$ (\cref{e1:prop:rest}(v)). A closure with two rates exists, with
the rotor acting by $\operatorname{diag}(\kappa_1,\kappa_2)\otimes G$, $\kappa_1\neq\kappa_2$ (\cref{d7:prop:rotor}), and is admitted, not selected.
\item (\cref{cl:ladder}.) Under \cref{cl:quant}, every leaf through the balance surface is flat, $U\equiv0$, so the
conditional ladder $\kappa_n=(n+\delta)\hbar/I_0$ holds exactly with both constituents massive, where at $n=1$ it needs null data.
\end{enumerate}
\end{proposition}

\begin{proof}
(i) The proof of \cref{cl:frozen} with $G_2$ in place of $G$; $zz^\dagger$ is a rank-one Hermitian matrix on $\C^6$, with $2\cdot6-1=11$ real
parameters. (ii) The equations for $(\theta,h)$ are those of one constituent with $R_0e^{i\varphi_0}/m$ replaced by $F(p(0))$. (iii) As in \cref{cl:halfturn}: a product of all twelve anticommuting generators anticommutes with each, and $G_2^{\!\top}\eta_2=\eta_2G_2$, $G_2^{\!\top}\eta_2G_2=-\eta_2$.
(iv) and (v) are direct.
\end{proof}

Items (i)--(iii) restate the cited results for two constituents (fixed there) and are fixed in their $n=2$ form (Jacobian rank $11$ and $5$, Lie--Poisson
rank $2$, exact Clifford identities; e5-5h). Item (iv) restates \cref{e1:prop:rest} and \cref{d7:prop:rotor} and is constructed for the two-rate
closure. Item (v) is constructed and conditional on the assumption. The copy rotor of \cref{d6:rem:copyrotor}, a second rotor that exists because the
uniqueness argument of \cref{thm:closureunique} fails under $\lG^\Delta$, and the two-rate rotor of (iv) are the structures of the rest closure that exist only from $n=2$ on.

% --- source: emergence/e5_closure_two#7 (The ladder: one, two and three constituents)
\subsection{The ladder: one, two and three constituents}
\label{e5:sec:ladder}

\begin{table}[htbp]
\centering
\footnotesize
\begin{tabular}{@{}>{\raggedright\arraybackslash}p{3.4cm}>{\raggedright\arraybackslash}p{2.9cm}>{\raggedright\arraybackslash}p{4.3cm}>{\raggedright\arraybackslash}p{3.9cm}@{}}
\toprule
 & $n=1$ & $n=2$ & $n=3$\\
\midrule
closure symmetry $\Cent(G_n)$ & $\lG\cong\so{3,\C}$, $6$ & $\so{6,\C}$, $30$ & $\so{9,\C}$, $72$\\
$\Cent(G_n)\cap\lint^{(1)}$ & $0$ & $0$ & $0$\\
commuting with $G_n$ and $\lint^{(1)}$ & $0$ & $\lG^{(2)}$, $6$ & $\so{6,\C}$ of $2,3$, $30$\\
$S_n\otimes\C$ under $\lG^{(1)}$ & $2\,(\half,\half)$ & $16\,(\half,\half)$ & $128\,(\half,\half)$\\
Pati--Salam algebra & no & no & in $\so8\oplus\so6$\\
binding-compatible compact part of the spectator algebra & $0$ & $\su2$, $3$; vector-like, isospin $0,1$ & $\so6$ of $2,3$, $15$; contains $\so4_\Delta$\\
the same for the diagonal connection algebra & $0$ & $\R C_v$, $1$; no fermions & copy rotations $\so3$, $3$\\
under the projectors & --- & --- & $\Pfam$: $\uu1^2$; $P_*$: $\so4_\Delta\oplus\R\,\mathrm{rot}(8,9)$, chiral\\
anomalies of what survives & --- & none & $\Pfam$: $\mathcal A(H^3)=-12\alpha\beta^2$; $P_*$: free on $\so4_\Delta$\\
inertia, closure symmetry & $(2,0)$ & $(2,0)$ & $(2,0)$\\
inertia, per-constituent $\lG$ & $(2,0)$ & $(4,0)$ & $(6,0)$\\
inertia, binding-compatible part & --- & $(14,10)$ spectator; $(18,18)$ diagonal & $(18,13)$ $\Pfam$; $(8,6)$ $P_*$\\
inertia, none & $(12,9)$ & $(42,36)$ & $(90,81)$\\
balance surface & $\{w=0\}$: null data only, dim $4$ & codim $2$, connected; massive part dense, dim $10$ & codim $2$, dim $16$\\
\bottomrule
\end{tabular}
\caption{The closure along the ladder. Dimensions are real; the inertia entries are (Lorentz class, duality class). The $n=3$ column restates \crefrange{thm:so9C}{thm:inertia3} and \cref{sec:spectatorcompact} and \cref{d6:thm:balance}(iii), except the spectator and diagonal rows, which
are computed here (e5-6a).}
\label{e5:tab:ladder}
\end{table}

\begin{table}[htbp]
\centering
\footnotesize
\begin{tabular}{@{}>{\raggedright\arraybackslash}p{8.0cm}>{\raggedright\arraybackslash}p{4.7cm}>{\raggedright\arraybackslash}p{2.1cm}@{}}
\toprule
statement & class & checks\\
\midrule
\cref{e5:prop:so6C}: $\Cent(G_2)=\so{6,\C}$, paired rotations and boosts, $R-iS$, compact part $\so6$ & restated (\cref{d5:prop:so6C}(i); \cref{thm:so9C} at $n=2$) & e5-1a\\
\cref{e5:prop:quad}: quadratic invariants $q_{\rm tot}$; $q_1,q_2$; $q_1,q_2,z_1\cdot z_2$ & inherited (invariant theory); dimensions fixed & e5-1b\\
\cref{e5:thm:pair}(i): $M_{\rm pair}$ through the Gram entries; diagonal against per-constituent invariance & fixed (partly restated, \cref{d6:thm:pair}) & e5-1e\\
\cref{e5:thm:pair}(ii): stabiliser of $M_{\rm pair}$ is $\lG^\Delta\oplus\R C_v$, dimension $7$ & fixed (exact rational rank $23$) & z1-5b, e5-1c\\
\cref{e5:thm:pair}(iii),(iv): $\so{6,\C}$ does not preserve $M_{\rm pair}$; range $[\half|c|,\infty)$ on each orbit & fixed (\cref{d6:ex:R1} restated; transitivity inherited) & e5-1d\\
\cref{e5:prop:nocommon}: no common Lorentz algebra; $\lG^{(2)}$; $16(\half,\half)$; integer spin under $\lG^\Delta$ & fixed (restated: \cref{d5:prop:so6C}(iv), \cref{d7:prop:exist}(iii); \cref{thm:nocommon}) & e5-2a to e5-2c\\
\cref{e5:prop:noPS}: no Pati--Salam algebra at $n=2$ & obstructed (dimension count) & e5-3a, 3e, 3h, 3i\\
\cref{e5:prop:spectator}: binding-compatible $\su2$; kinetic condition; integer isospin, vector-like, anomaly free & fixed (Weyl block restated, \cref{prop:spectator}; branching inherited) & e5-3a to e5-3d, 3j\\
\cref{e5:prop:diagonal}: $\R C_v$; no canonical kinetic kernel; content, vector-like & fixed ((i) restated, \cref{d7:prop:exist}(ii), \cref{d5:prop:cmrel}) & e5-3e to e5-3g\\
\cref{e5:prop:mixed}: mixed rung, nothing survives $G_2$; covariant rotors $G^{(2)},J_1$ & fixed (restated: \cref{prop:mixed}, \cref{prop:mixedcost}); rotor $G^{(2)}$ constructed & e5-3h\\
binding actions: $\lG^{(2)}$ and $\so{2,\C}$ survive & restated (\cref{d5:prop:so6C}(iii), \cref{d5:tab:actions}) & e5-3i\\
\cref{e5:thm:inertia}: inertia freedom at $n=2$; $(2,0)$ at $n=1,2,3$ & fixed (rows 1, 3, 4 restated, \cref{d6:sec:pairs}) & e5-4a, e5-4b\\
\cref{e5:prop:rotorlaw}: rotor law with pair inertias & fixed ((i) restated, \cref{thm:inertiareq}(ii)) & e5-5a\\
\cref{e5:prop:Bgeom}: codimension two, connected, one orbit, parts & fixed ((i) restated, \cref{d6:thm:balance}(iii)) & e5-5f\\
\cref{e5:thm:stab}: neutral surface, stable rate, secular frames, $h(0)\neq0$ needed, libration drift & fixed ((i), (ii) restated) & e5-5b to e5-5e, 5i\\
\cref{e5:prop:exchange}: exchange, $\det\sigma=-1$, balanced iff $w$ exchange-odd & fixed & e5-5g, 5j\\
\cref{e5:prop:carry}: Casimirs, pendulum, half turn, rates, ladder at $n=2$ & (i)--(iii) restated (\cref{sec:restclosure}), fixed at $n=2$; (iv), (v) constructed & e5-5h\\
\cref{e5:tab:ladder}: the ladder $n=1,2,3$ & fixed; $n=3$ mostly restated (\cref{sec:closurethree}) & e5-6a, 6b, 2a, 4b\\
\cref{e5:sec:ladder}: answer & interpretive (drawn from fixed statements) & ---\\
\bottomrule
\end{tabular}
\caption{Claims of \cref{e5:sec:main}, their classes and their checks, which are \texttt{PASS} lines of the script
\texttt{e5\_closure\_two.py}.}
\label{e5:tab:claims}
\end{table}

The new entries of \cref{e5:tab:ladder} at $n=1$ and $n=3$ are fixed (e5-6a). At $n=1$, $\Cent(\lint)\cap\so{3,3}=\R J$ meets $\Cent(G)$ in $0$, so the
binding root leaves no gauge algebra on one constituent. At $n=3$ the compact spectator algebra meets $\Cent(G_3)$ in the $15$ paired rotations of
constituents $2$ and $3$, which contain the $\so4_\Delta$ of \cref{thm:PSbinding}, and the diagonal connection algebra meets it in the three copy
rotations. At every rung the binding-compatible parts are $\so{3n-3}$, the paired rotations of constituents $2,\dots,n$, and $\so n$, the copy
rotations: the proofs of \cref{e5:prop:spectator}(i) and \cref{e5:prop:diagonal}(i) apply verbatim (fixed; exact at $n=4$, e5-6b). The balance condition $F=0$ has codimension two at every rung, because $dF$
has rank two at every nonzero point. At $n=1$ it is the null cone and carries no massive datum.

\paragraph{Answer.} \emph{Is $n=2$ a closure rung in its own right?} Yes, in a precise sense (interpretive, drawn from the fixed statements cited here):
every structure of the three-constituent analysis of \cref{sec:closurethree} except the two projectors, which exist only at $n=3$, has a computed $n=2$
analogue, and four of them first appear at $n=2$. The balance surface carries massive constituents at
an exactly constant rate (\cref{d6:thm:balance}). The closure symmetry mixes the constituents and does not see their pair invariant
(\cref{e5:thm:pair}). The closure symmetry compatible with the matter Lorentz algebra is nonzero, and is itself a Lorentz algebra (\cref{e5:prop:nocommon}).
And the binding root leaves a compact gauge algebra, non-abelian under the spectator action (\cref{e5:prop:spectator}). \emph{What $n=2$ adds to $n=1$} is
therefore a constant rate with mass, a gauge part that survives the binding root, a pair invariant and an exchange symmetry. \emph{What it lacks}, and
\emph{what $n=3$ adds}, is everything chiral. At $n=2$ every binding-compatible gauge part acts vector-like, with zero anomaly coefficients and integer
isospin, and no Pati--Salam algebra fits. At $n=3$ the projectors $\Pfam$ and $P_*$ exist, the surviving $\so4_\Delta$ acts on inequivalent halves
with doublets, the abelian charges of $\Pfam$ carry nonzero anomaly coefficients, and fermions exist under the diagonal action. The inertia question does
not distinguish the rungs: the closure symmetry fixes the inertia up to one complex scale at $n=1,2,3$ alike, and the gauge parts never do.

For the question of this section the closure gives a definite but modest answer (interpretive, drawn from the fixed statements above). The closure at $n=2$
distinguishes two constituents only through structures it does not select. Equal inertias and the closure symmetry see one field strength in $\C^6$. A
pair invariant needs a diagonal Lorentz action and the copy rotation. Per-constituent rest energies need the copies to be distinguished, and then nothing in
the closure moves rest energy between them (\cref{d6:prop:gauge} and \cref{d6:cor:nopair}). What two constituents can do that one cannot is to run at
an exactly constant rate with both momentum fields massive. That holds on a set of codimension two which is invariant but whose frames are secularly
unstable, which is closer to a balanced pair carried by the dynamics than to two particles produced by it.

% --- source: emergence/e6_bfield_triples#0 ((head))
% Track e6 -- The b-field on the temporal triple, c on the spatial triple, and the mass form T = e1e2e3.
% New references: sections/bib_e6.tex.

% --- source: emergence/e6_bfield_triples#1 (The b-field on the temporal triple, c on the spatial triple)
\section{\texorpdfstring{The $b$-field on the temporal triple, $c$ on the spatial triple}{The b-field on the temporal triple, c on the spatial triple}}
\label{e6:sec:main}

An earlier assignment places the arena of $c$ (of $c^2$) on the negative triple $\langle e_4,e_5,e_6\rangle$, of signature $(-,-,-)$, and the
arena of the $b$-field on the positive triple $\langle e_1,e_2,e_3\rangle$, of signature $(+,+,+)$, with the remark that the mass is unique in
$T=e_1e_2e_3$. In the conventions of this manuscript the positive triple is the temporal triple $V^+$ and carries $\bB$, the negative triple is
the spatial triple $V^-$ and carries $\bE$ ($X=\sum_i(B_ie_i+E_ie_{i+3})$, \cref{thm:emfield}), the clock is $\gamma^0=e_1$, the spatial
gamma matrices are $\gamma^a=e_{a+3}$, and $c$ is carried by the conversion $C\colon V^+\to V^-$, of which $C=c\,\id$ is the unique choice that
commutes with the diagonal rotations (\cref{thm:onec}; \cref{en:onec}). We therefore read the assignment as follows:
the $b$-field takes values in the temporal rotations $\so3_+=\operatorname{span}\{e_2e_3,e_3e_1,e_1e_2\}$, and $c^2$ is a symmetric part of $C$.
The trivector $T=\gamma^0J$ ($J=e_2e_3$) is the Dirac mass form, the unique antisymmetric kernel invariant under all fifteen bivectors
(\cref{cor:fullmass}), kept by $\mathrm{Spin}(3,3)$ at group level (\cref{mb:Tinv}).

The answers, with their classes, are these. (1) As a form $T$ is invariant under all of $\so{3,3}$; as an operator it commutes exactly with
$\so3_+\oplus\so3_-$, and its commutant is $M_4(\C)$ (fixed; \cref{e6:sec:Tform}). (2) The temporal rotations keep the mass form and the rest
flow, and the kinetic term breaks them to $\uu1_J$, in proportion to the speed of the clock. They and the Pauli--G\"ursey $\su2$ of the massless
kinetic term are two diagonals of one $\so4$ meeting in $\R J$ (fixed; \cref{e6:sec:dirac}). (3) A temporal $b$-field is half a field-strength
rotation and half a counter-rotation. It keeps the rates, splits the lines into $\kappa$ and $s\pm\omega/2$, $s=(\kappa^2+\omega^2/4)^{1/2}$, cannot be
removed along a worldline, and tears the binding (fixed; \cref{e6:sec:binding}). (4) The binding root $G_C$ preserves
$\tfrac12(c^2|\bB|^2+|\bE|^2)$ iff $C=cR$, $R\in O(3)$. A temporal rotation keeps the symmetric factor $CC^{\!\top}$ on the spatial triple and moves
only the orthogonal factor (fixed; the reading ``$c^2$ on $(-,-,-)$, $b$ on $(+,+,+)$'' is interpretive; \cref{e6:sec:polar}). (5) Under a temporal
$b$-field the four-dimensional mass read by a fixed clock oscillates at fixed six-dimensional mass, unless the clock or the temporal energy is on
the axis (fixed; reading interpretive; \cref{e6:sec:clock}). (6) Between two constituents the temporal copy rotation rotates the graded masses at
fixed $M=(m_1^2+m_2^2)^{1/2}$, half of it creates quanta in pairs, and in the ungraded reading no continuous transfer at fixed spectrum exists,
only the exchange of the two masses (fixed;
\cref{e6:sec:copies}). (7) At every rung the $\so{3n,3n}$-invariant forms are spanned by the temporal and the spatial volume; their symmetry has
period four in $n$, so the invariant Grassmann mass kernel is the temporal volume only for $n\equiv1\pmod4$ (fixed; \cref{e6:sec:ladder}).
\Cref{e6:sec:compare} compares the assignment with the $b$-field of \cref{sec:bfield,d7:sec:main}. Identifiers e6-$n$ refer to the \texttt{PASS} lines of
\texttt{e6\_bfield\_triples.py} (\cref{e6:tab:claims}).

\paragraph{Notation.} $S=\R^8$ with the realisation of equation~\cref{eq:realisation}; $\omega_T=T=e_1e_2e_3$, $\omega_S=e_4e_5e_6=\gamma^1\gamma^2\gamma^3$,
$\Gamma=\omega_T\omega_S=e_1\cdots e_6$, $J=e_2e_3$, $\Omega=e_1e_4e_5e_6$ (\cref{sec:dirac} and \cref{thm:subdichotomy}); $K^\mu=-\gamma^0\gamma^\mu$ are the kinetic kernels
(\cref{prop:kernels1}) and $\mathfrak g_{\rm massless}=\operatorname{span}\{\Omega,J,e_2\Omega,e_3\Omega\}=\uu1_\Omega\oplus\su2_{\rm PG}$ is the
massless kinetic algebra (\cref{prop:gauge1}). The temporal generators are symmetric and the spatial ones antisymmetric
(\cref{lem:blade}, \eqref{eq:transpose}). $\so3_\pm$ are the rotations of $V^\pm$, the six compact bivectors, and the nine bivectors $e_ie_{a+3}$ are the
boosts (\cref{thm:sig}). A bivector $Y$ acts on vectors by $v\mapsto Y(v):=\tfrac12[Y,v]$. On $V=V^+\oplus V^-$ we use the blocks $G=\bigl(\begin{smallmatrix}0&-\id\\ \id&0\end{smallmatrix}\bigr)$,
$G_C=\bigl(\begin{smallmatrix}0&-C^{-1}\\ C&0\end{smallmatrix}\bigr)$ (\cref{en:onec}), $N_C=\operatorname{diag}(C^{\!\top}C,\id)$ and
$E_C(X)=\tfrac12X^{\!\top}N_CX=\tfrac12(|C\bB|^2+|\bE|^2)$, with $E_c$ for $C=c\,\id$; $\lG\oplus\mathfrak p_G$ is the splitting of $\so{3,3}$ into the
parts commuting and anticommuting with $G$ (\cref{bf:lp}). The interior shift is $S=b+\kappa G$ with phase-space generator $M$
(equation~\cref{bf:eq:L} and \cref{bf:M}). Following \cref{d7:sec:main}, the \emph{rates} are the spectrum of the $\eta$-self-adjoint part of
$S$ and the \emph{lines} are the frequencies of $S$ itself. On the tower $E_{(k,a)}$, $N$ and the ungraded generators $u_a=\id\otimes e_a$ are as in
\cref{e1:sec:main}; $\hat C=\sigma(\varepsilon\otimes\id_6)=-\tfrac12\sum_a\eta_{aa}E_{(1,a)}E_{(2,a)}=-C_b$ is the spin lift of the copy rotation, as in \cref{e2:sec:main}; $C_b$ has the sign fixed in \cref{e1:sec:main}. Units are $\hbar=c=1$ unless $c$ is shown. 

% --- source: emergence/e6_bfield_triples#2 (T as a form and as an operator)
\subsection{\texorpdfstring{$T$ as a form and as an operator}{T as a form and as an operator}}
\label{e6:sec:Tform}

\begin{proposition}[The two invariances of the mass form]
\label{e6:prop:Tform}
\begin{enumerate}[label=(\roman*),leftmargin=2.2em]
\item \emph{As a form.} For every $Y\in\so{3,3}$, $Y^{\!\top}T+TY=0$; hence $Y^{\!\top}TY=-TY^2$, and $(e^{sY})^{\!\top}Te^{sY}=T$ for every real $s$. (The
ordering matters: $-Y^2T$ equals $-TY^2$ only when $[T,Y^2]=0$, which holds for a blade but fails for $Y=e_1e_2+e_3e_4$, whose square
$2e_1e_2e_3e_4-\id+\id$ anticommutes with $T$.) The same holds for
$\omega_S$. The forms on $S$ invariant under $\so{3,3}$ are $\operatorname{span}\{T,\omega_S\}$, with $T$ antisymmetric and $\omega_S$ symmetric.
\item \emph{As an operator.} $T$ commutes with the six compact bivectors and anticommutes with the nine boosts. Hence $e^{sY}$ commutes with $T$ for
$Y\in\so3_+\oplus\so3_-$, while $e^{sY}Te^{-sY}=e^{2sY}T$ for a boost.
\item \emph{The commutant.} The commutant of $T$ in $M_8(\R)$ is spanned by the $32$ blades with an even number of spatial indices. It is
$\Cl(V^+)\cdot\Cl^0(V^-)$, a product of two commuting subalgebras with $\Cl(V^+)\cong\Cl(3,0)\cong M_2(\C)$, whose centre is $\operatorname{span}\{\id,T\}$, and
$\Cl^0(V^-)\cong\HH$. So it is $M_2(\C)\otimes_\R\HH\cong M_4(\C)$, of dimension $32$ and with centre $\operatorname{span}\{\id,T\}$, and it meets $\so{3,3}$ in
$\so3_+\oplus\so3_-\cong\so4$. The $32$ blades with an odd number of spatial indices anticommute with $T$.
\item $\Ad_T$ acts on $V$ as $\eta$: on the frame it is parity.
\end{enumerate}
\end{proposition}

\begin{proof}
By the blade rule (\cref{lem:blade}) a blade $e_I$ commutes with $T$ iff $3|I|-|I\cap\{1,2,3\}|$ is even, that is iff $|I\cap\{4,5,6\}|$ is even,
and anticommutes otherwise. Conjugation by $T$ is diagonal in the blade basis, so its commutant is spanned by the blades it contains, as in
\cref{e1:prop:copies}(ii). This gives (ii) and the first sentence of (iii). A temporal generator commutes with a spatial bivector, so
$\Cl(V^+)$ and $\Cl^0(V^-)$ commute, and their products are the $8\times4$ commuting blades. $T$ is central in $\Cl(V^+)$; $e_4e_5$, $e_5e_6$ and their product
$-e_4e_6$ square to $-\id$. The isomorphisms $\Cl(3,0)\cong M_2(\C)$ and $M_2(\C)\otimes_\R\HH\cong M_4(\C)$ are inherited; equivalently, $T^2=-\id$ makes
$S\cong\C^4$, and the commutant of a complex structure is $\mathfrak{gl}(4,\C)$. For (i): on the fifteen basis bivectors, a compact blade is
antisymmetric and commutes with $T$, so $Y^{\!\top}T+TY=[T,Y]=0$, and a boost blade is symmetric and anticommutes with $T$, so
$Y^{\!\top}T+TY=\{T,Y\}=0$. The condition is linear in $Y$, so it holds on all of $\so{3,3}$. Then $Y^{\!\top}TY=-TYY$, and
$\tfrac{d}{ds}\bigl[(e^{sY})^{\!\top}Te^{sY}\bigr]=(e^{sY})^{\!\top}(Y^{\!\top}T+TY)e^{sY}=0$. $\omega_S$ has the same commutation pattern. The dimension two is
\cref{e6:thm:ladder}(i) at $n=1$. (iv) $T$ commutes with $e_1,e_2,e_3$ and anticommutes with $e_4,e_5,e_6$.
\end{proof}

Part (i) restates \cref{cor:fullmass} and \cref{mb:Tinv}(i) (fixed there), with the symmetric partner $\omega_S$ of
\cref{sec:selector}. Part (iv) restates \cref{en:parity}(1). Parts (ii) and (iii) are fixed (e6-1a to e6-1c); that $\omega_T$ and
$\omega_S$ are fixed by all six compact rotations is noted after \cref{thm:subdichotomy}. The two invariances are one statement read twice.
Because the compact bivectors are antisymmetric and the boosts symmetric, invariance of the form is commutation with the first and
anticommutation with the second. As a form, $T$ breaks nothing in $\so{3,3}$. As an operator it singles out the maximal compact subalgebra, the
rotations of the two triples, and among these it does not distinguish the temporal from the spatial ones. Two things distinguish $T$ from
$\omega_S$: as a Grassmann kernel $T$ is antisymmetric and $\omega_S$ symmetric, and as an operator $T^2=-\id$ is a complex structure while
$\omega_S^2=\id$ gives a projector (\cref{thm:subdichotomy}).

% --- source: emergence/e6_bfield_triples#3 (The temporal rotations on the Dirac data)
\subsection{The temporal rotations on the Dirac data}
\label{e6:sec:dirac}

The Dirac data of one constituent are the kinetic kernels $K^\mu$, the rest Hamiltonian $H(0)=m\gamma^0=me_1$ of the real first-order form,
and the rest flow $e^{-m\lambda T}$, with $JH(0)=mT$ (\cref{sec:firstorder}; \cref{mb:restflow} and~\cref{sec:resttworot}).

\begin{proposition}[What the temporal rotations do to the Dirac data]
\label{e6:prop:dirac}
Let $Y=\alpha J+\beta\,e_3e_1+\gamma\,e_1e_2\in\so3_+$.
\begin{enumerate}[label=(\roman*),leftmargin=2.2em]
\item \emph{Kinetic forms.} $Y^{\!\top}K^0+K^0Y=0$ and $Y^{\!\top}K^a+K^aY=2\,Y(\gamma^0)\,\gamma^a$, where $Y(\gamma^0)=\beta e_3-\gamma e_2$ is the velocity
of the clock. Hence $\sum_a\|Y^{\!\top}K^a+K^aY\|^2=96(\beta^2+\gamma^2)$ (Frobenius norm), and $Y$ maps the kinetic forms into their span,
$Y^{\!\top}K^\mu+K^\mu Y=\Lambda^\mu{}_\nu K^\nu$, iff $\beta=\gamma=0$: $\so3_+\cap(\lint\oplus\R J)=\R J$.
\item \emph{Rest Hamiltonian.} $e^{\varphi e_1e_2}e_1e^{-\varphi e_1e_2}=\cos2\varphi\,e_1-\sin2\varphi\,e_2$, and likewise with $e_1e_3$ and $e_3$. The stabiliser of
$H(0)$ in $\so3_+$ is $\R J$. A rotor of $\mathrm{Spin}(V^+)$ that carries $e_1$ to a clock $\hat n$ carries $J$ to $J(\hat n)=T\hat n$, $K^a$ to $-\hat n\gamma^a$ and
$H(0)$ to $mn$, and it fixes $K^0=-\id$ and $T$. Moreover $\hat nJ(\hat n)=J(\hat n)\hat n=T$.
\item \emph{Rest flow.} $e^{-m\lambda T}$ commutes with $\mathrm{Spin}(3)_+\times\mathrm{Spin}(3)_-$. The spin lift $\tfrac12e_je_k$ of $\so3_+$ has Casimir
$-\tfrac34\id$ on $S$, so $S$ carries isospin $\tfrac12$, and $e^{\pi e_je_k}=-\id$. The clock rotations $e_1e_2$ and $e_3e_1$ anticommute with $J$ and commute
with $T$. Their half turn $\Ad_{e_1e_2}=\Ad_{e_1}\Ad_{e_2}$ maps $(e_1,J,T,\gamma^a)$ to $(-e_1,-J,T,\gamma^a)$.
\item \emph{Field octants.} With the field-octant map $\Phi\colon V\to S$ of \cref{r3:thm:phi}, $-T\Phi=\Phi G$, so
$e^{-m\lambda T}\Phi=\Phi\,e^{m\lambda G}$: on the six field octants the Dirac rest flow is the duality rotation at angular frequency $m$.
\end{enumerate}
\end{proposition}

\begin{proof}
(i) $Y$ is antisymmetric and commutes with each $\gamma^a$ (a temporal bivector and a spatial generator), so
$Y^{\!\top}K^a+K^aY=[K^a,Y]=-[e_1,Y]\gamma^a=2Y(e_1)\gamma^a$, and $K^0=-\id$. With $v=Y(e_1)$ temporal, $(v\gamma^a)^{\!\top}(v\gamma^a)=|v|^2\id$, so each term has
squared norm $4\cdot8|v|^2$. The covariance statement intersects the seven-dimensional solution space $\lint\oplus\R J$ of \cref{mb:Tinv}(ii) with $\so3_+$; it is also solved directly (e6-2a). (ii) $e_1$ anticommutes with $e_1e_2$, so $e^{\varphi e_1e_2}e_1e^{-\varphi e_1e_2}=e^{2\varphi e_1e_2}e_1$. The
stabiliser condition is $Y(e_1)=0$. A rotor $R$ of $\mathrm{Spin}(V^+)$ commutes with $T$ and with every $\gamma^a$, and $J=Te_1$, so $\Ad_RJ=T\hat n$ and
$\Ad_RK^a=-\hat n\gamma^a$. Since $\hat n$ commutes with $T$ and $\hat n^2=\id$, $\hat nJ(\hat n)=J(\hat n)\hat n=T$. (iii) $T$ commutes with the compact bivectors by
\cref{e6:prop:Tform}(ii); the Casimir is $\tfrac14\sum Y_j^2$ with $Y_j^2=-\id$; the rest is the blade rule. (iv) Iterate $-T\Phi=\Phi G$.
\end{proof}

Part (i) is fixed (e6-2a); its covariance half restates \cref{mb:Tinv}(ii). Part (ii) is fixed (e6-2b, exact at $\cos\varphi=\tfrac35$
and at the clock $\hat n=-\tfrac7{25}e_1-\tfrac{24}{25}e_2$, where (i) holds with $Y(\hat n)$ in place of $Y(\gamma^0)$). Its transport of the data restates
\cref{r2:prop:J} and \cref{r2:thm:i}, by which the Dirac mass form $\gamma^0J(\hat n)=T$ is the same for every clock while $J(\hat n)$ turns with $\hat n$.
Part (iii) is fixed (e6-2c). Part (iv) restates \cref{r3:prop:odd}(i), re-verified with $\Phi$ built from the Fock vacuum (e6-2c).

So the temporal $SU(2)$ is a symmetry of the mass form and of the rest flow that the kinetic term breaks to $\uu1_J$. That reverses the usual
order, kinetic term symmetric and mass breaking, as \cref{sec:massdirac} found for $\mathrm{Spin}(3,3)$. The broken directions are the two clock
rotations: the kinetic term pairs one temporal direction with the spatial triple, and the breaking is the speed $|Y(\hat n)|$ of the clock. They flip the
charge, and their half turn acts as $\mathsf{CP}$ ($\Ad_{e_2}$ is the charge conjugation of \cref{thm:C}, $\Ad_{e_1}$ the parity), reversing
the sense of the rest rotation as in \cref{mb:senses}(iii). Interpretively, the temporal $SU(2)$ joins particle and antiparticle through
the sphere of clocks at fixed mass form.

\begin{proposition}[The mass symmetry and the kinetic symmetry]
\label{e6:prop:twosu2}
Let $P_\pm=\tfrac12(\id\pm\omega_S)$.
\begin{enumerate}[label=(\roman*),leftmargin=2.2em]
\item $\so3_+\cap\mathfrak g_{\rm massless}=\R J$.
\item $e_2\Omega=-(e_1e_2)\,\omega_S$ and $e_3\Omega=(e_3e_1)\,\omega_S$. $\omega_S$ commutes with $\so3_+$, and $P_\pm$ are orthogonal projectors of rank four.
\item $\so3_+=\{YP_++YP_-\}$ and $\su2_{\rm PG}=\{YP_++(JYJ^{-1})P_-\}$, $Y\in\so3_+$: both are graphs of automorphisms of $\so3$ over the two
eigenspaces of $\omega_S$, the identity and the half turn $\Ad_J$ about $J$. The Lie algebra they generate is $\so3_+P_+\oplus\so3_+P_-\cong\so4$, and they
meet in the fixed algebra of $\Ad_J$, which is $\R J$.
\item $g=P_++JP_-=\exp(\tfrac\pi2JP_-)$ is orthogonal and $g\,\so3_+\,g^{-1}=\su2_{\rm PG}$; $2g=\id+J+\omega_S-J\omega_S$ has components of grades $0,2,3,5$,
so $g\notin\mathrm{Pin}(3,3)$.
\item $\so3_+$ preserves $T$ and the kinetic term breaks it to $\R J$; $\su2_{\rm PG}$ preserves the kinetic forms and the mass breaks it to $\R J$.
\end{enumerate}
\end{proposition}

\begin{proof}
(ii) $e_2\Omega=e_2e_1e_4e_5e_6=-e_1e_2\,\omega_S$, and similarly for $e_3$. $\omega_S$ is disjoint from each temporal bivector, and the grades $3$ and $2$
give a commuting sign. (iii) Substitute $\omega_S=P_+-P_-$. $J$ commutes with $J$ and anticommutes with $e_1e_2$ and $e_3e_1$, so
$\Ad_J$ is the half turn about $J$. A bracket of two graphs of automorphisms $\phi_1,\phi_2$ contains $[Y,Y']P_++[\phi_1Y,\phi_2Y']P_-$, and the
generated algebra is computed (e6-2d). (iv) $(JP_-)^2=-P_-$ and $gYg^{-1}=YP_++JYJ^{-1}P_-$, because $P_\pm$ commute with $Y$ and $J$. (i) follows from (iii),
and (v) is \cref{e6:prop:dirac}(i) together with \cref{mb:massless}(i)--(ii).
\end{proof}

This is fixed (e6-2d); (v) restates \cref{prop:gauge1} and \cref{mb:massless}. The comparison has a sharp answer.
The two $\su2$'s are not the same, and neither contains the other. They are the untwisted and the $\Ad_J$-twisted diagonals of $\so3_+P_+\oplus\so3_+P_-$,
and they share the Cartan direction $J$. One is the symmetry of the mass that the kinetic term breaks, the other the symmetry of the kinetic term
that the mass breaks. Both leave the same $\uu1_J$ \cite{Pauli57,Gursey58}. The orthogonal map that exchanges them is not in the Clifford group.

% --- source: emergence/e6_bfield_triples#4 (A temporal bb-field and the binding root)
\subsection{A temporal \texorpdfstring{$b$}{b}-field and the binding root}
\label{e6:sec:binding}

\begin{proposition}[A $b$-field on the temporal triple]
\label{e6:prop:binding}
Let $b=(R,0)$ on $V^+\oplus V^-$, with $R=[\mathbf r]_\times\neq0$, $\omega=|\mathbf r|$, $\kappa>0$ and $S=b+\kappa G$.
\begin{enumerate}[label=(\roman*),leftmargin=2.2em]
\item $b=b_l+b_p$ with $b_l=\tfrac12(R,R)\in\lG$, a field-strength rotation, and $b_p=\tfrac12(R,-R)\in\mathfrak p_G$, a counter-rotation; they commute. In
the notation of equation~\cref{bf:eq:compact}, $\mathbf a=\mathbf c=\mathbf r/2$, so $\mathbf a\cdot\mathbf c=\omega^2/4$.
\item $e^{\theta b}Ge^{-\theta b}=G_{C(\theta)}$ with $C(\theta)=e^{-\theta R}$. $G_C^2=-\id$ for every $C$, and $G_C$ is $\eta$-self-adjoint iff $C$ is
orthogonal; so $\kappa G_{C(\theta)}$ stays in the self-adjoint summand and the rates, $\pm i\kappa$, are unchanged. $[b,G]=2b_pG\neq0$, so $b$ is not removable along a worldline.
\item $\det(\lambda-S)=(\lambda^2+\kappa^2)\bigl(\lambda^4+(2\kappa^2+\omega^2)\lambda^2+\kappa^4\bigr)$. The lines are $\kappa$, on the plane
$\langle\hat{\mathbf r},G\hat{\mathbf r}\rangle$, and $s\pm\omega/2$, with $s=(\kappa^2+\omega^2/4)^{1/2}$ and $(s+\omega/2)(s-\omega/2)=\kappa^2$; at $(\kappa,\omega)=(2,3)$
they are $1,2,4$.
\item Not every solution is bounded. At $(\kappa,\omega)=(2,3)$ and $(6,5)$, $M$ has Jordan blocks exactly at the frequencies $s\pm\omega/2$ and is
semisimple at $\kappa$.
\item On spinors $e_je_k=\tfrac12(e_je_k-e_{j+3}e_{k+3})+\tfrac12(e_je_k+e_{j+3}e_{k+3})$ for temporal $j,k$. The first summand, a paired rotation, commutes with
$N$; the second changes $N$ by $\pm2$ only.
\end{enumerate}
\end{proposition}

\begin{proof}
(i) $G$ commutes with $\operatorname{diag}(R,R)$ and anticommutes with $\operatorname{diag}(R,-R)$ in the $(\bB,\bE)$ blocks. (ii)
$e^{\theta b}=\operatorname{diag}(e^{\theta R},\id)$, and $\operatorname{diag}(Q,\id)\,G\,\operatorname{diag}(Q^{-1},\id)=G_{Q^{-1}}$. $G_C^\dagger=\eta G_C^{\!\top}\eta$ has blocks $-C^{\!\top}$ and
$C^{-\top}$, which equal those of $G_C$ iff $C^{\!\top}=C^{-1}$. The rates are the spectrum of the self-adjoint part, a tensor (\cref{d7:thm:schr}(iii)). Removability fails by \cref{bf:remove}. (iii) $b_l$ commutes with $b_p$ and with $G$. By Larmor's theorem (\cref{bf:larmor}) the substitution $X=e^{-b_l\tau}Y$ removes $b_l$ and shifts the exponents by the eigenvalues $0,\pm i\omega/2$ of $-b_l$. On the support of $b_p$ the
remaining shift $b_p+\kappa G$ has the frequency $(\kappa^2+\omega^2/4)^{1/2}$ of \cref{bf:pG}(ii), and on its kernel $\langle\hat{\mathbf r},G\hat{\mathbf r}\rangle$ the
frequency $\kappa$. The characteristic polynomial is also computed symbolically in $\mathbf r$. (iv) By \cref{bf:screen}(i), boundedness needs
$\mathbf a\cdot\mathbf c=0$. The Jordan structure is that of $b_p+\kappa G$ (\cref{bf:pG}(ii)), transported by the bounded frame $e^{-b_l\tau}$, and is
computed exactly. (v) With $e=a^++a^-$ and $e'=a^+-a^-$, the sum is $2(a^+_ja^+_k+a^-_ja^-_k)$ and the difference $2(a^+_ja^-_k+a^-_ja^+_k)$.
\end{proof}

This is fixed (e6-3a to e6-3e); the statement on the rates restates \cref{d7:thm:schr}(iii). A $b$-field on the temporal triple alone is
therefore half binding-compatible and half binding-odd. Its paired half is the compact part of $\lG$, a Zeeman term removable by a rotating frame
\cite{Larmor1897}. Its counter half lies in $\mathfrak p_G$. The latter gives the gauge-covariant remainder $D_\tau G=[b,G]=2b_pG$, it moves the binding root
to the root $G_{e^{-\theta R}}$ of a rotated conversion, and it tears the binding, with the screened frequency $s$ and Jordan blocks. The rate
$\kappa=c/\ell$ is never touched. In \cref{e6:prop:binding}(iii) the two halves act independently, because $\mathbf a\parallel\mathbf c$: the lines are
the screened frequency $\sqrt{\kappa^2+|\mathbf c|^2}$ shifted by the Zeeman amounts $\pm|\mathbf a|$. \cref{bf:screen} needed
$\mathbf a\perp\mathbf c$ for screening; a field on one triple has the opposite geometry. On spinors the same split is the split of
\cref{e1:prop:Nbiv}(ii): the paired half hops quanta within a constituent and the counter half creates or annihilates them in pairs, so the
temporal rotations, $J$ included, do not conserve $N$ (\cref{e1:prop:Nlorentz}(ii); \cref{e2:lem:N}).

% --- source: emergence/e6_bfield_triples#5 (c)
\subsection{\texorpdfstring{$c$}{c} on the spatial triple: the polar form of the conversion}
\label{e6:sec:polar}

\begin{proposition}[Conversions and the temporal rotations]
\label{e6:prop:polar}
Let $C\colon V^+\to V^-$ be invertible, $z_C=C\bB+i\bE$ and $q_C=z_C\cdot z_C$.
\begin{enumerate}[label=(\roman*),leftmargin=2.2em]
\item $G_C^2=-\id$, $G_C$ acts on $z_C$ as multiplication by $i$, and $e^{\theta G_C}$ preserves $E_C$, which depends on $C$ only through $C^{\!\top}C$.
\item $G_C^{\!\top}N_c+N_cG_C=\bigl(\begin{smallmatrix}0&C^{\!\top}-c^2C^{-1}\\ C-c^2C^{-\top}&0\end{smallmatrix}\bigr)$, so $G_C$ preserves
$E_c=\tfrac12(c^2|\bB|^2+|\bE|^2)$ iff $C^{\!\top}C=c^2\id$, that is iff $C=cR$ with $R\in O(3)$.
\item For $C=cR$, the elements of $\so3\oplus\so3$ commuting with $G_C$ are the twisted diagonal rotations $(Q,RQR^{-1})$. $G_{cR}=\Ad_{(R^{-1},\id)}G_c=\Ad_{(\id,R)}G_c$
for $R\in SO(3)$, and $\Ad_{(R_+,R_-)}G_c=G_{cR_-R_+^{-1}}$ for $(R_+,R_-)\in O(3)\times O(3)$. The Lie algebra preserving $\eta$ and $E_c$ is
$\so3\oplus\so3$, and adding $[Y,G_c]=0$ leaves the diagonal $\so3$. Hence $\{G_C:C^{\!\top}C=c^2\id\}$ is one orbit of $O(3)\times O(3)$, of dimension three.
\item A temporal rotation $(R_+,\id)$ carries $G_C$ to $G_{CR_+^{-1}}$ and $N_C$ to $N_{CR_+^{-1}}$. It conjugates $C^{\!\top}C$ to $R_+C^{\!\top}CR_+^{-1}$ and fixes
$CC^{\!\top}$; in the polar form $C=(CC^{\!\top})^{1/2}U$ it moves only $U$. A spatial rotation $(\id,R_-)$ carries $C$ to $R_-C$, fixes $C^{\!\top}C$ and
moves only $U'$ in $C=U'(C^{\!\top}C)^{1/2}$.
\item For $C=cR$: $E_c$ and $\Re q_C=c^2|\bB|^2-|\bE|^2$ do not depend on $R$, while $\Im q_C=2c\,(R\bB)\cdot\bE$ does. At $c=1$, $\bB=e_1$, $\bE=e_2$ and $R$ the
rotation through $\theta$ about $e_3$, $q_C=2i\sin\theta$: the same field point is null for $\theta=0$ and at rest for $\theta=\tfrac\pi2$, where
$E_0=E_6=1$.
\end{enumerate}
\end{proposition}

\begin{proof}
(i) $G_C^2$ has blocks $-C^{-1}C$ and $-CC^{-1}$. On $z_C$, $G_CX$ has $C\bB'=-\bE$ and $\bE'=C\bB$, so $z_C(G_CX)=iz_C(X)$, and $|z_C|^2=2E_C$. (ii) is a
block computation. Its vanishing means $C^{\!\top}C=c^2\id$, and then $C/c$ is orthogonal. (iii) The lower-left block of $[G_C,\operatorname{diag}(Q,Q')]$ is
$CQ-Q'C$. The conjugation formulas are block products. An $\eta$-isometry preserving the positive form $N_c$ preserves the eigenspaces $V^\pm$ of
$\eta^{-1}N_c=\operatorname{diag}(c^2,-1)$, and is orthogonal on each. The stabiliser dimensions are exact null spaces (e6-4c). (iv)
$(CR_+^{-1})(CR_+^{-1})^{\!\top}=CC^{\!\top}$. (v) Expand $z_C\cdot z_C$; in the example $R\bB=(\cos\theta,\sin\theta,0)$.
\end{proof}

This is fixed (e6-4a to e6-4e). Parts (i)--(iii) extend \cref{en:onec} and \cref{thm:onec} from $C=c\,\id$ to the orbit. They are
elementary linear algebra: Schur's lemma applied to a twisted diagonal action gives $C=cR$. Part (v) is the passive form of \cref{en:transport}(2), where the anti-diagonal part of a compact transport turns a null field into one at rest at fixed $E_6$.

\begin{remark}[What ``$c^2$ on $(-,-,-)$, $b$ on $(+,+,+)$'' fixes]
\label{e6:rem:csym}
\begin{enumerate}[label=(\alph*),leftmargin=2.2em]
\item \emph{Fixed.} The complete invariant of a conversion under the temporal rotations is $CC^{\!\top}$, a symmetric form on the spatial triple:
$C$ and $C'$ lie in one orbit of the orthogonal group of the temporal triple iff $C'C'^{\!\top}=CC^{\!\top}$ (\cref{e6:prop:polar}(iv); e6-4f). A $b$-field on the temporal triple,
a flow in $SO(3)$, can change anything about the conversion except this and the sign of $\det C$ (e6-4f). The paper's $c$ is the isotropic case $CC^{\!\top}=c^2\id$, which is also $C^{\!\top}C=c^2\id$, and that
is exactly the case in which the binding root preserves the weighted energy (\cref{e6:prop:polar}(ii)). In this sense, which is the reading of (c),
$c^2$ sits on the spatial triple and the temporal $b$-field cannot reach it. The energy $E_C$ sees the other symmetric factor $C^{\!\top}C$, a form on
the temporal triple that the temporal rotations conjugate (\cref{e6:prop:polar}(i), (iv)). After the two triples are identified, the two factors coincide
exactly when $C$ is normal, $CC^{\!\top}=C^{\!\top}C$, for instance $C=\operatorname{diag}(2,3,4)$: isotropy is sufficient, not necessary. What the temporal $b$-field moves is the orthogonal factor $R$. The energy $E_c$ does not see
$R$; the binding root $G_{cR}$ (which spatial direction is paired with which temporal one), the pseudoscalar $\Im q_C$ and the rest energy
$\tfrac12|q_C|$ do.
\item \emph{Fixed: a mirror.} With the triples exchanged the statement still holds: a spatial $b$-field keeps $C^{\!\top}C$ on the temporal triple and
moves the other orthogonal factor. The conversion alone therefore does not decide which triple carries $c^2$ and which the $b$-field. The Dirac data
decide it, through the clock in $V^+$ and the kinetic term (\cref{e6:prop:dirac}): in the interior reading $\so3_-\subset\lint$ is part of the
Lorentz algebra, while $\so3_+$ meets $\lint\oplus\R J$ only in $J$. The frame reading does not decide it, since its Lorentz algebra $\lG$ contains
only the diagonal rotations.
\item \emph{Relation to the clause ``$c^2$ is the symmetric part of $E=mc^2$''.} \Cref{sec:energy-c} gives the clause its exact forms: $c$ is
the unique isotropic conversion, its weight in $E_6$ a unit convention (\cref{en:onec,en:SI}), and the shift splits as
$\End(V)=\so{3,3}\oplus\{Y^\dagger=Y\}$ with $\kappa G$ in the self-adjoint summand, $(\kappa G)^2=-c^2/\ell^2$, and the connections in $\so{3,3}$ (\cref{sec:bfield}). The polar decomposition of the conversion is a further, compatible form: its symmetric factor carries $c^2$ and is invariant under the
temporal $b$-field, its orthogonal factor is moved by it, and for every orthogonal $C$, $\kappa G_C$ stays in the self-adjoint summand with the same
spectrum (\cref{e6:prop:binding}(ii)). It adds nothing else: it does not fix $c$, does not make $c$ a symmetry (the stabiliser $O(3)\times O(3)$ of $\eta$
and $E_c$ is not a Lorentz group), and selects no rest energy. A constant $R$ is a choice of identification of the triples (\cref{thm:onec} selects $R=\id$ once a common $SO(3)$ is fixed); only a $\tau$-dependent $R(\tau)$, a temporal $b$-field, has gauge-invariant content, its
$D_\tau G\neq0$. Reading the symmetric factor as ``$c^2$'' and the orthogonal factor as the place of the $b$-field is interpretive.
\end{enumerate}
\end{remark}

% --- source: emergence/e6_bfield_triples#6 (The six-momentum read by a clock)
\subsection{The six-momentum read by a clock}
\label{e6:sec:clock}

Let $P=\sum_i(\mathsf E_i/c)e_i+\sum_ap_ae_{a+3}$ with $\eta(P,P)=M_6^2c^2$, so $M_6^2c^4=|\boldsymbol{\mathsf E}|^2-|\bp|^2c^2$, as in \cref{en:temporal}.
For a clock $\hat n$ the energy read is $E(\hat n)=\hat n\cdot\boldsymbol{\mathsf E}$ and $m_4(\hat n)^2c^4=E(\hat n)^2-|\bp|^2c^2=M_6^2c^4-|\hat n\times\boldsymbol{\mathsf E}|^2$ (\cref{r2:prop:rotate}). A temporal $b$-field of constant rate $\nu$ about a unit axis $\hat{\mathbf e}$ acts on $P$ by $\boldsymbol{\mathsf E}(t)=R_{\hat{\mathbf e}}(\nu t)\boldsymbol{\mathsf E}$,
with $\bp$ fixed.

\begin{proposition}[The four-dimensional mass under a temporal $b$-field]
\label{e6:prop:clock}
\begin{enumerate}[label=(\roman*),leftmargin=2.2em]
\item The flow keeps $M_6$ and $|\boldsymbol{\mathsf E}|$ constant.
\item For a fixed clock $\hat n$, $E(\hat n,t)=A+B\cos(\nu t-\phi)$, where $A=(\hat n\cdot\hat{\mathbf e})(\hat{\mathbf e}\cdot\boldsymbol{\mathsf E})$ and
$B=|\hat n_\perp|\,|\boldsymbol{\mathsf E}_\perp|$ with $\hat n_\perp$ and $\boldsymbol{\mathsf E}_\perp$ the parts orthogonal to $\hat{\mathbf e}$. Hence $m_4(\hat n)^2c^4=M_6^2c^4-|\boldsymbol{\mathsf E}|^2+E(\hat n,t)^2$ oscillates. Its maximum is
$M_6^2c^4-|\boldsymbol{\mathsf E}|^2+(|A|+B)^2\le M_6^2c^4$. Its minimum is $M_6^2c^4-|\boldsymbol{\mathsf E}|^2+(|A|-B)^2$ if $|A|\ge B$ and
$M_6^2c^4-|\boldsymbol{\mathsf E}|^2$ otherwise. It is constant iff $B=0$, that is iff $\hat n\parallel\hat{\mathbf e}$ or $\boldsymbol{\mathsf E}\parallel\hat{\mathbf e}$.
\item The rotation about the clock, whose spin lift is proportional to $J(\hat n)$ (\cref{r2:prop:J}(iii)), keeps $E(\hat n)$. On the reduced shell
$\boldsymbol{\mathsf E}\parallel \hat n$ of \cref{thm:wlgauge} (\cref{r2:prop:shellII}) the flow preserves the constraint iff $\hat{\mathbf e}\parallel \hat n$.
\item With $c=1$, $\boldsymbol{\mathsf E}=(2,3,6)$ and $\bp=(0,0,3)$, so that $M_6^2=40$: the rotation about $e_3$ read by $e_1$ gives $E(\hat n)=2\cos\nu t-3\sin\nu t$ and
$m_4^2\in[-9,4]$; read by the symmetric clock $\widehat W=(e_1+e_2+e_3)/\sqrt3$ (\cref{prop:clock}) it gives $m_4^2\in[\tfrac{35}3-4\sqrt{26},\tfrac{35}3+4\sqrt{26}]$. The rotation about $\widehat W$ read by $\widehat W$ keeps
$m_4^2=\tfrac{94}3$ while the relative-time energies $\boldsymbol{\mathsf E}_\perp$ turn.
\end{enumerate}
\end{proposition}

\begin{proof}
Rodrigues' formula gives $\hat n\cdot\boldsymbol{\mathsf E}(t)=A+(\hat n\cdot\boldsymbol{\mathsf E}-A)\cos\nu t+\hat n\cdot(\hat{\mathbf e}\times\boldsymbol{\mathsf E})\sin\nu t$, and the squared
amplitude $(\hat n\cdot\boldsymbol{\mathsf E}-A)^2+(\hat n\cdot(\hat{\mathbf e}\times\boldsymbol{\mathsf E}))^2$ equals $(|\hat n|^2-(\hat n\cdot\hat{\mathbf e})^2)(|\boldsymbol{\mathsf E}|^2-(\hat{\mathbf e}\cdot\boldsymbol{\mathsf E})^2)$
(Lagrange's identity in the plane orthogonal to $\hat{\mathbf e}$). $|A|+B\le|\boldsymbol{\mathsf E}|$ by the Cauchy--Schwarz inequality. The rest is
substitution; the reduced-shell statement says that a vector stays on the line $\R \hat n$ under all rotations about $\hat{\mathbf e}$ iff the line is the axis.
\end{proof}

This is fixed (e6-5a, symbolic for three rational axes; e6-5b, exact). The reading needs care. ``The four-dimensional mass oscillates'' holds only
if the clock is held fixed while the temporal $b$-field turns the temporal momentum. Equivalently, the clock turns against a fixed momentum along the
circle of the clock sphere of \cref{r2:prop:rotate} about the axis $\hat{\mathbf e}$ through $\hat n$, a great circle only when $\hat n\perp\hat{\mathbf e}$ (e6-5c). Read in the sense of \cref{r2:rem:projection}, the energy read by the clock is
the projection of the undivided temporal energy, and the $b$-field moves energy periodically between the clock and the relative times. This is the
dynamical form of the obstruction of \cref{en:temporal}: splitting the temporal energy lowers the mass. The six-dimensional mass is
constant throughout. The rotation about the clock is the charge $J(\hat n)$ and changes nothing; the two clock rotations are what make $m_4$ oscillate, and
they are also what the kinetic term breaks (\cref{e6:prop:dirac}). On the reduced shell, the gauge-fixed reading of \cref{sec:constraints}, a temporal $b$-field
transverse to the clock is not admissible at fixed clock. Two clocks must be kept apart. For a \emph{fixed} clock $\hat n$, $m_4(\hat n)^2c^4=M_6^2c^4-|\hat n\times\boldsymbol{\mathsf E}(t)|^2$ oscillates and is negative whenever $|\hat n\times\boldsymbol{\mathsf E}|>M_6 c^2$, as in (iv), where $m_4^2$ reaches $-9$; then the projected four-momentum is spacelike and $m_4$ is not a rest mass of anything. For a \emph{co-rotating} clock, $\hat n(t)=R_{\hat{\mathbf e}}(\nu t)\hat n$, $|\hat n\times\boldsymbol{\mathsf E}|$ is constant and so is $m_4$. Which clock is physical is the open clock-selection question; the oscillation is a statement about a fixed reading, not a derived variation of a particle mass.

% --- source: emergence/e6_bfield_triples#7 (Between constituents)
\subsection{Between constituents}
\label{e6:sec:copies}

Write $\hat C_+=-\tfrac12\sum_{i\le3}E_{(1,i)}E_{(2,i)}$ and $\hat C_-=\tfrac12\sum_{a\ge4}E_{(1,a)}E_{(2,a)}$ for the temporal and the spatial parts of the copy rotation,
$\hat C=\hat C_++\hat C_-$, and $C_\times=\hat C_+-\hat C_-$. We write $C_\pm$ for $\hat C_\pm$.

\begin{proposition}[Temporal and spatial copy rotations]
\label{e6:prop:copyrot}
\begin{enumerate}[label=(\roman*),leftmargin=2.2em]
\item $C_\pm$ are antisymmetric and commute, and each is a sum of three commuting terms of square $-\tfrac14\id$. Hence $e^{2\pi C_\pm}=-\id$ and $e^{2\pi \hat C}=\id$.
The eigenvalues of $C_\pm$ are $\pm\tfrac i2$ ($24$ times each) and $\pm\tfrac{3i}2$ ($8$ times each).
\item On vectors neither $C_+$ nor $C_-$ commutes with $G_2$, $\hat C$ does, and $C_\times$ is compact and anticommutes with $G_2$.
\item $C_+=\tfrac12\hat C+\tfrac12C_\times$. $\hat C$ conserves $N$ and $C_\times=-\sum_i(a^+_{1,i}a^+_{2,i}+a^-_{1,i}a^-_{2,i})$ changes it by $\pm2$, one quantum in each
constituent. From the Fock vacuum, $e^{\theta C_+}|0\rangle$ has weight $\binom3k\cos^{2(3-k)}\tfrac\theta2\,\sin^{2k}\tfrac\theta2$ in $N=2k$, none in odd $N$,
and $N^{(1)}=N^{(2)}$ throughout. So $\langle N\rangle=3(1-\cos\theta)$, and $\theta=\pi$ fills all six modes.
\item As a constant field between two binding frames, $S=\kappa G_2+\beta V_+$, where $V_+$ is the vector form of $C_+$, has
$\det(\lambda-S)=\bigl(\lambda^4+(2\kappa^2+\beta^2)\lambda^2+\kappa^4\bigr)^3$. The rates are $\kappa$ and the lines are $s\pm\beta/2$. At $(\kappa,\beta)=(2,3)$, $M$ has Jordan blocks
at both lines.
\end{enumerate}
\end{proposition}

\begin{proof}
(i) $E_{(1,i)}E_{(2,i)}$ squares to $-\id$ for every $i$, and the terms for distinct $i$ involve four distinct anticommuting generators. Then
$e^{2\pi C_\pm}=\prod(\cos\pi\mp\sin\pi\,E_{(1,i)}E_{(2,i)})=(-1)^3$. The multiplicities follow from the minimal polynomial $(\lambda^2+1)(\lambda^2+9)$ of
$2C_\pm$ and $\Tr(2C_\pm)^2=-192$. (ii) $V_\pm$ maps $e^{(1)}_a$ to $e^{(2)}_a$ on the temporal (respectively spatial) directions, and $G_2$ exchanges the two
kinds; $V_+-V_-=\varepsilon\otimes\eta$. (iii) As in \cref{e6:prop:binding}(v): $E_{(1,i)}E_{(2,i)}\pm E_{(1,i+3)}E_{(2,i+3)}$ is a pair term or a hopping term. Since
$[\hat C,C_\times]=0$ and $\hat C|0\rangle=0$, $e^{\theta C_+}|0\rangle=e^{\theta C_\times/2}|0\rangle$, a product of three rotations in the planes
$\langle|0\rangle,a^+_{1,i}a^+_{2,i}|0\rangle\rangle$; this is checked exactly at $\cos\tfrac\theta2=\tfrac35$. (iv) As in \cref{e6:prop:binding}(iii), with
$\varepsilon\otimes\id_3$ in place of $[\hat{\mathbf r}]_\times$.
\end{proof}

Part (i) restates \cref{d7:prop:isospin}(i); the rest is fixed (e6-6a, e6-6c, e6-6e). Part (iii) is the compact counterpart of the copy boost
of \cref{e2:prop:copyboost}. The pairs emerge unitarily, with probability $\binom3k\cos^{2(3-k)}\tfrac\theta2\sin^{2k}\tfrac\theta2$ in place of the hyperbolic
weights there. As \cref{e2:sec:quanta} says, compactness does not protect the number of quanta; commutation with the binding root does, and $C_+$
fails it. Part (iv) puts $C_+$ outside the binding-compatible class of \cref{d7:thm:schr}: as a constant field it tears the binding of two
manifolds as a temporal $b$-field tears that of one.

\begin{proposition}[Mass transfer between constituents]
\label{e6:prop:transfer}
Let $H_g=m_1E_{(1,1)}+m_2E_{(2,1)}$, $M=(m_1^2+m_2^2)^{1/2}$, $u_1=e_1\otimes\id$ and $u_2=\id\otimes e_1$.
\begin{enumerate}[label=(\roman*),leftmargin=2.2em]
\item $\Ad(e^{\theta C_+})E_{(1,1)}=\cos\theta\,E_{(1,1)}+\sin\theta\,E_{(2,1)}$ and $\Ad(e^{\theta C_+})E_{(2,1)}=\cos\theta\,E_{(2,1)}-\sin\theta\,E_{(1,1)}$. Hence
$\Ad(e^{\theta C_+})H_g=m_1(\theta)E_{(1,1)}+m_2(\theta)E_{(2,1)}$, where $(m_1(\theta),m_2(\theta))$ is $(m_1,m_2)$ rotated through $\theta$. The spectrum stays $\{\pm M\}$, and
$(M,0)$ goes to $(M\cos\theta,M\sin\theta)$. $C_-$ commutes with $H_g$.
\item The ungraded $H_u=m_1u_1+m_2u_2$ has spectrum $\{\pm m_1\pm m_2\}$, each value $16$ times, and this multiset determines $\{|m_1|,|m_2|\}$. Rotating
the initial pair $(M,0)$ inside the ungraded family changes the spectrum unless $\sin2\theta=0$, from $\{\pm M\}$ to
$\{\pm M(\cos\theta+\sin\theta),\pm M(\cos\theta-\sin\theta)\}$. In general the spectrum is unchanged exactly when the unordered pair $(|m_1|,|m_2|)$ is: $(2,1)$
rotated with $\cos\theta=\tfrac45$ becomes $(1,2)$, with the same spectrum, although $\sin2\theta=\tfrac{24}{25}$.
$(\cos\theta\,u_1+\sin\theta\,u_2)^2=\id+\sin2\theta\,u_1u_2$, so for $\sin2\theta\neq0$ no invertible matrix conjugates $u_1$ into $\cos\theta\,u_1+\sin\theta\,u_2$,
whereas $(\cos\theta\,E_{(1,1)}+\sin\theta\,E_{(2,1)})^2=\id$. The product $u_1u_2$ of the two ungraded clocks squares to $\id$ and commutes with both. The flip $P(x\otimes y)=y\otimes x$ is orthogonal with
$Pu_1P^{-1}=u_2$: the exchange $(m_1,m_2)\mapsto(m_2,m_1)$ keeps the spectrum, it is the only change of $(|m_1|,|m_2|)$ that does, and no continuous
family at fixed spectrum changes $(|m_1|,|m_2|)$.
\item $\Ad(e^{\theta C_+})$ keeps the spectrum of $H_u$, but it maps $u_1$ to $\cos\theta\,u_1+\sin\theta\,\omega\otimes e_1$. This lies in
$\Cl(3,3)\otimes\id+\id\otimes\Cl(3,3)$ only for $\sin\theta=0$.
\end{enumerate}
\end{proposition}

\begin{proof}
(i) For $i=2,3$, $E_{(k,1)}$ anticommutes with both factors of $E_{(1,i)}E_{(2,i)}$, so these terms commute with $E_{(k,1)}$. The term
$X=E_{(1,1)}E_{(2,1)}$ anticommutes with $E_{(1,1)}$ and $E_{(2,1)}$, so $e^{-\theta X/2}E_{(1,1)}e^{\theta X/2}=e^{-\theta X}E_{(1,1)}=\cos\theta\,E_{(1,1)}+\sin\theta\,E_{(2,1)}$.
$C_-$ commutes with $E_{(k,1)}$ by the same count. (ii) $u_1$ and $u_2$ are commuting involutions whose products have trace zero, so each joint sign
occurs on $16$ dimensions. The largest and smallest absolute values in the multiset are $|m_1|+|m_2|$ and $||m_1|-|m_2||$. Conjugation preserves the
spectrum, and $\{\pm\cos\theta\pm\sin\theta\}\neq\{\pm1\}$ for $\sin2\theta\neq0$. $P$ exchanges the tensor factors, and since the spectrum fixes
$\{|m_1|,|m_2|\}$, a continuous family at fixed spectrum keeps $(|m_1|,|m_2|)$. (iii) $\omega\otimes e_1$ is a tensor of two non-identity blades, independent of every
$A\otimes\id$ and $\id\otimes B$.
\end{proof}

This is fixed (e6-6b, e6-6d, e6-6f). The checks are exact at $\cos\tfrac\theta2=\tfrac35$, where the rotated ungraded spectrum is
$\{\pm\tfrac{17}{25},\pm\tfrac{31}{25}\}M$. Part (i) is \cref{e1:prop:rest}(ii) realised by the temporal copy field. The single bivector $B=E_{(2,1)}E_{(1,1)}$
used there is one of the three terms of $2C_+$, and the other two leave the clocks alone.

\begin{remark}[Does a second particle acquire mass from the first?]
\label{e6:rem:transfer}
The algebra is fixed. The temporal copy rotation turns the clock $E_{(1,1)}$ of the first constituent into the clock $E_{(2,1)}$ of the second, and
from $(M,0)$ it produces $(M\cos\theta,M\sin\theta)$ at fixed graded mass: a massless second copy acquires mass from the first. Reading this as ``a
second particle acquires mass from the first through the $b$-field'' needs two further structures. One is a clock per copy, $E_{(k,1)}$ as the
$\gamma^0$ of copy $k$. The other is the graded reading, and there \cref{e1:prop:rest} applies: the copies admit no simultaneous nonzero-mass Dirac constraints of the graded form, and the rest energy is
the single value $M$. In that reading the process is a rotation of the mass direction of one particle between two internal clock directions,
the masses adding in quadrature as in \cref{d6:prop:dirac}(iv). It is not the creation of a second massive particle. In the ungraded
reading, where energies add and there are two particles (\cref{e1:prop:rest}(iii)--(iv)), the continuous transfer is obstructed: no conjugation turns
the clock $u_1$ into $\cos\theta\,u_1+\sin\theta\,u_2$ with $\sin2\theta\neq0$, and the additive spectrum fixes the unordered pair of masses
(\cref{e6:prop:transfer}(ii)). The only spectrum-preserving change is the exchange of the two masses by the flip of the tensor factors, a
relabelling of the constituents rather than a transfer. The transfer is moreover not
binding-compatible, and half of it creates quanta in pairs across the constituents (\cref{e6:prop:copyrot}(iii)). So the mass transfer between
constituents is fixed algebra; the reading of it as the emergence of a second massive particle is interpretive in the one-particle reading and,
as a continuous process, obstructed in the two-particle reading.
\end{remark}

% --- source: emergence/e6_bfield_triples#8 (Up the ladder: the mass forms of a rung)
\subsection{Up the ladder: the mass forms of a rung}
\label{e6:sec:ladder}

\begin{theorem}[Mass forms of a rung]
\label{e6:thm:ladder}
At rung $n$ let $T_n$ be the product, in increasing order, of the $3n$ temporal generators $E_{(k,i)}$ ($i\le3$), $S_n$ that of the $3n$ spatial ones, and
$\Gamma_n$ the volume element of $\Cl(3n,3n)$ (here $S_n$ is the spatial volume, not the spinor module $(\R^8)^{\otimes n}$ of equation~\cref{eq:tower}).
\begin{enumerate}[label=(\roman*),leftmargin=2.2em]
\item For every $n$, $Y^{\!\top}T_n+T_nY=0=Y^{\!\top}S_n+S_nY$ for every bivector $Y$, $T_nS_n=\pm\Gamma_n$, and the forms on $(\R^8)^{\otimes n}$ invariant under
$\so{3n,3n}$ are exactly $\operatorname{span}\{T_n,S_n\}$.
\item $T_n^{\!\top}=(-1)^{3n(3n-1)/2}T_n$ and $S_n^{\!\top}=(-1)^{3n+3n(3n-1)/2}S_n$. For $n\equiv1,2,3,0\pmod4$ the pair $(T_n,S_n)$ is (antisymmetric,
symmetric), (antisymmetric, antisymmetric), (symmetric, antisymmetric), (symmetric, symmetric). Hence the Grassmann mass kernel invariant under
$\so{3n,3n}$ is unique iff $n$ is odd. It is $T_n$ for $n\equiv1$ and $S_n$ for $n\equiv3\pmod4$. For $n\equiv2$ the kernels form the plane $\operatorname{span}\{T_n,S_n\}$,
and for $n\equiv0$ there is none.
\item At $n=2$, $T_2=-\omega_S\otimes T$ and $S_2=-T\otimes\omega_S$. The four products $M_1\otimes M_2$, $M_k\in\{T,\omega_S\}$, span the forms invariant under
$\so{3,3}^{(1)}\oplus\so{3,3}^{(2)}$; $T\otimes T$ and $\omega_S\otimes\omega_S$ are symmetric, and $\so{6,6}$ keeps exactly the antisymmetric ones,
$T_2$ and $S_2$. The per-copy $T\otimes\id$ is not invariant even under the boosts of copy~$2$.
\item The forms invariant under the temporal rotations $\so{3n}_+$ are their commutant, of dimension $16$, $128$ and $1024$ at $n=1,2,3$, with
(symmetric, antisymmetric) dimensions $(6,10)$, $(64,64)$ and $(528,496)$. They contain $T_n$ and $S_n$.
\item The forms invariant under the diagonal $\so{3,3}$, the same bivector on every constituent, have dimension $2$, $14$ and $168$ at $n=1,2,3$, with
(symmetric, antisymmetric) dimensions $(1,1)$, $(10,4)$ and $(84,84)$.
\end{enumerate}
\end{theorem}

\begin{proof}
(i) At every rung the temporal generators are symmetric and the spatial ones antisymmetric, $E_A^{\!\top}=\eta_{AA}E_A$, since $\omega^{\!\top}=\omega$. By the
blade rule in $\Cl(3n,3n)$, $T_nE_AT_n^{-1}=(-1)^{3n-1}\eta_{AA}E_A=(-1)^{3n-1}E_A^{\!\top}$ for every $A$, so for $Y=E_AE_B$,
$T_nYT_n^{-1}=E_A^{\!\top}E_B^{\!\top}=(E_BE_A)^{\!\top}=-Y^{\!\top}$, which is the invariance. The same holds for $S_n$ with the sign $(-1)^{3n}$. Conversely, if $M$ is
invariant, then $Y^{\!\top}=-T_nYT_n^{-1}$ gives $Y(T_n^{-1}M)=(T_n^{-1}M)Y$ for every bivector. The commutant of the bivectors is spanned by the blades it
contains. A blade $e_I$ commutes with $E_AE_B$ iff $|I\cap\{A,B\}|$ is even, and this holds for all $A\neq B$ iff $I=\emptyset$ or $I$ is everything. So
$M\in T_n\operatorname{span}\{\id,\Gamma_n\}=\operatorname{span}\{T_n,S_n\}$. (ii) The transpose of a product of $m$ symmetric (respectively antisymmetric)
generators is the reversed product times $1$ (respectively $(-1)^m$), and reversal costs $(-1)^{m(m-1)/2}$. (iii) $E_{(2,1)}E_{(2,2)}E_{(2,3)}=\omega^3\otimes T=\omega\otimes T$
and $T\omega=-\omega_S$. Within each copy the invariant forms are $\operatorname{span}\{T,\omega_S\}$, and blade conjugation is diagonal. (iv) A compact $Y$
is antisymmetric, so invariance is commutation; the dimensions are exact orbit counts. (v) Let $c(\mu)$ be the weight multiplicities of a module for
the split Cartan subalgebra spanned by the pole operators $e_ie_{i+3}$, which are diagonal. The multiplicity of the trivial module is
$\sum_{w\in W(D_3)}\det(w)\,c(\rho-w\rho)$ \cite{Humphreys72}. It is evaluated on $\mathrm{Sym}^2$ and $\Lambda^2$ of $(\R^8)^{\otimes n}$ with the exact weights; at $n=2$ a
direct exact solution on the $216$ entries of opposite weight agrees.
\end{proof}

The theorem is fixed (e6-7a to e6-7e). At $n=1$, (i) restates \cref{cor:fullmass} and the pairing of \cref{sec:selector}, and the
antisymmetric count $10$ of (iv) restates the kernel space of $\so T$ in \cref{thm:subkernels}. The pattern of
(ii) is Chevalley's: on the spinors of $\R^{m,m}$ the two invariant forms have symmetries $(-1)^{m(m-1)/2}$ and $(-1)^{m(m+1)/2}$
\cite{Chevalley54,LawsonMichelsohn89}, here with $m=3n$ (inherited). What the arena adds is that the two forms are the temporal and the spatial volume
of the rung, with the Klein factors of (iii). The uniqueness of the mass in $T=e_1e_2e_3$ is therefore a statement of the first rung. At the three-constituent rung the unique $\so{9,9}$-invariant Grassmann kernel is the spatial volume $S_3$, and $T_3$ is a symmetric (bosonic) form. Under
the diagonal action the mass is not unique from $n=2$ on, and under the temporal rotations alone not even at $n=1$. At the second rung the
antisymmetric invariant kernels pair the temporal volume of one constituent with the spatial volume of the other (iii): no invariant mass kernel
belongs to one constituent alone.

The invariant bilinear forms of spinor modules of $\mathrm{Spin}(p,q)$, with their symmetry types by signature, are classified by Alekseevsky and
Cort\'es \cite{AlekseevskyCortes97}; (i) and (ii) are that classification in the realisation of the tower.

% --- source: emergence/e6_bfield_triples#9 (The two assignments compared)
\subsection{The two assignments compared}
\label{e6:sec:compare}

\begin{remark}[The temporal $b$-field and the $b$-field between constituents]
\label{e6:rem:compare}
\begin{enumerate}[label=(\alph*),leftmargin=2.2em]
\item \emph{How they relate} (fixed). In \cref{sec:bfield,d7:sec:main} the $b$-field is the $\eta$-antisymmetric part of the shift $S=\kappa G_n+b$, and
$\kappa G_n$ is its $\eta$-self-adjoint part, $(\kappa G)^2=-c^2/\ell^2$; $\kappa G$ maps $V^+$ to $V^-$ and back, and that is where $c$ enters. \Cref{d7:sec:main} studies the binding-compatible fields, $[b,G_n]=0$. A $b$-field on the temporal triple is part of the connection, compact and
gauge-admissible (\cref{cor:noboost}), but it is not binding-compatible: half of it lies in $\mathfrak p_G$
(\cref{e6:prop:binding}(i)). It belongs to the compact class of \cref{sec:bfield} with $\mathbf a=\mathbf c$, it transports $\kappa G$ to
$\kappa G_{e^{-\theta R}}$ without changing the rates, and it tears the binding (\cref{e6:prop:binding}). Between constituents the same holds for the
temporal copy rotation (\cref{e6:prop:copyrot}(iv)), which lies outside the Schr\"odinger class of \cref{d7:thm:schr}.
\item \emph{A non-abelian Yang--Mills $\su2$ with Cartan direction $J$} (fixed; the Cartan choice is the clock's). $\so3_+\cong\su2$ is compact. It acts on
$V$ as isospin one on $V^+$ and trivially on $V^-$, and on spinors as isospin one-half (\cref{e6:prop:dirac}(iii)). The algebra fixes no Cartan direction,
since all Cartan directions of $\so3_+$ are conjugate (all its elements of a fixed nonzero norm are; a12-2). The kinetic term fixes it to $J=J(e_1)$, the rotation about the clock, which it keeps. The two
off-diagonal directions are the clock rotations; they flip the charge, and their half turn acts as $\mathsf{CP}$.
\item \emph{A mass term for the off-diagonal components} (identity fixed; reading interpretive; mechanism open). The breaking of the kinetic forms
by $Y\in\so3_+$ is $Y^{\!\top}K^a+K^aY=2Y(\hat n)\gamma^a$, with total weight $96|Y(\hat n)|^2$ (\cref{e6:prop:dirac}(i)). This is a quadratic form on $\so3_+$ with
kernel $\R J(\hat n)$, invariant under the stabiliser of the clock: an invariant statement of the breaking. It has the form of the adjoint-Higgs mass
term $g^2|[b,\langle\Phi\rangle]|^2$ of the Georgi--Glashow $SO(3)$ model \cite{GeorgiGlashowPRL72}, with the clock in the role of the Higgs vector
\cite{EnglertBrout64,Higgs64,Kibble67}: the two charge-flipping components would be massive and $J(\hat n)$ massless. That needs a clock field $\hat n(x)$ with a
covariant kinetic term, and the arena supplies none: a dynamics that selects the clock is open (\cref{r2:sec:main}). Gauge invariance forbids a
bare mass term (inherited), and the Dirac mass gives none, since $\so3_+$ preserves $T$ (\cref{e6:prop:Tform}). In \cref{d7:rem:verdict}(e) the rate matrix
played the role of the adjoint vacuum value; here the clock does. The pattern is that of an adjoint $SO(3)\to U(1)$ breaking, not of the electroweak $SU(2)_L\times U(1)_Y\to U(1)_{\rm em}$ by a doublet, and nothing here derives an electroweak Higgs mechanism.
\item \emph{What the assignment fixes, constructs, obstructs and leaves open.} Fixed: every finite-dimensional statement of
\cref{e6:sec:Tform,e6:sec:dirac,e6:sec:binding,e6:sec:polar,e6:sec:clock,e6:sec:copies,e6:sec:ladder}. Constructed: the clock, which fixes the Cartan
direction and the kinetic term, and the temporal $b$-field itself. Obstructed: its removal along a worldline, bounded binding under it, continuous mass
transfer at fixed spectrum in the ungraded reading, and a unique Grassmann mass kernel at even rungs. Open: a dynamics for the clock, hence a mass for the
off-diagonal components.
\item \emph{What it implies for the emergence of a second particle} (interpretive, drawn from the fixed statements). The temporal $b$-field gives the
second rung a process that moves mass from one constituent to the other, but only as a rotation of one mass vector in the graded reading, in which
the rung carries one particle. In the reading that has two particles it is obstructed as a continuous process; only
the exchange of the two masses, a relabelling, keeps the spectrum. The uniqueness of the mass in $T$, which the assignment invokes,
does not survive the ladder in a form that assigns a mass to each particle: the invariant kernels of the second rung pair the temporal volume of one
constituent with the spatial volume of the other.
\end{enumerate}
\end{remark}

% --- source: emergence/e6_bfield_triples#10 (Answer)
\subsection{Answer}
\label{e6:sec:answer}

The assignment ``$c$ on $(-,-,-)$, $b$ on $(+,+,+)$'' has a precise and correct form. The temporal rotations cannot change the symmetric factor
$CC^{\!\top}$ of the conversion, which lives on the spatial triple and is $c^2$ in the isotropic case; they keep the mass form $T$ and the rest
flow, and they move the orthogonal factor of the conversion, the clock and the binding root. The price is the binding and the kinetic term: a
temporal $b$-field keeps the rates, splits the lines, cannot be gauged away and tears the binding, and the kinetic term breaks it to $\uu1_J$ in
proportion to the speed of the clock. The mass is unique in $T=e_1e_2e_3$ on the first rung only; up the ladder the unique invariant Grassmann kernel is the temporal volume
$T_n$ for $n\equiv1$ and the spatial volume $S_n$ for $n\equiv3\pmod4$, and at even rungs there is a plane of them or none. Between two constituents the temporal $b$-field rotates
graded mass from one clock to the other at fixed total, a statement about one particle, not two.

\begin{table}[htbp]
\centering
\footnotesize
\begin{tabular}{@{}>{\raggedright\arraybackslash}p{8.0cm}>{\raggedright\arraybackslash}p{4.7cm}>{\raggedright\arraybackslash}p{2.1cm}@{}}
\toprule
statement & class & checks\\
\midrule
conventions; the triples, $T$, $\omega_S$, Fock modes & restated (\cref{eq:realisation}, \eqref{eq:tower}, \cref{thm:subdichotomy}) & e6-0\\
\cref{e6:prop:Tform}(i): $T$ and $\omega_S$ invariant forms under $\mathrm{Spin}(3,3)$ & restated (\cref{cor:fullmass}, \cref{sec:selector}; \cref{mb:Tinv}) & e6-1a\\
\cref{e6:prop:Tform}(ii), (iii): commutes with $\so3_+\oplus\so3_-$ only; commutant $M_4(\C)$, $32$ blades & fixed ($M_4(\C)$ inherited) & e6-1b\\
\cref{e6:prop:Tform}(iv): $\Ad_T=\eta$; boosts conjugate $T$ to $e^{2sY}T$ & restated (\cref{en:parity}); fixed & e6-1c\\
\cref{e6:prop:dirac}(i): kinetic breaking $2Y(\hat n)\gamma^a$, weight $96|Y(\hat n)|^2$; $\so3_+\cap(\lint\oplus\R J)=\R J$ & fixed (covariance restated, \cref{mb:Tinv}(ii)) & e6-2a\\
\cref{e6:prop:dirac}(ii): clock rotation of $H(0)$; transport to the clock $\hat n$ & fixed (transport restated, \cref{r2:prop:J}, \cref{r2:thm:i}) & e6-2b\\
\cref{e6:prop:dirac}(iii), (iv): rest flow kept; isospin $\tfrac12$; $\mathsf{CP}$ half turn; octant duality & fixed; (iv) restated (\cref{r3:prop:odd}) & e6-2c\\
\cref{e6:prop:twosu2}: $\so3_+$ and $\su2_{\rm PG}$ two diagonals of $\so4$, meeting in $\R J$ & fixed ((v) restated, \cref{prop:gauge1}) & e6-2d\\
\cref{e6:prop:binding}(i), (ii): half $\lG$, half $\mathfrak p_G$; $G\to G_{e^{-\theta R}}$; rates kept; not removable & fixed (rates restated, \cref{d7:thm:schr}(iii)) & e6-3a, e6-3b\\
\cref{e6:prop:binding}(iii), (iv): lines $\kappa$, $s\pm\omega/2$; Jordan blocks & fixed & e6-3c, e6-3d\\
\cref{e6:prop:binding}(v): paired half conserves $N$, counter half changes it by $\pm2$ & fixed & e6-3e\\
\cref{e6:prop:polar}: $C=cR$; twisted diagonal; orbit $O(3)$; $CC^{\!\top}$ invariant; $q_C$ & fixed & e6-4a to e6-4e\\
\cref{e6:rem:csym}: ``$c^2$ on $(-,-,-)$, $b$ on $(+,+,+)$'' & (a), (b) fixed; reading interpretive & e6-4d, e6-4f\\
\cref{e6:prop:clock}: $m_4$ oscillates at fixed $\mu$; constant iff on the axis & fixed; the reading interpretive & e6-5a to e6-5c\\
\cref{e6:prop:copyrot}: $C_\pm$, $e^{2\pi C_\pm}=-\id$; pairs of quanta; tearing between frames & fixed ((i) restated, \cref{d7:prop:isospin}(i)) & e6-6a, e6-6c, e6-6e\\
\cref{e6:prop:transfer}: graded mass transfer at fixed $M$; ungraded obstruction & fixed & e6-6b, e6-6d, e6-6f\\
\cref{e6:rem:transfer}: ``a second particle acquires mass'' & interpretive (one-particle reading); obstructed as a continuous process (two-particle reading) & ---\\
\cref{e6:thm:ladder}(i), (ii): invariant forms $\operatorname{span}\{T_n,S_n\}$; period four & fixed (Chevalley's pattern inherited) & e6-7a, e6-7b\\
\cref{e6:thm:ladder}(iii)--(v): Klein factors; temporal and diagonal counts & fixed & e6-7a, e6-7c to e6-7e\\
\cref{e6:rem:compare}: relation to RE; Yang--Mills $\su2$; clock as adjoint Higgs & (a), (b) fixed; (c) identity fixed, reading interpretive, mechanism open; (e) interpretive & e6-2a\\
\bottomrule
\end{tabular}
\caption{Claims of \cref{e6:sec:main}, their classes and their checks, which are \texttt{PASS} lines of \texttt{e6\_bfield\_triples.py}. ``Restated'': fixed in an earlier part of the manuscript, re-verified here.}
\label{e6:tab:claims}
\end{table}


\section{The ladder: a table, and what changes along it}
\label{sec:ladder}

The two subsections below collect, rung by rung, what \cref{ch:ch05_tower} and this chapter establish about the tower: the table of the
ladder, and what changes, and what does not, in going from $n$ constituents to $n+1$.

% --- source: emergence/e3_ladder#7 (The ladder table)
\subsection{The ladder table}
\label{e3:sec:table}

\begin{table}[t]
\centering
\footnotesize
\setlength{\tabcolsep}{4pt}
\begin{tabular}{@{}>{\raggedright\arraybackslash}p{4.5cm}ccccc>{\raggedright\arraybackslash}p{2.9cm}@{}}
\toprule
structure & $n=1$ & $2$ & $3$ & $4$ & $5$ & source\\
\midrule
$\dim V_n$, $\dim S_n$ & $6$, $8$ & $12$, $64$ & $18$, $512$ & $24$, $8^4$ & $30$, $8^5$ & ---\\
$\Cl(3n,3n)$ & \multicolumn{5}{c}{$M_{8^n}(\R)$ at every rung} & e3-1a ($n\le4$)\\
Fock modes & $3$ & $6$ & $9$ & $12$ & $15$ & e3-1c\\
spectator centraliser $\so{3n-1,3n-3}$ & $1$ & $28$ & $91$ & $190$ & $325$ & \cref{thm:centraliser}; e3-6a\\
commutant of diagonal $\lint$, $\so n\oplus\so{2n}$ & $1$ & $7$ & $18$ & $34$ & $55$ & \cref{thm:algebra}; e3-6a\\
commutant of diagonal $\lG$, $\so{n,\C}$ & $0$ & $2$ & $6$ & $12$ & $20$ & \cref{prop:commutantsn}; e3-6a\\
its compact part, $\so n$ & $0$ & $1$ & $3$ & $6$ & $10$ & e3-6a\\
rotor algebra $\uu n$ (in $\so{6n}$) & $1$ & $4$ & $9$ & $16$ & $25$ & \cref{prop:commutantsn}; e3-6a\\
copy algebra, commutant of diagonal $\so{3,3}$ & $0$ & $1$ & $3$ & $6$ & $10$ & e3-4a\\
$C_{\rm cop}+C_{\rm diag}$ & $\tfrac{15}4$ & $9$ & $\tfrac{63}4$ & $24$ & $\tfrac{135}4$ & e3-4c\\
graded rest mass & $m$ & $\sqrt2\,m$ & $\sqrt3\,m$ & $2m$ & $\sqrt5\,m$ & e3-2a\\
ungraded rest energies $|m(n-2j)|$ & $m$ & $0,2m$ & $m,3m$ & $0,2m,4m$ & $m,3m,5m$ & e3-2b\\
$2\pi$ rotation, diagonal $\lint$ & $-\id$ & $+\id$ & $-\id$ & $+\id$ & $-\id$ & e3-3a\\
$2\pi$ rotation, diagonal $\lG$ & $+\id$ & $+\id$ & $+\id$ & $+\id$ & $+\id$ & e3-3b\\
copy-invariant four-vector & yes & no & yes & no & yes & \cref{thm:copyparity}; e3-6b\\
selected imaginary unit (diagonal $\lint$) & yes & no & yes & no & yes & \cref{e1:prop:select}(iv)\\
invariant number of quanta (diagonal $\lG$) & no & yes & no & yes & no & \cref{e2:thm:invnumber}\\
invariant Grassmann mass kernel & $T_1$ & plane & $S_3$ & none & $T_5$ & \cref{e6:thm:ladder}(ii)\\
copy isospin & --- & integral & integral & $j_\pm\in\half\Z$ & --- & \cref{d7:prop:isospin}; e3-4d\\
projector $P_*$ & --- & --- & yes & --- & --- & \cref{thm:PS}\\
$\Lambda^n_\C$ of the Dirac module on $W_n$ & Dirac & $2(0,0){\oplus}(\half,\half)$ & Dirac & $(0,0)$ & $0$ & e3-5d\\
Koszul-antisymmetric part of $S_n$ & $8$ & $32$ & $88$ & $192$ & $360$ & e3-5e\\
\bottomrule
\end{tabular}
\caption{The ladder. Dimensions of Lie algebras are real. ``Plane'' and ``none'': at $n\equiv2\pmod4$ the invariant Grassmann kernels form a plane,
at $n\equiv0$ there is none; $T_n$ and $S_n$ there are the temporal and spatial volumes of \cref{e6:thm:ladder}, not the spinor module. The copy
isospin is that of the six-direction copy algebra.}
\label{e3:tab:ladder}
\end{table}

\Cref{e3:tab:ladder} collects the rung-by-rung data. The Lie algebra rows are fixed (e3-6a, exact over $\Q$ at $n\le5$) and restate
\cref{thm:centraliser}, \cref{thm:algebra} and \cref{prop:commutantsn}; the row of copy-invariant four-vectors restates \cref{thm:copyparity}, with
the case $n=3$ and the absence at $n=2$ re-verified (e3-6b). Read as a whole, the table says which structures single out which rungs.

\emph{No rung is singled out by the algebra.} $\Cl(3n,3n)=M_{8^n}(\R)$, its module is $S_n$, and every Lie algebra row is a polynomial in $n$. The
first structure with a period in $n$ is the symmetry of the invariant bilinear forms on $S_n$, of period four (\cref{e6:thm:ladder}(ii)). It is
Chevalley's pattern for spinors of $\R^{3n,3n}$, the dimension $6n$ running through the residues $6,4,2,0$ modulo $8$ with period four in $n$
(inherited \cite{Chevalley54}). Through it the $\so{3n,3n}$-invariant Grassmann mass kernel is unique exactly at odd $n$. That is the only
selection the algebra makes by itself, and it does not distinguish $n=3$ from the other odd rungs.

\emph{Odd rungs are singled out by the interior Lorentz action.} Under the diagonal $\lint$, odd $n$ is where the composite is a fermion
(\cref{e3:prop:spin}(i)), where an imaginary unit is selected (\cref{e1:prop:select}(iv)), where copy-invariant four-vectors and the Dirac block
exist (\cref{thm:copyparity} and \cref{prop:diagtower}). \emph{Even rungs
are singled out by the binding action}: only there does a Lorentz-invariant number of quanta exist under the diagonal $\lG$
(\cref{e2:thm:invnumber}). \emph{The third rung is singled out by added projectors and by the low-rank coincidences of the diagonal algebra}: the
chiral Pati--Salam projector $P_*$ is constructed at $n=3$ (\cref{thm:PS}), the diagonal connection algebra there is
$\so3\oplus\so6\cong\su2\oplus\su4$ with colour $\su3$ inside $\so6$, and the copy isospin is integral (\cref{prop:diagtower}(iii);
\cref{d7:prop:isospin}). At $n=5$ the Dirac block carries one Standard-Model family per copy multiplet (\cref{prop:diagtower}(iv)). Apart from the period of the forms, every one of these selections comes from a Lorentz action, a projector or a gauge-algebra
identification added to the rung. None comes from the ladder.

% --- source: emergence/e8_synthesis#3 (What changes up the ladder)
\subsection{What changes up the ladder}

Rung by rung the algebra changes polynomially in $n$, and no rung is singled out by the algebra alone. The only periodic structure is the
period four of the invariant spinor forms (Chevalley's pattern), through which the $\so{3n,3n}$-invariant Grassmann mass kernel is unique
exactly at odd $n$. Parity selections come from added structures: odd $n$ from the diagonal interior action (spin parity $(-1)^n$, selected
imaginary unit, copy-invariant four-vectors), even $n$ from the common binding action (an invariant number of quanta), and $n=3$ from the
projector $P_*$ and the low-rank coincidences of the gauge algebra. Between rungs there is no process in the algebra: rung $n+1$ is not a
state or an element of rung $n$, and the tensor embedding $\psi\mapsto\psi\otimes\phi_0$ needs a state $\phi_0$ that nothing selects.

\begin{table}[!htbp]
\centering
\footnotesize
\begin{tabular}{@{}>{\raggedright\arraybackslash}p{6.1cm}>{\raggedright\arraybackslash}p{2.4cm}>{\raggedright\arraybackslash}p{3.0cm}>{\raggedright\arraybackslash}p{3.0cm}@{}}
\toprule
claim & class & location & evidence\\
\midrule
$\Cl(3n,3n)=M_{8^n}(\R)$; step with multiplicity $8$; commuting copy unique up to conjugation & proved algebra & \cref{e1:prop:copies,e3:prop:rung} & e1-1a--d, e3-1a--c, y1-1a--f\\
graded masses in quadrature, no simultaneous nonzero-mass graded Dirac constraints; ungraded spectra additive & proved algebra & \cref{e1:prop:rest,e3:prop:mass} & e1-2a--g, e3-2a--b, y1-2a--g, z1-1c\\
common complex structure only on $W_n$; selection by actions; none at even $n$ & proved algebra; obstruction & \cref{e1:prop:W,e1:prop:select,e1:cor:nocx} & e1-3a--h, y1-3a--4f\\
$N$ conserved exactly by $\mathfrak{gl}(3n,\R)$; compact $\Leftrightarrow$ paired rotation (in $\so{3n}\oplus\so{3n}$); no odd change & proved algebra & \cref{e1:prop:Nbiv,e1:thm:Nrule,e2:lem:N} & e1-4a--g, e2-10--13, y2-11--14, z1-1a\\
compact fields outside $\Cent(G_n)$ ($J$, $e_1e_2$, $S_k$) create pairs of quanta & proved algebra & \cref{e1:prop:Nlorentz,e2:prop:bogo} & e2-14--16, z1-1b\\
$N$ not Lorentz invariant; invariant number iff $n$ even (binding action) & proved algebra & \cref{e2:thm:invnumber} & e2-19--22, y2-20--24\\
no equivariant vertex $S\to S_n$ for even $n$ (interior action) & proved algebra (selection rule inherited) & \cref{e2:prop:vertex} & e2-2--6, y2-2--6\\
$C_{\rm cop}+C_{\rm diag}=\tfrac34n(n+4)$; Howe duality; Koszul $\mathfrak S_n$ & proved algebra (constant and duality inherited, Neerg{\aa}rd, Howe) & \cref{e3:thm:dual,e3:prop:koszul} & e3-4a--f, e3-5a--e, z1-3a--d\\
compact copy field $=d\Gamma(H)$: transfer, not creation; no binding in this class & proved algebra & \cref{e4:thm:fock,e4:cor:transfer} & e4-1a--f\\
$c_n$ triangle signs; $SO(4)$ doublets on odd $N$; no copy spinor & proved algebra; CP reading not established & \cref{e4:thm:cycles,e4:thm:so4,e4:cor:nospinor} & e4-2a--3e, z1-4a--b\\
pair invariant: stabiliser exactly $\lG^\Delta\oplus\R C_v$ & proved (exact certificate) & \cref{e5:thm:pair} & e5-1c, y5, z1-5b\\
mass form: $Y^{\!\top}T+TY=0$ on $\so{3,3}$, $Y^{\!\top}TY=-TY^2$ & proved (corrected ordering) & \cref{e6:prop:Tform} & e6-1a--c, y6, z1-5a\\
$W_n$, $H_u$, plain exchange with antisymmetry; $n$ fermions at one label, exclusion $n\le4$ & constructed; not selected & \cref{e1:sec:answer,e3:prop:ungraded} & e3-5d, z1-3c\\
mixed-action vertex; Rabi transfer with external $i$; non-abelian frames on a harmonic ladder & constructed & \cref{e2:prop:vertex,e2:prop:hop,e4:prop:rotor} & e2-5, e2-23--25, e4-5a--e, z1-4c\\
no selected two-particle complex structure at $n=2$; no PS in the five candidate algebras & obstruction (stated class) & \cref{e1:cor:nocx,e5:prop:noPS} & e1-3c--g, e5-3a--i\\
graded composite below free pair at equal labels; not a bound state & proved inequality; binding not supported & \cref{e2:prop:threshold} & e2-8, y2-8, z1-2a\\
uphill decays in $(3,3)$; none after reduction to one clock & inherited kinematics; conditional on the reduction & \cref{e7:prop:uphill} & e7-16, e7-17\\
``one particle'' or ``$n$ particles'' & interpretive & \cref{e1:sec:answer,e3:sec:answer} & ---\\
Poincar\'e state space, asymptotic states, interacting theory, dynamical selection, spin--statistics & missing & --- & ---\\
\bottomrule
\end{tabular}
\caption{Claims about particle multiplicity collected along the ladder, their classes and their evidence. A \texttt{PASS} line shows that the script's
assertion holds; the class column records what the claim establishes. The checks are listed in \cref{app:methods}.}
\label{e8:tab:claims}
\end{table}

